\chaptertitle{Air on G's String}

                                  Air on G's String
       The Tortoise and Achilles have just completed a tour of a porridge factory.

Achilles: You don't mind if I change the subject, do you? Tortoise: Be my guest.
Achilles: Very well, then. It concerns an obscene phone call I received a few days ago.
Tortoise: Sounds interesting.
Achilles: Yes. Well-the problem was that the caller was incoherent, at least as far as I
  could tell. He shouted something over the line and then hung up-or rather, now that I
  think of it, he shouted something, shouted it again, and then hung up.
Tortoise: Did you catch what that thing was?
Achilles: Well, the whole call went like this:

       Myself: Hello?
       Caller (shouting wildly): Yields falsehood when preceded by its quotation! Yields
              falsehood when preceded by its quotation!
       (Click.)

Tortoise: That is a most unusual thing to say to somebody on an obscene phone call.
Achilles: Exactly how it struck me.
Tortoise: Perhaps there was some meaning to that seeming madness.
Achilles: Perhaps.

   (They enter a spacious courtyard framed by some charming three-story stone
     houses. At its center stands a palm tree, and to one side is a tower. Near the
     tower there is a staircase where a boy sits, talking to a young woman in a
     window.)

Tortoise: Where are you taking me, Achilles?
Achilles: I would like to show you the pretty view from the top of this tower.
Tortoise: Oh, how nice.

  (They approach the boy, who watches them with curiosity, then says something to
    the young woman-they both chuckle. Achilles and Mr. T, instead of going up the
    boy's staircase, turn left and head down a short flight of stairs which leads to a
    small wooden door.)
  Achilles: We can just step inside right here. Follow me.




Air on G's String                                                                     439

FIGURE 74. Above and Below, by M.C. Escher (lithograph 1947).




Air on G's String                                               440

   (Achilles opens the door. They enter, and begin climbing the steep helical staircase
      inside the tower.)

Tortoise (puffing slightly): I'm a little out of shape for this sort of exercise,
Achilles. How much further do we have to go?
Achilles: Another few flights ... but I have an idea. Instead of walking on the top side of
  these stairs, why don't you walk on the underside?
Tortoise: How do I do THAT?
Achilles: Just hold on tightly, and climb around underneath-there's room enough for you.
  You'll find that the steps make just as much sense from below as from above ...
Tortoise (gingerly shifting himself about): Am I doing it right?
Achilles: You've got it!
Tortoise (his voice slightly muffled): Say-this little maneuver has got me confused.
  Should I head upstairs or downstairs, now?
Achilles: Just continue heading in the same direction as you were before. On your side of
  the staircase, that means go DOWN, on mine it means UP.
Tortoise: Now you're not going to tell me that I can get to the top of the tower by going
  down, are you?
Achilles: I don't know, but it works ...

   (And so they begin spiraling in synchrony, with A always on one side, and T
     matching him on the other side. Soon they reach the end of the staircase.)

   Now just undo the maneuver, Mr. T. Here-let me help you up.

   (He lends an arm to the Tortoise, and hoists him back to the other side of the
     stairs.)

Tortoise: Thanks. It was a little easier getting back up.

   (And they step out onto the roof, overlooking the town.)

  That's a lovely view, Achilles. I'm glad you brought me up here-or rather, DOWN
      here.
Achilles: I figured you'd enjoy it.
Tortoise: I've been thinking about that obscene phone call. I think I understand it a little
  better now.
Achilles: You do? Would you tell me about it?
Tortoise: Gladly. Do you perchance feel, as I do, that that phrase "preceded by its
  quotation" has a slightly haunting quality about it?
Achilles: Slightly, yes-extremely slightly.
Tortoise: Can you imagine something preceded by its quotation?
Achilles: I guess I can conjure up an image of Chairman Mao walking into a banquet
  room in which there already hangs a large banner with some of his own writing on it.
  Here would be Chairman Mao, preceded by his quotation.
Tortoise: A most imaginative example. But suppose we restrict the word



Air on G's String                                                                       441

  "preceded" to the idea of precedence on a printed sheet, rather than elaborate entries
  into a banquet room.
Achilles: All right. But what exactly do you mean by "quotation" here? Tortoise: When
  you discuss a word or a phrase, you conventionally put it in quotes. For example, I can
  say,

                           The word "philosopher" has five letters.

  Here, I put "philosopher" in quotes to show that I am speaking about the WORD
  "philosopher" rather than about a philosopher in the flesh. This is called the USE-
  MENTION distinction.
Achilles: Oh?
Tortoise: Let me explain. Suppose I were to say to you,

                            Philosophers make lots of money.

  Here, I would be USING the word to manufacture an image in your mind of a twinkle-
  eyed sage with bulging moneybags. But when I put this word-or any word-in quotes, I
  subtract out its meaning and connotations, and am left only with some marks on paper,
  or some sounds. That is called "MENTION". Nothing about the word matters, other
  than its typographical aspects-any meaning it might have is ignored.
Achilles: It reminds me of using a violin as a fly swatter. Or should I say mentioning"?
  Nothing about the violin matters, other than its solidity-any meaning or function it
  might have is being ignored. Come to think of it, I guess the fly is being treated that
  way, too.
Tortoise: Those are sensible, if slightly unorthodox, extensions of the use-mention
  distinction. But now, I want you to think about preceding something by its own
  quotation.
Achilles: All right. Would this be correct?

                                    "HUBBA" HUBBA

Tortoise: Good. Try another.
Achilles: All right.

      "'PLOP' IS NOT THE TITLE OF ANY BOOK. SO FAR AS I KNOW"'
       'PLOP' IS NOT THE TITLE OF ANY BOOK, SO FAR AS I KNOW.

Tortoise: Now this example can be modified into quite an interesting specimen, simply
  by dropping `Plop'. Achilles: Really? Let me see what you mean. It becomes

      "IS NOT THE TITLE OF ANY BOOK, SO FAR AS I KNOW"
       IS NOT THE TITLE OF ANY BOOK, SO FAR AS I KNOW.

Tortoise: You see, you have made a sentence.
Achilles: So I have. It is a sentence about the pjrase “is not the toitle of any book, as far
  as I know”, and quite a silly one too.
Tortoise: Why silly?


Air on G's String                                                                        442

Achilles: Because it's so pointless. Here's another one for you:
                           “WILL BE BOYS" WILL BE BOYS.
  Now what does that mean? Honestly, what a silly game.
Tortoise: Not to my mind. It's very earnest stuff, in my opinion. In fact this operation of
  preceding some phrase by its quotation is so overwhelmingly important that I think I'll
  give it a name.
Achilles: You will? What name will you dignify that silly operation by?
Tortoise: I believe I'll call it "to quine a phrase", to quine a phrase.
Achilles: "Quine"? What sort of word is that?
Tortoise: A five-letter word, if I'm not in error.
Achilles: What 1 was driving at is why you picked those exact five letters in that exact
  order.
Tortoise: Oh, now I understand what you meant when you asked me "What sort of word
  is that?" The answer is that a philosopher by the name of "Willard Van Orman Quine"
  invented the operation, so I name it in his honor. However, I cannot go any further
  than this in my explanation. Why these particular five letters make up his name-not to
  mention why they occur in this particular order-is a question to which I have no ready
  answer. However, I'd be perfectly willing to go and
Achilles: . Please don't bother! I didn't really want to know everything about Quine's
  name. Anyway, now I know how to quine a phrase. It's quite amusing. Here's a quined
  phrase:
   ”IS A SENTENCE FRAGMENT" IS A SENTENCE FRAGMENT.
   It's silly but all the same I enjoy it. You take a sentence fragment, quine
   it, and lo and behold, you've made a sentence! A true sentence, in this case.
Tortoise: How about quining the phrase "is a king with without no subject”?
Achilles: A king without a subject would be---
Tortoise: -an anomaly, of course. Don't wander from the point. Let's have quines first,
  and kings afterwards!
Achilles: I'm to quine that phrase, am I? All right
   "IS A KING WITH NO SUBJECT" IS A KING WITH NO SUBJECT.
  It seems to me that it might make more sense if it said "sentence" instead of "king".
  Oh, well. Give me another!
Tortoise: All right just one more. Try this one:

   "WHEN QUINED, YIELDS A TORTOISE'S LOVE SONG"
Achilles: That should be easy ... I'd say the quining gives this:
   "WHEN QUINED, YIELDS A TORTOISE'S LOVE SONG"
   WHEN QUINED, YIELDS A TORTOISE'S LOVE SONG

   Hmmm… There´s something just a little peculiar here. Oh, I see what it is! The
     sentence is talking about itself! Do you see that?


Air on G's String                                                                      443

  Tortoise: What do you mean? Sentences can't talk.
  Achilles: No, but they REFER to things-and this one refers directly unambiguously-
    unmistakably-to the very sentence which it is! You just have to think back and
    remember what quining is all about.
  Tortoise: I don't see it saying anything about itself. Where does it say "me", or: "this
    sentence", or the like?
  Achilles: Oh, you are being deliberately thick-skulled. The beauty of it lies in just that:
    it talks about itself without having to come right out and say so!
  Tortoise: Well, as I'm such a simple fellow, could you just spell it all out for me,
  Achilles: Oh, he is such a Doubting Tortoise ... All right, let me see ... Suppose I make
    up a sentence-I'll call it "Sentence P"-with a blank in it.
  Tortoise: Such as?
  Achilles: Such as ...

     “\_\_\_\_\_\_\_\_WHEN QUINED, YIELDS A TORTOISE'S LOVE SONG".

    Now the subject matter of Sentence P depends on how you fill in the blank. But
    once you've chosen how to fill in the blank, then the subject matter is determined: it
    is the phrase which you get by QUINING the blank. Call that "Sentence Q", since it
    is produced by an act of quining.
  Tortoise: That makes sense. If the blank phrase were "is written on old jars of mustard
    to keep them fresh", then Sentence Q would have to be

     "IS WRITTEN ON OLD JARS OF MUSTARD TO KEEP THEM FRESH"
     IS WRITTEN ON OLD JARS OF MUSTARD TO KEEP THEM FRESH.

  Achilles: True, and Sentence P makes the claim (though whether it is valid or not, I do
    not know) that Sentence Q is a Tortoise's love song. In any case, Sentence P here is
    not talking about itself, but rather about Sentence Q. Can we agree on that much?
  Tortoise: By all means, let us agree-and what a beautiful song it is, too.
  Achilles: But now I want to make a different choice for the blank, namely

     : "WHEN QUINED, YIELDS A TORTOISE'S LOVE SONG".

  Tortoise: Oh, heavens, you're getting a little involved here. I hope this all isn't going to
    be too highbrow for my modest mind.
  Achilles: Oh, don't worry-you'll surely catch on. With this choice, Sentence Q
    becomes .. .

     "WHEN QUINED, YIELDS A TORTOISE'S LOVE-SONG"
     WHEN QUINED, YIELDS A TORTOISE'S LOVE-SONG.

  Tortoise: Oh, you wily old warrior you, I catch on. Now Sentence Q is just the same as
    Sentence P.
  Achilles: And since Sentence Q is always the topic of Sentence P, there is a loop now,
    P points back to itself. But you see, the self-reference is a



Air on G's String                                                                        444

    sort of accident. Usually Sentences Q and P are entirely unlike each other; but with
    the right choice for the blank in Sentence-P, quining will do this magic trick for
    you.
  Tortoise: Oh, how clever. I wonder why I never thought of that myself. Now tell me:
    is the following sentence self-referential?

     "IS COMPOSED OF FIVE WORDS" IS COMPOSED OF FIVE WORDS.

  Achilles: Hmm ... I can't quite tell. The sentence which you just gave is not really
    about itself, but rather about the phrase "is composed of five words". Though, of
    course, that phrase is PART of the sentence ...
  Tortoise: So the sentence refers to some part of itself-so what? Achilles: Well,
    wouldn't that qualify as self-reference, too?
  Tortoise: In my opinion, that is still a far cry from true-self-reference. But don't worry
    too much about these tricky matters. You'll have ample time to think about them in
    the future. Achilles: I will?
  Tortoise: Indeed you will. But for now, why don't you try quining the phrase "yields
    falsehood when preceded by its quotation"?
  Achilles: I see what you're getting at-that old obscene phone call. Quining it produces
    the following:

     "YIELDS FALSEHOOD WHEN PRECEDED BY ITS QUOTATION"
     YIELDS FALSEHOOD WHEN PRECEDED BY ITS QUOTATION.

      So this is what that caller was saying! I just couldn't make out where the quotation
      marks were as he spoke. That certainly is an obscene remark! People ought to be
      jailed for saying things like that!
  Tortoise: Why in the world?
  Achilles: It just makes me very uneasy. Unlike the earlier examples, I can't quite make
      out if it is a truth or a falsehood. And the more I think about it, the more I can't
      unravel it. It makes my head spin. I wonder what kind of a lunatic mind would
      make something like that up, and torment innocent people in the night with it?
  Tortoise: I wonder ... Well, shall we go downstairs now?
  Achilles: We needn't go down-we're at ground level already. Let's go back inside-
      you'll see. (They go into the tower, and come to a small wooden door.) We can just
      step outside right here. Follow me.
  Tortoise: Are you sure? I don't want to fall three floors and break my shell.
  Achilles: Would I fool you?
      (And he opens the door. In front of them sits, to all appearances, the same boy,
      talking to the same young woman. Achilles and Mr. T walk up what seem to be the
      same stairs they walked down to enter the tower, and find themselves in what looks
      like      just     the     same      courtyard      they     first   came      into.)
      Thank you, Mr. T, for your lucid clarification of that obscene telephone call.
Tortoise: And thank you, Achilles, for the pleasant promenade. I hope we meet again
      soon.




Air on G's String                                                                       445


                         On Formally Undecidable
                           Propositions of TNT
                          and Related Systems'
                           The Two Ideas of the "Oyster"
THIS CHAPTER'S TITLE is an adaptation of the title of Gödel’s famous 1931 paper-
"TNT" having been substituted for "Principia Mathematica". Gödel’s paper was a
technical one, concentrating on making his proof watertight and rigorous; this Chapter
will be more intuitive, and in it I will stress the two key ideas which are at the core of the
proof. The first key idea is the deep discovery that there are strings of TNT which can be
interpreted as speaking about other strings of TNT; in short, that TNT, as a language, is
capable of "introspection", or self-scrutiny. This is what comes from Gödel-numbering.
The second key idea is that the property of self scrutiny can be entirely concentrated into
a single string; thus that string's sole focus of attention is itself. This "focusing trick" is
traceable, in essence, to the Cantor diagonal method.
        In my opinion, if one is interested in understanding Gödel’s proof. in a deep way,
then one must recognize that the proof, in its essence, consists of a fusion of these two
main ideas. Each of them alone is a master stroke; to put them together took an act of
genius. If I were to choose, however, which of the two key ideas is deeper, I would
unhesitatingly pick the first one-the idea of Gödel-numbering, for that idea is related to
the whole notion of what meaning and reference are, in symbol-manipulating systems.
This is an idea which goes far beyond the confines of mathematical logic, whereas the
Cantor trick, rich though it is in mathematical consequences, has little if any relation to
issues in real life.

                             The First Idea: Proof-Pairs
Without further ado, then, let us proceed to the elaboration of the proof itself. We have
already given a fairly careful notion of what the Gödel isomorphism is about, in Chapter
IX. We now shall describe a mathematical notion which allows us to translate a statement
such as "The string 0=0 is a theorem of TNT into a statement of number theory. This will
involve the notion of proof-pairs. A proof-pair is a pair of natural numbers related in a
particular way. Here is the idea:




On Formally Undecidable Propositions                                                       446

   Two natural numbers, m and n respectively, form a TNT proof-pair if and only if m
   is the Gödel number of a TNT derivation whose bottom line is the string with
   Gödel number n.

The analogous notion exists with respect to the MIU-system, and it is a little easier on the
intuition to consider that case first. So, for a moment, let us back off from TNT-proof-
pairs, and look at MIU-proof-pairs. Their definition is parallel:

   Two natural numbers, m and n respectively, form a MIU-proof pair if and only if m
   is the Gödel number of a MIU-system derivation whose bottom line is the string
   with Gödel number n.

Let us see a couple of examples involving MIU-proof-pairs. First, let m =
3131131111301, n = 301. These values of m and n do indeed form a MIU-proof-pair,
because m is the Gödel number of the MIU-derivation

                                       MI
                                      MII
                                     MIIII
                                     MUI

whose last line is MUI, having Gödel number 301, which is n. By contrast, let m =
31311311130, and n = 30. Why do these two values not form a MIU-proof-pair? To see
the answer, let us write out the alleged derivation which m codes for:

                                       MI
                                      MII
                                      MIII
                                      MU

There is an invalid step in this alleged derivation! It is the step from the second to the
third line: from MII to MIII. There is no rule of inference in the MIU-system which
permits such a typographical step. Correspondingly-and this is most crucial-there is no
arithmetical rule of inference which carries you from 311 to 3111. This is perhaps a
trivial observation, in light of our discussion in Chapter IX, yet it is at the heart of the
Gödel isomorphism. What we do in any formal system has its parallel in arithmetical
manipulations.
        In any case, the values m = 31311311130, n = 30 certainly do not form a MIU-
proof-pair. This in itself does not imply that 30 is not a MIU-number. There could be
another value of m which forms a MIU proof-pair with 30. (Actually, we know by earlier
reasoning that MU is not a MIU-theorem, and therefore no number at all can form a
MIU-proof-pair with 30.)
        Now what about TNT proof pairs? Here are two parallel examples, one being
merely an alleged TNT proof-pair, the other being a valid TNT proof-pair. Can you spot
which is which? (Incidentally, here is where




On Formally Undecidable Propositions                                                    447

the `611' codon comes in. Its purpose is to separate the Gödel numbers of successive lines
in a TNT-derivation. In that sense, '611' serves as a punctuation mark. In the MIU-
system, the initial '3' of all lines is sufficient-no extra punctuation is needed.)

   (1) m = 626.262,636,223,123,262,111,666,611,223,123.666.111,666
         n = 123,666.111,666
   (2) m=626,262.636,223.123,262,111,666,611223,333,262.636,123.262,111,666
         n = 223,333,262,636,123,262.111,666

It is quite simple to tell which one is which, simply by translating back to the old
notation, and making some routine examinations to see

   (1) whether the alleged derivation coded for by m is actually a legitimate derivation;
   (2) if so, whether the last line of the derivation coincides with the string which n codes
            for.

Step 2 is trivial; and step 1 is also utterly straightforward, in this sense: there are no open-
ended searches involved, no hidden endless loops. Think of the examples above
involving the MIU-system, and now just mentally substitute the rules of TNT for the
MIU-system's rules, and the axioms of TNT for the MIU-system's one axiom. The
algorithm in both cases is the same. Let me make that algorithm explicit:

     Go down the lines in the derivation one by one. Mark those which are axioms.
     For each line which is not an axiom, check whether it follows by any of the
           rules of inference from earlier lines in the alleged derivation.
     If all nonaxioms follow by rules of inference from earlier lines, then you have a
           legitimate derivation; otherwise it is a phony derivation.

At each stage, there is a clear set of tasks to perform, and the number of them is quite
easily determinable in advance.

                     Proof-Pair-ness Is Primitive Recursive...
       The reason I am stressing the boundedness of these loops is, as you may have
sensed, that I am about to assert

   FUNDAMENTAL FACT 1: The property of being a proof-pair is a primitive
   recursive number-theoretical property, and can therefore be tested for by a BlooP
   program.

        There is a notable contrast to be made here with that other closely related number-
theoretical property: that of being a theorem-number. To




On Formally Undecidable Propositions                                                       448

assert that n is a theorem-number is to assert that some value of m exists which forms a
proof-pair with n. (Incidentally, these comments apply equally well to TNT and to the
MIU-system; it may perhaps help to keep both in mind, the MIU-system serving as a
prototype.) To check whether n is a theorem-number, you must embark on a search
through all its potential proof-pair "partners" m-and here you may be getting into an
endless chase. No one can say how far you will have to look to find a number which
forms a proof-pair with n as its second element. That is the whole problem of having
lengthening and shortening rules in the same system: they lead to a certain degree of
unpredictability.
         The example of the Goldbach Variation may prove helpful at this point. It is
trivial to test whether a pair of numbers (m,n) form a Tortoise pair: that is to say, both m
and n + m should be prime. The test is easy because the property of primeness is
primitive recursive: it admits of a predictably terminating test. But if we want to know
whether n possesses the Tortoise property, then we are asking, "Does any number m form
a Tortoise-pair with n as its second element?"-and this, once again, leads us out into the
wild, MU-loopy unknown.

                    ... And Is Therefore Represented in TNT
The key concept at this juncture, then, is Fundamental Fact 1 given above, for from it we
can conclude

   FUNDAMENTAL FACT 2: The property of forming a proof-pair is testable in
   BlooP, and consequently, it is represented in TNT by some formula having two
   free variables.

        Once again, we are being casual about specifying which system these proof-pairs
are relative to; it really doesn't matter, for both Fundamental Facts hold for any formal
system. That is the nature of formal systems: it is always possible to tell, in a predictably
terminating way, whether a given sequence of lines forms a proof, or not-and this carries
over to the corresponding arithmetical notions.

                              The Power of Proof-Pairs
Suppose we assume we are dealing with the MIU-system, for the sake of concreteness.
You probably recall the string we called "MUMON", whose interpretation on one level
was the statement "MU is a theorem of the MIU-system". We can show how MUMON
would be expressed in TNT, in terms of the formula which represents the notion of MIU-
proof-pairs. Let us abbreviate that formula, whose existence we are assured of by
Fundamental Fact 2, this way:

                               MIU-PROOF-PAIR {α,α´}




On Formally Undecidable Propositions                                                     449

Since it is a property of two numbers, it is represented by a formula with two free
variables. (Note: In this Chapter we shall always use austere TNT-so be careful to
distinguish between the variables a, a', a".) In order to assert "MU is a theorem of the
MIU-system", we would have to make the isomorphic statement "30 is a theorem-
number of the MIU-system", and then translate that into TNT-notation. With the aid of
our abbreviation, this is easy (remember also from Chapter VIII that to indicate the
replacement of every a' by a numeral, we write that numeral followed by "/a' 1):

       Зa:MIU-PROOF- PAIRja,SSSSSSSSSSSSSSSSSSSSSSSSSSSSSSO/a'}

Count the S's: there are 30. Note that this is a closed sentence of TNT, because one free
variable was quantified, the other replaced by a numeral. A clever thing has been done
here, by the way. Fundamental Fact 2 gave us a way to talk about proof-pairs; we have
figured out how to talk about theorem-numbers, as well: you just add an existential
quantifier in front! A more literal translation of the string above would be, "There exists
some number a that forms a MlIJ-proof-pair with 30 as its second element".
        Suppose that we wanted to do something parallel with respect to TNT-say, to
express the statement "0=0 is a theorem of TNT". We may abbreviate the formula which
Fundamental Fact 2 assures us exists, in an analogous way (with two free variables,
again):

                              TNT- PROOF- PAIR{a,a'}

(The interpretation of this abbreviated TNT-formula is: "Natural numbers a and a' form a
TNT-proof-pair.") The next step is to transform our statement into number theory,
following the MUMON-model above. The statement becomes "There exists some
number a which forms a TNT proof-pair with 666,111,666 as its second element". The
TNT-formula which expresses this is:

                   Зa:TNT-PROOF-PAI R{a,SSSSS SSSSSO/a' }
                                   many, many 5's!
                           (in fact, 666,111,666 of them)

-a closed sentence of TNT. (Let us call it "JOSHtU", for reasons to appear
momentarily.) So you see that there is a way to talk not only about the primitive recursive
notion of TNT-proof-pairs, but also about the related but trickier notion of TNT-
theorem-numbers.
        To check your comprehension of these ideas, figure out how to translate into TNT
the following statements of meta-TNT:.

   (1) 0=0 is not a theorem of TNT.
   (2) ~0=0 is a theorem of TNT.
   (3) ~0=0 is not a theorem of TNT.




On Formally Undecidable Propositions                                                   450

How do the solutions differ from the example done above, and from each other' Here are
a few more translation exercises.


(4) JOSHU is a theorem of TNT. (Call the TNT-string which expresses this
   ":METAJOSH t'".)
   (5) META-JOSH[. is a theorem of TNT. (Call the TNT-string which expresses this
   "META-META-JOSHC".)
   (6)META-META-JOSHU is a theorem of TNT
   (7)META-META- ME IA -JOSHU is a theorem of TNT

   (etc., etc.)

Example 5 shows that statements of meta-meta-TNT can be translated into TNT-notation;
example 6 does the same for meta-meta-meta-TNT, etc.
        It is important to keep in mind the difference between expressing a property, and
representing it, at this point. The property of being a TNT theorem-number, for instance,
is expressed by the formula

                             Зa:TNT- PROOF- PAI R{a,a' }

Translation: "a' is a TNT-theorem-number". However, we have no guarantee that this
formula represents the notion, for we have no guarantee that this property is primitive
recursive-in fact, we have more than a sneaking suspicion that it isn't. (This suspicion is
well warranted. The property of being a TNT-theorem-number is not primitive recursive,
and no TNT-formula can represent the property!) By contrast, the property of being a
proof-pair, by virtue of its primitive recursivity, is both expressible and representable, by
the formula already introduced.

                      Substitution Leads to the Second Idea
The preceding discussion got us to the point where we saw how TNT can "introspect" on
the notion of TNT-theoremhood. This is the essence of the first part of the proof. We
now wish to press on to the second major idea of the proof, by developing a notion which
allows the concentration of this introspection into a single formula. To do this, we need to
look at what happens to the Gödel number of a formula when you modify the formula
structurally in a simple way. In fact, we shall consider this specific modification:

                  replacement of all free variables by a specific numeral

. Below are shown a couple of examples of this operation in the left hand column, and in
the right hand column are exhibited the parallel changes in Gödel numbers.




On Formally Undecidable Propositions                                                     451

       Formula                                          Gödel number

       a=a                                              262,1 11,262

We now replace all
free variables by
the numeral for 2:
                                                 123,123,666,111.123,123,666
SSO=SSO
                                                 223,333,262,636,333,262,163,636,
3a:3a':a"=(SSa•SSa')
                                                 262,163,163,111,362,123,123,262,
                                                 236,123,123,262,163,323

We now replace all
 free variables by
the numeral for 4:
                                                 223,333,262,636,333,262,163,636,
---3a:3a':SSSSO=(SSa•SSa')                       123,123,123,123,666,111,362,123,
                                                 123,262,236,123,123,262,163,323

         An isomorphic arithmetical process is going on in the right-hand column, in
which one huge number is turned into an even huger number. The function which makes
the new number from the old one would not be too difficult to describe arithmetically, in
terms of additions, multiplications, powers of 10 and so on-but we need not do so. The
main point is this: that the relation among (1) the original Gödel number, (2) the number
whose numeral is inserted, and (3) the resulting Gödel number, is a primitive recursive
relation. That is to say, a BlooP test could be written which, when fed as input any three
natural numbers, says YES if they are related in this way, and NO if they aren't. You may
test yourself on your ability to perform such a test-and at the same time convince yourself
that there are no hidden open-ended loops to the process-by checking the following two
sets of three numbers:

(1)    362,262,112,262,163,323,111,123,123,123,123,666;
       2:
       362,123,123,666,112,123,123,666,323,111,123,123,123,123,666.

(2)    223,362,123,666,236,123,666,323,111,262,163.
       223,362,262,236,262,323,111,262,163;

As usual, one of the examples checks, the other does not. Now this relationship between
three numbers will be called the substitution relationship. Because it is primitive
recursive, it is represented by some formula of TNT having three free variables. Lets us
abbreviate that TNT – formula by the following notation
                                      SUB (a,a´,a´´)



On Formally Undecidable Propositions                                                   452

      Because this formula represents the substitution relationship, the formula shown
below must be a TNT-theorem:

                SU B{SSSSS SSSSSO/a,SSO/a',SSSSSS SSSSO/a"}
                262,111,262 S's 123,123,666,111,123,123,666 S's

(This is based on the first example of the substitution relation shown in the parallel
columns earlier in this section.) And again because the SUB formula represents the
substitution relation, the formula shown below certainly is not a TNT-theorem:

                              SU B{SSSO/a,SSO/a',S0/a"}

                                     Arithmoquining

We now have reached the crucial point where we can combine all of our disassembled
parts into one meaningful whole. We want to use the machinery of the TNT-PROOF-
PAIR and SUB formulas in some way to construct a single sentence of TNT whose
interpretation is: "This very string of TNT is not a TNT-theorem." How do we do it\%
Even at this point, with all the necessary machinery in front of us, the answer is not easy
to find.
         A curious and perhaps frivolous-seeming notion is that of substituting a formula's
own Gödel number into itself. This is quite parallel to that other curious, and perhaps
frivolous-seeming, notion of "quining" in the Air on G's String. Yet quining turned out to
have a funny kind of importance, in that it showed a new way of making a self-referential
sentence. Self reference of the Quine variety sneaks up on you from behind the first time
you see it-but once you understand the principle, you appreciate that it is quite simple and
lovely. The arithmetical version of quining-let's call it arithmoquining-will allow us to
make a TNT-sentence which is "about itself ".
         Let us see an example of arithmoquining. We need a formula with at least one
free variable. The following one will do:

                                          a=SO

This formula's Gödel number is 262,111,123,666, and we will stick this number into the
formula itself-or rather, we will stick its numeral in. Here is the result:

                               SSSSS        SSSSSO=SO
                                  262,111,123,666 S's

This new formula a asserts a silly falsity-that 262.111.123.666 equals 1: If we had begun
with the string ~a=S0 and then arthmoquined, we would have cone up with a true
statement—as you can see for yourself.
       When you arithmoquine, you are of course performing a special case




On Formally Undecidable Propositions                                                    439

of the substitution operation we defined earlier. If we wanted to speak about
arithmoquining inside TNT, we would use the formula

                                      SUB{a" a" a'}

where the first two variables are the same. This comes from the fact that we are using a
single number in two different ways (shades of the Cantor diagonal method!). The
number a" is both (1) the original Gödel number, and (2) the insertion-number. Let us
invent an abbreviation for the above formula:

                               ARITHMOQUINE{a", a'}

What the above formula says, in English, is:

      a' is the Gödel number of the formula gotten by arithmoquining the formula with
      Gödel number a".
Now the preceding sentence is long and ugly. Let's introduce a concise and elegant term
to summarize it. We'll say

       a' is the arithmoquinification of a"

to mean the same thing. For instance, the arithmoquinification of 262,111,123,666 is this
unutterably gigantic number:

                      123,123,123 123,123,123,666,111,123,666
                           262,111,123,666 copies of '1231

(This is just the Gödel number of the formula we got when we arithmoquined a=SO.) We
can speak quite easily about arithmoquining inside TNT.

                                    The Last Straw
Now if you look back in the Air on G's String, you will see that the ultimate trick
necessary for achieving self-reference in Quine's way is to quine a sentence which itself
talks about the concept of quining. It's not enough just to quine-you must quine a quine-
mentioning sentence! All right, then the parallel trick in our case must be to arithmoquine
some formula which itself is talking about the notion of arithmoquining!
        Without further ado, we'll now write that formula down, and call it G's uncle:

          -3a:3a':<TNT-PROOF-PAIR{a,a'}A.ARITHMOQUINE{a",a'}>

You can see explicitly how arithmoquinification is thickly involved in the plot, Now this
“uncle” has a Gödel number, of course, which we’ll call `u´




On Formally Undecidable Propositions                                                   440

The head and tail of u's decimal expansion, and even a teeny bit of its midsection, can be
read off directly:

               u = 223,333,262,636,333,262,163,636,212, ... ,161, ... ,213

For the rest, we'd have to know just how the formulas TNT-PROOF-PAIR and
ARITHMOQUINE actually look when written out. That is too complex, and it is quite
beside the point, in any case.
        Now all we need to do is-arithmoquine this very uncle! What this entails is
"booting out" all free variables-of which there is only one, namely a"-and putting in the
numeral for u everywhere. This gives us:

   -3a:3a':<TNT-PROOF-PAIR{a,a'}^,ARITHMOQUINE{SSS ... SSSO/a",a'}>
                                                          u S's

And this, believe it or not, is Gödel’s string, which we can call 'G'. Now there are two
questions we must answer without delay. They are

       (1) What Is G's Gödel number?
       (2) What is the interpretation of G?

Question 1 first. How did we make G? Well, we began with the uncle, and arithmoquined
it. So, by the definition of arithmoquinification, G's Gödel number is
                              the arithmoquinification of u.

Now question 2. We will translate G into English in stages, getting gradually more
comprehensible as we go along. For our first rough try, we make a pretty literal
translation:

       "There do not exist numbers a and a' such that both (1) they form a TNT-proof-
       pair. and (2) a' is the arithmoquinification of u."

Now certainly there is a number a' which is the arithmoquinification of u-so the problem
must lie with the other number, a. This observation allows us to rephrase the translation
of G as follows:

       "There is no number a that forms a TNT-proof-pair with the arithmoquinification
       of u."

(This step, which can be confusing, is explained below in more detail.) Do you see what
is happening? G is saying this:

       "The formula whose Gödel number is the arithmoquinification
       of u is not a theorem of TNT."

But-and this should come as no surprise by now-that formula is none other than G itself;
whence we can make the ultimate translation of G; as

                              “G is not a theorem of TNT.”


On Formally Undecidable Propositions                                                  441

-or if you prefer,

                              "I am not a theorem of TNT."

We have gradually pulled a high-level interpretation-a sentence of meta-TNT-out of what
was originally a low-level interpretation-a sentence of number theory.   -

                                   TNT Says "Uncle!"

The main consequence of this amazing construction has already been delineated in
Chapter IX: it is the incompleteness of TNT. To reiterate the argument:

   Is G a TNT-theorem? If so, then it must assert a truth. But what in fact does G
   assert? Its own nontheoremhood. Thus from its theoremhood would follow its
   nontheoremhood: a contradiction.
            Now what about G being a nontheorem? This is acceptable, in that it doesn't
   lead to a contradiction. But G's nontheoremhood is what G asserts-hence G asserts
   a truth. And since G is not a theorem, there exists (at least) one truth which is not a
   theorem of TNT.

   Now to explain that one tricky step again. I will use another similar example. Take this
   string:

        --3a:3a': <TORTOISE-PAIR{a,a' }.TENTH-POWER{SSO/a ", a' J >

where the two abbreviations are for strings of TNT which you can write down yourself.
TENTH-POWER{a",a'} represents the statement "a' is the tenth power of a"". The
literal translation into English is then:

   "There do not exist numbers a and a' such that both (1) they form a Tortoise-pair,
   and (2) a' is the tenth power of 2."

But clearly, there is a tenth power of 2-namely 1024. Therefore, what the string is really
saying is that

        "There is no number a that forms a Tortoise-pair with 1024"

which can be further boiled down to:

                       "1024 does not have the Tortoise property."

The point is that we have achieved a way of substituting a description of a number, rather
than its numeral, into a predicate. It depends on using one °extra quantified variable (a').
Here, it was the number 1024 that was described as the “tenth power of 2”; above it was
the number described as the arithmoquinification of a”.




On Formally Undecidable Propositions                                                     442

                  "Yields Nontheoremhood When Arithmoquined"

Let us pause for breath for a moment, and review what has been done. The best way I
know to give some perspective is to set out explicitly how it compares with the version of
the Epimenides paradox due to Quine. Here is a map:

             Falsehood                  <==>                 nontheoremhood

       quotation of a phrase            <==>

       preceding a predicate            <==>        definite term) into an open formula
           by a subject

       preceding a predicate            <==>        substituting the Gödel number of a
        by a quoted phrase                             string into an open formula

       preceding a predicate            <==>       substituting the Gödel number of an
        by itself, in quotes                       open formula into the formula itself
            ("quining")                                     ("arithmoquining")

   yields falsehood when quined         <==>                  "uncle" of G”
   (a predicate without a subject)                     the(an open formula of TNT

  "yields falsehood when quined"        <==>         the number a (the Gödel number
   (the above predicate. quoted)                        of the above open formula)

  "yields falsehood when quined"        <==>                    G itself
   yields falsehood when quined                        (sentence of TNT formed
   (complete sentence formed by                   by substituting a into the uncle, i.e.,
    quining the above predicate)                       arithmoquining the uncle)

                                Gödel’s Second Theorem

Since G's interpretation is true, the interpretation of its negation --G is false. And we
know that no false statements are derivable in TNT. Thus neither G nor its negation -G
can be a theorem of TNT. We have found a "hole" in our system-an undecidable
proposition. This has a number of ramifications. Here is one curious fact which follows
from G's undecidability: although neither G nor -G is a theorem, the formula <GV -G> is
a theorem, since the rules of the Propositional Calculus ensure that all well-formed
formulas of the form <Pv--P> are theorems.
This is one simple example where an assertion inside the system and an assertion about
the system seem at odds with each other. It makes one wonder if the system really
reflects itself accurately. Does the "reflected metamathematics" which exists inside TNT
correspond well to the metamathematics which we do? This was one of the questions
which intrigued Gödel when he wrote his paper. In particular, he was interested in
whether it was possible, in the “reflected metamathematics”, to prove TNT’s consistency.
Recall that this was a great philosophical dilemma of



On Formally Undecidable Propositions                                                  443

the day: how to prove a system consistent. Gödel found a simple way to express the
statement "TNT is consistent" in a TNT formula; and then he showed that this formula
(and all others which express the same idea) are only theorems of TNT under one
condition: that TNT is inconsistent. This perverse result was a severe blow to optimists
who expected that one could find a rigorous proof that mathematics is contradiction-free.
        How do you express the statement "TNT is consistent" inside TNT It hinges on
this simple fact: that inconsistency means that two formulas, x and x, one the negation of
the other, are both theorems. But if both x and -- x are theorems, then according to the
Propositional Calculus, all well-formed formulas are theorems. Thus, to show TNT's
consistency, it would suffice to exhibit one single sentence of TNT which can be proven
to be a nontheorem. Therefore, one way to express "TNT is consistent" is to say "The
formula -0=0 is not a theorem of TNT". This was already proposed as an exercise a few
pages back. The translation is:

       ---3a:TNT-PROOF- PAIR{a,SSSSS              SSSSSOIa'}
                                    223,666,111,666 S's

It can be shown, by lengthy but fairly straightforward reasoning, that-as long as TNT is
consistent-this oath-of-consistency by TNT is not a theorem of TNT. So TNT's powers
of introspection are great when it comes to expressing things, but fairly weak when it
comes to proving them. This is quite a provocative result, if one applies it metaphorically
to the human problem of self-knowledge.

                               TNT Is ω-Incomplete
Now what variety of incompleteness does TNT "enjoy? We shall see that TNT's
incompleteness is of the "omega" variety-defined in Chapter VIII. This means that there
is some infinite pyramidal family of strings all of which are theorems, but whose
associated "summarizing string" is a nontheorem. It is easy to exhibit the summarizing
string which is a nontheorem:

                                               u S's
     Va: 3a':<TNT-PROOF- PAIR{a,a'}nARITHMOQUINE{SSS ... SSSO/a",a'}>

To understand why this string is a nontheorem, notice that it is extremely similar to G
itself-in fact, G can be made from it in one step (viz., according to TNT's Rule of
Interchange). Therefore, if it were a theorem, so would G be. But since G isn't a theorem,
neither can this be.
         Now we want to show that all of the strings in the related pyramidal family are
theorems. We can write them own easily enough:




On Formally Undecidable Propositions                                                   444

                                                       u S's
      --3a': <TNT-PROOF- PAIR,O/a,a'} ARITHMOQUINE;SSS ... SSSO/a" a'I>
    -3a': <TNT-PROOF- PAIR) SO/a,a'} ARITHMOQUINEISSS ... SSSO/a", a'}>
    -3a': <TNT-PROOF-PAIR{SSO/a,a'} ARITHMOQUINEI5SS ... SSSO/a", a'}>
   -3a': <TNT-PROOF- PAIR}SSSO/a,a'} ARITHMOQUINE{SSS ... SSSO/a", a'}>

What does each one assert? Their translations, one by one, are:

           "0 and the arithmoquinification of u do not form a TNT-proof-pair."
           "1 and the arithmoquinification of u do not form a TNT-proof-pair."
           "2 and the arithmoquinification of u do not form a TNT-proof-pail."
           "3 and the arithmoquinification of u do not form a TNT-proof-pair."

Now each of these assertions is about whether two specific integers form a proof-pair or
not. (By contrast, G itself is about whether one specific integer is a theorem-number or
not.) Now because G is a nontheorem, no integer forms a proof-pair with G's Gödel
number. Therefore, each of the statements of the family is true. Now the crux of the
matter is that the property of being a proof-pair is primitive recursive, hence represented,
so that each of the statements in the list above, being true, must translate into a theorem
of TNT-which means that everything in our infinite pyramidal family is a theorem. And
that shows why TNT is w-incomplete.

                       Two Different Ways to Plug Up the Hole

Since G's interpretation is true, the interpretation of its negation -G is false. And, using
the assumption that TNT is consistent, we know that no false statements are derivable in
TNT. Thus neither G nor its negation -G is a theorem of TNT. We have found a hole in
our system-an undecidable proposition. Now this need be no source of alarm, if we are
philosophically detached enough to recognize what this is a symptom of. It signifies that
TNT can be extended, just as absolute geometry could be. In fact, TNT can be extended
in two distinct directions, just as absolute geometry could be. It can be extended in a
standard direction-which corresponds to extending absolute geometry in the Euclidean
direction; or, it can be extended in a nonstandard direction-which corresponds, of course,
to extending absolute geometry in the non-Euclidean direction. Now the standard type of
extension would involve

                                Adding G as a new axiom.




On Formally Undecidable Propositions                                                    445

This suggestion seems rather innocuous and perhaps even desirable, since, after all, G
asserts something true about the natural number system. But what about the nonstandard
type of extensions If it is at all parallel to the case of the parallel postulate, it must
involve

                        adding the negation of G as a new axiom.

But how can we even contemplate doing such a repugnant, hideous thing? After all, to
paraphrase the memorable words of Girolamo Saccheri, isn't what --G says "repugnant to
the nature of the natural numbers'?

                                 Supernatural Numbers

I hope the irony of this quotation strikes you. The exact problem with Saccheri's approach
to geometry was that he began with a fixed notion of what was true and what was not
true, and he set out only to prove what he'd assessed as true to start with. Despite the
cleverness of his approach-which involved denying the fifth postulate, and then proving
many "repugnant" propositions of the ensuing geometry-Saccheri never entertained the
possibility of other ways of thinking about points and lines. Now we should be wary of
repeating this famous mistake. We must consider impartially, to the extent that we can,
what it would mean to add -G as an axiom to TNT. Just think what mathematics would
be like today if people had never considered adding new axioms of the following sorts:

                              3a:(a+a)=S0
                              3a:Sa=O
                              3a:(a•a)=SSO
                              3a:S(a•a) =0

While each of them is "repugnant to the nature of previously known number systems",
each of them also provides a deep and wonderful extension of the notion of whole
numbers: rational numbers, negative numbers, irrational numbers, imaginary numbers.
Such a possibility is what -G is trying to get us to open our eyes to. Now in the past, each
new extension of the notion of number was greeted with hoots and catcalls. You can hear
this particularly loudly in the names attached to the unwelcome arrivals, such as
"irrational numbers", "imaginary numbers". True to this tradition, we shall name the
numbers which -'-G is announcing to us the supernatural numbers, showing how we feel
they violate all reasonable and commonsensical notions.
        If we are going to throw -G in as the sixth axiom of TNT, we had better
understand how in the world it could coexist, in one system, with the infinite pyramidal
family we just finished discussing. To put it bluntly, -G says:

               “There exists some number which forms a TNT-proof-pair with the
                               arithmoquinification of u”




On Formally Undecidable Propositions                                                    446

-but the various members of the pyramidal family successively assert:

                                  "0 is not that number"
                                  "1 is not that number"
                                  "2 is not that number"

This is rather confusing, because it seems to be a complete contradiction (which is why it
is called "ω-inconsistency"). At the root of our confusion-much as in the case of the
splitting of geometry-is our stubborn resistance to adopt a modified interpretation for the
symbols, despite the fact that we are quite aware that the system is a modified system.
We want to get away without reinterpreting any symbols-and of course that will prove
impossible.
        The reconciliation comes when we reinterpret 3 as "There exists a generalized
natural number", rather than as "There exists a natural number". As we do this, we shall
also reinterpret V in the corresponding way. This means that we are opening the door to
some extra numbers besides the natural numbers. These are the supernatural numbers.
The naturals and supernaturals together make up the totality of generalized naturals.
        The apparent contradiction vanishes into thin air, now, for the pyramidal family
still says what it said before: "No natural number forms a TNT-proof-pair with the
arithmoquinification of u." The family doesn't say anything about supernatural numbers,
because there are no numerals for them. But now, -G says, "There exists a generalized
natural number which forms a TNT-proof-pair with the arithmoquinification of u." It is
clear that taken together, the family and -G tell us something: that there is a supernatural
number which forms a TNT-proof-pair with the arithmoquinification of u. That is all-
there is no contradiction any more. TNT+-G is a consistent system, under an
interpretation which includes supernatural numbers.
        Since we have now agreed to extend the interpretations of the two quantifiers, this
means that any theorem which involves either of them has an extended meaning. For
example, the commutativity theorem

                                  Va:da':(a+a')=(a'+a)

now tells us that addition is commutative for all generalized natural numbers-in other
words, not only for natural numbers, but also for supernatural numbers. Likewise, the
TNT-theorem which says "2 is not the square of a natural number"--

                                     -3a:(a • a)=SSO

--now tells us that 2 is not the square of a supernatural number, either. In fact,
supernatural numbers share all the properties of natural numbers, as




On Formally Undecidable Propositions                                                    447

long as those properties are given to us in theorems of TNT. In other words, everything
that can be formally proven about natural numbers is thereby established also for
supernatural numbers. This means, in particular, that supernatural numbers are not
anything already familiar to you, such as fractions, or negative numbers, or complex
numbers, or whatever. The supernatural numbers are, instead, best visualized as integers
which are greater than all natural numbers-as infinitely large integers. Here is the point:
although theorems of TNT can rule out negative numbers, fractions, irrational numbers,
and complex numbers, still there is no way to rule out infinitely large integers. The
problem is, there is no way even to express the statement "There are no infinite
quantities".
        This sounds quite strange, at first. Just exactly how big is the number which
makes a TNT-proof-pair with G's Gödel number= (Let's call it 'I . for no particular
reason.) Unfortunately, we have not got any good vocabulary for describing the sizes of
infinitely large integers, so I am afraid I cannot convey a sense of I's magnitude. But then
just how big is i (the square root of -1)? Its size cannot be imagined in terms of the sizes
of familiar natural numbers. You can't say, "Well, i is about half as big as 14, and 9/10 as
big as 24." You have to say, "i squared is -1", and more or less leave it at that. A quote
from Abraham Lincoln seems a propos here. When he was asked, "How long should a
man's legs be?" he drawled, "Long enough to reach the ground." That is more or less how
to answer the question about the size of I-it should be just the size of a number which
specifies the structure of a proof of G-no bigger, no smaller.
        Of course, any theorem of TNT has many different derivations, so you might
complain that my characterization of I is nonunique. That is so. But the parallel with 1-
the square root of -1-still holds. Namely, recall that there is another number whose square
is also minus one: -i. Now i and -i are not the same number. They just have a property in
common. The only trouble is that it is the property which defines them! We have to
choose one of them-it doesn't matter which one-and call it "i". In fact there's no way of
telling them apart. So for all we know we could have been calling the wrong one "i" for
all these centuries and it would have made no difference. Now, like i, I is also
nonuniquely defined. So you just have to think of I as being some specific one of the
many possible supernatural numbers which form TNT-proof-pairs with the
arithmoquinification of u.

                Supernatural Theorems Have Infinitely Long Derivations.

We haven't yet faced head on what it means to throw -G in as an axiom. We have said it
but not stressed it. The point is that -G asserts that G has a proof. How can a system
survive, when one of its axioms asserts that its own negation has a proof? We must be in
hot water now! Well, it is not so bad as you might think. As long as we only construct
finite proofs, we will never prove G Therefore, no calamitous collision between G and its
negative ~G will ever take place. The supernatural number –I won’t cause any disaster.




On Formally Undecidable Propositions                                                    448

However, we will have to get used to the idea that ---G is now the one which asserts a
truth ("G has a proof "), while G asserts a falsity ("G has no proof"). In standard number
theory it is the other way around-but then, in standard number theory there aren't any
supernatural numbers. Notice that a supernatural theorem of TNT-namely G-may assert a
falsity, but all natural theorems still assert truths.

                       Supernatural Addition and Multiplication

There is one extremely curious and unexpected fact about supernaturals which I would
like to tell you, without proof. (I don't know the proof either.) This fact is reminiscent of
the Heisenberg uncertainty principle in quantum mechanics. It turns out that you can
"index" the supernaturals in a simple and natural way by associating with each
supernatural number a trio of ordinary integers (including negative ones). Thus, our
original supernatural number, I, might have the index set (9,-8,3), and its successor, I + 1,
might have the index set (9,-8,4). Now there is no unique way to index the supernaturals;
different methods offer different advantages and disadvantages. Under some indexing
schemes, it is very easy to calculate the index triplet for the sum of two supernaturals,
given the indices of the two numbers to be added. Under other indexing schemes, it is
very easy to calculate the index triplet for the product of two supernaturals, given the
indices of the two numbers to be multiplied. But under no indexing scheme is it possible
to calculate both. More precisely, if the sum's index can be calculated by a recursive
function, then the product's index will not be a recursive function; and conversely, if the
product's index is a recursive function, then the sum's index will not be. Therefore,
supernatural schoolchildren who learn their supernatural plus-tables will have to be
excused if they do not know their supernatural times-tables-and vice versa! You cannot
know both at the same time.

                               Supernaturals Are Useful ...

One can go beyond the number theory of supernaturals, and consider supernatural
fractions (ratios of two supernaturals), supernatural real numbers, and so on. In fact, the
calculus can be put on a new footing, using the notion of supernatural real numbers.
Infinitesimals such as dx and dy, those old bugaboos of mathematicians, can be
completely justified, by considering them to be reciprocals of infinitely large real
numbers! Some theorems in advanced analysis can be proven more intuitively with the
aid of "nonstandard analysis".

                                 But Are They Real?
Nonstandard number theory is a disorienting thing when you first meet up with it. But,
then, non-Euclidean geometry is also a disorienting subject. In




On Formally Undecidable Propositions                                                     449

both instances, one is powerfully driven to ask, "But which one of these two rival theories
is correct? Which is the truth?" In a certain sense, there is no answer to such a question.
(And vet, in another sense-to be discussed later-there is an answer.) The reason that there
is no answer to the question is that the two rival theories, although they employ the same
terms, do not talk about the same concepts. Therefore, they are only superficially rivals,
just like Euclidean and non-Euclidean geometries. In geometry, the words "point", "line",
and so on are undefined terms, and their meanings are determined by the axiomatic
system within which they are used.
         Likewise for number theory. When we decided to formalize TNT. we preselected
the terms we would use as interpretation words-for instance, words such as "number",
"plus", "times", and so on. By taking the step of formalization, we were committing
ourselves to accepting whatever passive meanings these terms might take on. But just like
Saccheri-we didn't anticipate any surprises. We thought we knew what the true, the real,
the only theory of natural numbers was. We didn't know that there would be some
questions about numbers which TNT would leave open, and which could therefore be
answered ad libitum by extensions of TNT heading off in different directions. Thus, there
is no basis on which to say that number theory "really" is this way or that, just as one
would be loath to say that the square root of -1 "really" exists, or "really" does not.

                       Bifurcations in Geometry, and Physicists

There is one argument which can be, and perhaps ought to be, raised against the
preceding. Suppose experiments in the real, physical world can be explained more
economically in terms of one particular version of geometry than in terms of any other.
Then it might make sense to say that that geometry is "true". From the point of view of a
physicist who wants to use the "correct" geometry, then it makes some sense to
distinguish between the "true" geometry, and other geometries. But this cannot be taken
too simplistically. Physicists are always dealing with approximations and idealizations of
situations. For instance, my own Ph.D. work, mentioned in Chapter V, was based on an
extreme idealization of the problem of a crystal in a magnetic field. The mathematics
which emerged was of a high degree of beauty and symmetry. Despite-or rather, because
of-the artificiality of the model, some fundamental features emerged conspicuously in the
graph. These features then suggest some guesses about the kinds of things that might
happen in more realistic situations. But without the simplifying assumptions which
produced my graph, there could never be such insights. One can see this kind of thing
over and over again in physics, where a physicist uses a "nonreal" situation to learn about
deeply hidden features of reality. Therefore, one should be extremely cautious in saying
that the brand of geometry which physicists might wish to use would represent “the




On Formally Undecidable Propositions                                                   450

true geometry", for in fact, physicists will always use a variety of different geometries,
choosing in any given situation the one that seems simplest and most convenient.
        Furthermore-and perhaps this is even more to the point-physicists do not study
just the 3-D space we live in. There are whole families of "abstract spaces" within which
physical calculations take place, spaces which have totally different geometrical
properties from the physical space within which we live. Who is to say, then, that "the
true geometry" is defined by the space in which Uranus and Neptune orbit around the
sun? There is "Hilbert space", where quantum-mechanical wave functions undulate; there
is "momentum space", where Fourier components dwell; there is "reciprocal space",
where wave-vectors cavort; there is "phase space", where many-particle configurations
swish; and so on. There is absolutely no reason that the geometries of all these spaces
should be the same; in fact, they couldn't possibly be the same! So it is essential and vital
for physicists that different and "rival" geometries should exist.

                     Bifurcations in Number Theory, and Bankers

So much for geometry. What about number theory? Is it also essential and vital that
different number theories should coexist with each other? If you asked a bank officer, my
guess is that you would get an expression of horror and disbelief. How could 2 and 2 add
up to anything but 4? And moreover, if 2 and 2 did not make 4, wouldn't world
economies collapse immediately under the unbearable uncertainty opened up by that
fact? Not really. First of all, nonstandard number theory doesn't threaten the age-old idea
that 2 plus 2 equals 4. It differs from ordinary number theory only in the way it deals with
the concept of the infinite. After all, every theorem of TNT remains a theorem in any
extension of TNT! So bankers need not despair of the chaos that will arrive when
nonstandard number theory takes over.
        And anyway, entertaining fears about old facts being changed betrays a
misunderstanding of the relationship between mathematics and the real world.
Mathematics only tells you answers to questions in the real world after you have taken
the one vital step of choosing which kind of mathematics to apply. Even if there were a
rival number theory which used the symbols `2', `3', and `+', and in which a theorem said
"2 + 2 = 3", there would be little reason for bankers to choose to use that theory! For that
theory does not fit the way money works. You fit your mathematics to the world, and not
the other way around. For instance, we don't apply number theory to cloud systems,
because the very concept of whole numbers hardly fits. There can be one cloud and
another cloud, and they will come together and instead of there being two clouds, there
will still only be one. This doesn't prove that 1 plus 1 equals 1; it just proves that our
number theoretical concept of “one” is not applicable in its full power to cloud counting.




On Formally Undecidable Propositions                                                     451

              Bifurcations in Number Theory, and Metamathematicians

So bankers, cloud-counters, and most of the rest of us need not worry ,about the advent of
supernatural numbers: they won't affect our everyday perception of the world in the
slightest. The only people who might actually be a little worried are people whose
endeavors depend in some crucial way on the nature of infinite entities. There aren't too
many such people around-but mathematical logicians are members of this category. How
can the existence of a bifurcation in number theory affect them Well, number theory
plays two roles in logic: (1) when axiomatized, it is an object of study; and (2) when used
informally, it is an indispensable tool with which formal systems can be investigated.
This is the use-mention distinction once again, in fact: in role (1), number theory is
mentioned, in role (2) it is used.
        Now mathematicians have judged that number theory is applicable to the study of
formal systems even if not to cloud-counting, just as bankers have judged that the
arithmetic of real numbers is applicable to their transactions. This is an
extramathematical judgement, and shows that the thought processes involved in doing
mathematics, just like those in other areas, involve "tangled hierarchies" in which
thoughts on one level can affect thoughts on any other level. Levels are not cleanly
separated, as the formalist version of what mathematics is would have one believe.
        The formalist philosophy claims that mathematicians only deal with abstract
symbols, and that they couldn't care less whether those symbols have any applications to
or connections with reality. But that is quite a distorted picture. Nowhere is this clearer
than in metamathematics. If the theory of numbers is itself used as an aid in gaining
factual knowledge about formal systems, then mathematicians are tacitly showing that
they believe these ethereal things called "natural numbers" are actually part of reality not
just figments of the imagination. This is why I parenthetically remarked earlier that, in a
certain sense, there is an answer to the question of which version of number theory is
"true". Here is the nub of the matter: mathematical logicians must choose which version
of number theory to put their faith in. In particular, they cannot remain neutral on the
question of the existence or nonexistence of supernatural numbers, for the two different
theories may give different answers to questions in metamathematics.
        For instance, take this question: "Is -G finitely derivable in TNT?" No one
actually knows the answer. Nevertheless, most mathematical logicians would answer no
without hesitation. The intuition which motivates that answer is based on the fact that if -
G were a theorem, TNT would be w-inconsistent, and this would force supernaturals
down your throat if you wanted to interpret TNT meaningfully-a most unpalatable
thought for most people. After all, we didn't intend or expect supernaturals to be part of
TNT when we invented it. That is, we-or most of us-believe that it is possible to make a
formalization of number theory which does not force you into believing that supernatural
numbers are every bit as real as naturals. It is that intuition about reality which
determines which “fork” of number theory mathematicians will put their faith in, when
the chips are




On Formally Undecidable Propositions                                                    452

down. But this faith may be wrong. Perhaps every consistent formalization of number
theory which humans invent will imply the existence of supernaturals, by being co-
inconsistent. This is a queer thought, but it is conceivable.
       If this were the case-which I doubt, but there is no disproof available-then G
would not have to be undecidable. In fact, there might be no undecidable formulas of
TNT at all. There could simply be one unbifurcated theory of numbers-which necessarily
includes supernaturals. This is not the kind of thing mathematical logicians expect, but it
is something which ought not to be rejected outright. Generally, mathematical logicians
believe that TNT-and systems similar to it-are ω-consistent, and that the Gödel string
which can be constructed in any such system is undecidable within that system. That
means that they can choose to add either it or its negation as an axiom.

                          Hilbert's Tenth Problem and the Tortoise

        I would like to conclude this Chapter by mentioning one extension of Gödel’s
Theorem. (This material is more fully covered in the article "Hilbert's Tenth Problem" by
Davis and Hersh, for which see the Bibliography.) For this, I must define what a
Diophantine equation is. This is an equation in which a polynomial with fixed integral
coefficients and exponents is set to 0. For instance,

                                                    a=0

                                                    and

                                               5x+13y-1=0

                                                    And

                                        5p5 + 17q17 - 177 = 0

                                                    and

                         a123,666,111,666 + b123,.666,111,666 - c123,666, 111,666 = 0

are Diophantine equations. It is in general a difficult matter to know whether a given
Diophantine equation has any integer solutions or not. In fact, in a famous lecture at the
beginning of the century, Hilbert asked mathematicians to look for a general algorithm by
which one could determine in a finite number of steps if a given Diophantine equation
has integer solutions or not. Little did he suspect that no such algorithm exists!




On Formally Undecidable Propositions                                                    453

Now for the simplification of G. It has been shown that whenever you have a sufficiently
powerful formal number theory and a Gödel-numbering for it, there is a Diophantine
equation which is equivalent to G. The equivalence lies in the fact that this equation,
when interpreted on a metamathematical level, asserts of itself that it has no solutions.
Turn it around: if you found a solution to it, you could construct from it the Gödel
number of a proof in the system that the equation has no solutions! This is what the
Tortoise did in the Prelude, using Fermat's equation as his Diophantine equation. It is
nice to know that when you do this, you can retrieve the sound of Old Bach from the
molecules in the air!




On Formally Undecidable Propositions                                                 454

                           Birthday Cantatatata . .
   One (tine May day, the Tortoise and Achilles meet, wandering in the woods.
   The latter, all decked out handsomely, is doing a jiggish sort of thing to a
   tune which he himself is humming. On his vest he is wearing a great big
   button with the words "Today is my Birthday!"

Tortoise: Hello there, .Achilles. What makes you so joyful today? Is it your birthday, by
   any chance?
Achilles: Yes, yes! Yes it is, today is my birthday!
Tortoise: That is what I had suspected, on account of that button which you are wearing,
   and also because unless I am mistaken, you are singing a tune from a Birthday
   Cantata by Bach, one written in 1727 for the fifty-seventh birthday of Augustus, King
   of Saxony.
Achilles: You're right. And Augustus' birthday coincides with mine, so THIS Birthday
   Cantata has double meaning. However, I shan't tell you my age.
Tortoise: Oh, that's perfectly all right. However, I would like to know one other thing.
   From what you have told me so far, would it be correct to conclude that today is your
   birthday?
Achilles: Yes, yes, it would be. Today is my birthday.
Tortoise: Excellent. That's just as I suspected. So now, I WILL conclude it is your
   birthday, unless
Achilles: Yes-unless what?
Tortoise: Unless that would be a premature or hasty conclusion to draw, you know.
   Tortoises don't like to jump to conclusions, after all. (We don't like to jump at all, but
   especially not to conclusions.) So let me just ask you, knowing full well of your
   fondness for logical thought, whether it would be reasonable to deduce logically from
   the foregoing sentences, that today is in fact your birthday.
Achilles: I do believe I detect a pattern to your questions, Mr. T. But rather than jump to
   conclusions myself, I shall take your question at face value, and answer it
   straightforwardly. The answer is: YES.
Tortoise: Fine! Fine! Then there is only one more thing I need to know, to be quite
   certain that today is
Achilles: Yes, yes, yes, yes ... I can already see the line of your questioning, Mr. T. I'll
   have you know that I am not so gullible as I was when we discussed Euclid's proof, a
   while back.
Tortoise: Why, who would ever have thought you to be gullible? Quite to the contrary, I
   regard you as an expert in the forms of logical thought, an authority in the science of
   valid deductions. a fount of knowledge about certain correct methods of reasoning. . .
   To tell the truth, Achilles, you are, in my opinion, a veritable titan in the art of rational
   cogitation.




Birthday Cantatatata . .                                                                   455

And it is only for that reason that I would ask you, "Do the foregoing sentences present
   enough evidence that I should conclude without further puzzlement that today is your
   birthday
Achilles: You flatten me with your weighty praise, Mr. T-FLATTER, I mean. But I am
   struck by the repetitive nature of your questioning and in my estimation, you, just as
   well as I, could have answered `yes' each time.
Tortoise: Of course I could have, Achilles. But you see, to do so would have been to
   make a Wild Guess-and Tortoises abhor Wild Guesses. Tortoises formulate only
   Educated Guesses. Ah, yes-the power of the Educated Guess. You have no idea how
   many people fail to take into account all the Relevant Factors when they're guessing.
Achilles: It seems to me that there was only one relevant factor in this rigmarole, and that
   was my first statement.
Tortoise: Oh, to be sure, it's at least ONE of the factors to take into account, I'd say-but
   would you have me neglect Logic, that venerated science of the ancients? Logic is
   always a Relevant Factor in making Educated Guesses, and since I have with me a
   renowned expert in Logic, I thought it only Logical to take advantage of that fact, and
   confirm my hunches, by directly asking him whether my intuitions were correct. So
   let me finally come out and ask you point blank: "Do the preceding sentences allow
   me to conclude, with no room for doubt, that Today is your Birthday?"
Achilles: For one more time, YES. But frankly speaking, I have the distinct impression
   that you could have supplied that answer-as well as all the previous ones-yourself.
Tortoise: How your words sting! Would I were so wise as your insinuation suggests! But
   as merely a mortal Tortoise, profoundly ignorant and longing to take into account all
   the Relevant Factors, I needed to know the answers to all those questions.
Achilles: Well then, let me clear the matter up for once and for all. The answer to all the
   previous questions, and to all the succeeding ones which you will ask along the same
   line, is just this: YES.
Tortoise: Wonderful! In one fell swoop, you have circumvented the whole mess, in your
   characteristically inventive manner. I hope you won't mind if I call this ingenious
   trick an ANSWER SCHEMA. It rolls up yes-answers numbers 1, 2, 3, etc., into one
   single ball. In fact, coming as it does at the end of the line, it deserves the title
   "Answer Schema Omega", `w' being the last letter of the Greek alphabet-as if YOU
   needed to be told THAT!
Achilles: I don't care what you call it. I am just very relieved that you finally agree that it
   is my birthday, and we can go on to some other topic-such as what you are going to
   give me as a present.
Tortoise: Hold on—not so fast. I WILL agree it is your birthday, provided on thing
Achilles: What? That I Ask for no present?




Birthday Cantatatata . .                                                                   456

Tortoise: Not at all. In fact, Achilles, I am looking forward to treating you to a fine
   birthday dinner, provided merely that I am convinced that knowledge of all those yes-
   answers at once (as supplied by Answer Schema w) allows me to proceed directly and
   without any further detours to the conclusion that today is your birthday. That's the
   case, isn't it?
Achilles: Yes, of course it is.
Tortoise: Good. And now I have yes-answer ω + 1. Armed with it, I can proceed to
   accept the hypothesis that today is your birthday, if it is valid to do so. Would you be
   so kind as to counsel me on that matter, Achilles?
Achilles: What is this? I thought I had seen through your infinite plot. Now doesn't yes-
   answer ω + 1 satisfy you? All right. I'll give you not only yes-answer ω + 2, but also
   yes-answers ω + 3, ω + 4, and so on.
Tortoise: How generous of you, Achilles. And here it is your birthday, when I should be
   giving YOU presents instead of the reverse. Or rather, I SUSPECT it is your birthday.
   I guess I can conclude that it IS your birthday, now, armed with the new Answer
   Schema, which I will call "Answer Schema 2ω ". But tell me, Achilles: Does Answer
   Schema 2ω REALLY allow me to make that enormous leap, or am I missing
   something?
Achilles: You won't trick me any more, Mr. T. I've seen the way to end this silly game. I
   hereby shall present you with an Answer Schema to end all Answer Schemas! That is,
   I present you simultaneously with Answer Schemas ω, 2 ω, 3 ω, 4 ω, 5 ω, etc. With
   this Meta-Answer-Schema, I have JUMPED OUT of the whole system, kit and
   caboodle, transcended this silly game you thought you had me trapped in-and now we
   are DONE!
Tortoise: Good grief! I feel honored, Achilles, to be the recipient of such a powerful
   Answer Schema. I feel that seldom has anything so gigantic been devised by the mind
   of man, and I am awestruck by its power. Would you mind if I give a name to your
   gift?
Achilles: Not at all.
Tortoise: Then I shall call it "Answer Schema ω2". And we can shortly proceed to other
   matters-as soon as you tell me whether the possession of Answer Schema ω2 allows
   me to deduce that today is your birthday.
Achilles: Oh, woe is me! Can't I ever reach the end of this tantalizing trail? What comes
   next?
Tortoise: Well, after Answer Schema ω2 there's answer ω2 + 1. And then answer ω2 + 2.
   And so forth. But you can wrap those all together into a packet, being Answer
   Schema ω2 + ω. And then there are quite a few other answer-packets, such as ω2 + 2ω,
   and ω2 + 3ω…. Eventually you come to Answer Schema 2ω2, and after a while,
   Answer Schemas 3ω2 and 4ω2. Beyond them there




Birthday Cantatatata . .                                                               457

   are yet further Answer Schemas, such as ω3;, ω4, ω5, and so on. It goes on quite a
   ways, you know.
Achilles: I can imagine, I suppose it comes to Answer Schema are yet further Answer
   Schemas, such as w;, w4, w5, and so on. It goes on quite a ways, you know.
Achilles: I can imagine, I suppose it comes to Answer Schema ωω after a while.
Tortoise: Of course.
Achilles: And then ωωω, and ωωωω',
Tortoise: You're catching on mighty fast, Achilles. I have a suggestion for you, if you
   don't mind. Why don't you throw all of those together into a single Answer Schema?
Achilles: All right, though I'm beginning to doubt whether it will do any good.
Tortoise: It seems to me that within our naming conventions as so far set up, there is no
   obvious name for this one. So perhaps we should just arbitrarily name it Answer
   Schema Œo.
Achilles: Confound it all! Every time you give one of my answers a NAME, it seems to
   signal the imminent shattering of my hopes that that answer will satisfy you. Why
   don't we just leave this Answer Schema nameless?
Tortoise: We can hardly do that, Achilles. We wouldn't have any way to refer to it
   without a name. And besides, there is something inevitable and rather beautiful about
   this particular Answer Schema. It would be quite ungraceful to leave it nameless! And
   you wouldn't want to do something lacking in grace on your birthday, would you? Or
   is it your birthday? Say, speaking of birthdays, today is MY' birthday!
Achilles:       It is?
Tortoise: Yes, it is. Well, actually, it's my uncle's birthday, but that's almost the same.
   How would you like to treat me to a delicious birthday dinner this evening?
Achilles: Now just a cotton-picking minute, Mr. T. Today is MY birthday. You should do
   the treating!
Tortoise: Ah, but you never did succeed in convincing me of the veracity of that remark.
   You kept on beating around the bush with answers, Answer Schemas, and whatnot.
   All I wanted to know was if it was your birthday or not, but you managed to befuddle
   me entirely. Oh, well, too bad. In any case, I'll be happy to let you treat me to a
   birthday dinner this evening.
Achilles: Very well. I know just the place. They have a variety of delicious soups. And I
   know exactly what kind we should have ...




Birthday Cantatatata . .                                                               458


                        Jumping out of the System
                             A More Powerful Formal System

ONE OF THE things which a thoughtful critic of Gödel’s proof might do would be to
examine its generality. Such a critic might, for example, suspect that Gödel has just
cleverly taken advantage of a hidden defect in one particular formal system, TNT. If this
were the case, then perhaps a formal system superior to TNT could be developed which
would not be subject to the Gödelian trick, and Gödel’s Theorem would lose much of its
sting. In this Chapter we will carefully scrutinize the properties of TNT which made it
vulnerable to the arguments of last Chapter.
         A natural thought is this: If the basic trouble with TNT is that it contains a "hole"-
in other words, a sentence which is undecidable, namely G-then why not simply plug up
the hole? Why not just tack G onto TNT as a sixth axiom? Of course, by comparison to
the other axioms, G is a ridiculously huge giant, and the resulting system-TNT+G-would
have a rather comical aspect due to the disproportionateness of its axioms. Be that as it
may, adding G is a reasonable suggestion. Let us consider it done. Now, it is to be hoped,
the new system, TNT+G, is a superior formal system-one which is not only supernatural-
free, but also complete. It is certain that TNT+G is superior to TNT in at least one
respect: the string G is no longer undecidable in this new system, since it is a theorem.
         What was the vulnerability of TNT due to? The essence of its vulnerability was
that it was capable of expressing statements about itself-in particular, the statement

                       "I Cannot Be Proven in Formal System TNT"
or, expanded a bit,

      "There does not exist a natural number which forms a TNT-proof-pair with
      the Gödel number of this string."

Is there any reason to expect or hope that TNT+G would be invulnerable to Gödel’s
proof? Not really. Our new system is just as expressive as TNT. Since Gödel’s proof
relies primarily on the expressive power of a formal system, we should not be surprised
to see our new system succumb,




Jumping out of the System                                                                  459

too. The trick will be to find a string which expresses the statement

                    "I Cannot Be Proven in Formal System TNT+G."

Actually, it is not much of a trick, once you have seen it done for TNT. All the same
principles are employed: only, the context shifts slightly. (Figuratively speaking, we take
a tune we know and simply sing it again, only in a higher key.) As before, the string
which we are looking for-let us call it "G"'-is constructed by the intermediary of an
"uncle", But instead of being based on the formula which represents TNT-proof-pairs, it
is based on the similar but slightly more complicated notion of TNT+G-proofpairs. This
notion of TNT+G-proof-pairs is only a slight extension of the original notion of' TNT-
proof-pairs.
        A similar extension could be envisaged for the MIU-system. We have seen the
unadulterated form of MIU-proof-pairs. Were we now to add MU as a second axiom, we
would be dealing with a new system-the MIU+MU system. A derivation in this extended
system is presented:

                              MU              axiom
                              MUU             rule 2

There is a MIU+MU-proof-pair which corresponds-namely, m = 30300, n = 300. Of
course, this pair of numbers does not form a MIU-proof-pair-only a MIU+MU-proof-
pair. The addition of an extra axiom does not substantially complicate the arithmetical
properties of proof-pairs. The significant fact about them-that being a proof-pair is
primitive recursive-is preserved.

                             The Gödel Method Reapplied

Vow, returning to TNT+G, we will find a similar situation. TNT+G proof-pairs, like
their predecessors, are primitive recursive, so they are represented inside TNT+G by a
formula which we abbreviate in an obvious manner.

                             (TNT+G)-PROOF-PAIR{a,a' }

Vow we just do everything all over again. We make the counterpart of G by beginning
with an "uncle", just as before:

         3a:3a':<(TNT+G)-PROOF-PAIR{a,a'}ARITHMOQUINE{a",a'}>

.et us say its Gödel-number is u'. Now we arithmoquine this very uncle. That will give us
G':

                        3a:3a': < (TNT+G)-PROOF-PAIR{a,a' }
                        ARITHMOQUINE{SSS….SSSo/a´´,a´}>
                                               U´ S´s



Jumping out of the System                                                              460

Its interpretation is

                  "There is no number a that forms a TNT +G-proof-pair
                            with the arithmoquinification of u'."
More concisely,

                        "I Cannot Be Proven in Formal System TNT+G."

                                      Multifurcation
Well (yawn), the details are quite boring from here on out. G' is precisely to TNT+G as
G was to TNT Itself. One finds that either G' or -G' can be added to TNT+G, to yield a
further splitting of number theory. And, lest you think this only happens to the "good
guys", this very same dastardly trick can be played upon TNT+--G-that is, upon the
nonstandard extension of TNT gotten by adding G's negation. So now we see (Fig. 75)
that there are all sorts of bifurcations in number theory:




FIGURE 75. "Multifurcation" of TNT. Each extension of TNT has its very own Gödel
sentence; that sentence, or its negation, can be added on, so that from each extension
there sprouts a pair of further extensions, a process which goes on ad infinitum.

Of course, this is just the beginning. Let us imagine moving down the leftmost branch of
this downwards-pointing tree, where we always toss in the Gödel sentences (rather than
their negations). This is the best we can do by way of eliminating supernaturals. After
adding G, we add G'. Then we add G", and G"', and so on. Each time we make a new
extension of TNT, its vulnerability to the Tortoise's method-pardon me, I mean Gödel’s
method.. allows a new string to be devised, having the interpretation.

                            “I cannot be proven in formal system X”




Jumping out of the System                                                                461

Naturally, after a while, the whole process begins to seem utterly predictable and routine.
Why, all the "holes" are made by one single technique! This means that, viewed as
typographical objects, they are all cast from one single mold, which in turn means that
one single axiom schema suffices to represent all of them! So if this is so, why not plug
up all :he holes at once and be done with this nasty business of incompleteness 3nce and
for all? This would be accomplished by adding an axiom schema to TNT, instead of just
one axiom at a time. Specifically, this axiom schema would be the mold in which all of
G, G', G", G"', etc., are cast. By adding :his axiom schema (let's call it "G~"), we would
be outsmarting the "Gödelization" method. Indeed, it seems quite clear that adding G. to
TNT would :)e the last step necessary for the complete axiomatization of all of number-
theoretical truth.
         It was at about this point in the Contracrostipunctus that the Tortoise related the
Crab's invention of "Record Player Omega". However, readers were left dangling as to
the fate of that device, since before completing his tale, the tuckered-out Tortoise decided
that he had best go home to sleep; but not before tossing off a sly reference to Gödel’s
Incompleteness Theorem). Now, at last, we can get around to clearing up that dangling
detail ... Perhaps you already have an inkling, after reading the Birthday Cantatatata.

                                Essential Incompleteness

As you probably suspected, even this fantastic advance over TNT suffers the same fate.
And what makes it quite weird is that it is still for, in essence, the same reason. The
axiom schema is not powerful enough, and the Gödel construction can again be effected.
Let me spell this out a little. (One can do it much more rigorously than I shall here.) If
there is a way of capturing the various strings G, G', G", G"' . . in a single typographical
mold, then there is a way of describing their Gödel numbers in a single arithmetical mold.
And this arithmetical portrayal of an infinite class of numbers can then be represented
inside TNT+G. by some formula OMEGA-AXIOM{a} whose interpretation is: "a is the
Godel number of one of the axioms coming from G.". When a is replaced by any specific
numeral, the formula which results will be a theorem of TNT+G. if and only if the
numeral stands for the Gödel number of an axiom coming from the schema.
        With the aid of this new formula, it becomes possible to represent even such a
complicated notion as TNT+G.-proof-pairs inside TNT+Gω:

                            (TNT+G.)- PROOF- PAIR{a, a' )

sing this formula, we can construct a new uncle, which we proceed to Arithmoquine in
the by now thoroughly familiar way, making yet another undecidable string, which will
be called "TNT+Gω+i". At this point, you might well wonder, "Why isn't Gω+i among
the axioms created by the axiom schema Gω?” The answer is that G was not clever
enough to foresee its own embeddability inside number theory.




Jumping out of the System                                                               462

In the Contracrostipunctus, one of the essential steps in the Tortoise's making an
"unplayable record" was to get a hold of a manufacturer's blueprint of the record player
which he was out to destroy. This was necessary so that he could figure out to what kinds
of vibrations it was vulnerable, and then incorporate into his record such grooves as
would code for sounds which would induce those vibrations. It is a close analogue to the
Gödel trick, in which the system's own properties are reflected inside the notion of proof-
pairs, and then used against it. Any system, no matter how complex or tricky it is, can be
Gödel-numbered, and then the notion of its proof-pairs can be defined-and this is the
petard by which it is hoist. Once a system is well-defined, or "boxed", it becomes
vulnerable.
        This principle is excellently illustrated by the Cantor diagonal trick, which finds
an omitted real number for each well-defined list of reals between 0 and 1. It is the act of
giving an explicit list-a "box" of reals which causes the downfall. Let us see how the
Cantor trick can be repeated over and over again. Consider what happens if, starting with
some list L, you do the following:

       (la) Take list L, and construct its diagonal number d.
       (lb) Throw d somewhere into list L, making a new list L+d.

       (2a) Take list L +d, and construct its diagonal number d'.
       (2b) Throw d' somewhere into list L+d, making a new list L+d+d'.

Now this step-by-step process may seem a doltish way to patch up L, for we could have
made the entire list d, d', d", d"', ... at once, given L originally. But if you think that
making such a list will enable you to complete your list of reals, you are very wrong. The
problem comes at the moment you ask, "Where to incorporate the list of diagonal
numbers inside L?" No matter how diabolically clever a scheme you devise for
ensconcing the d-numbers inside L, once you have done it, then the new list is still
vulnerable. As was said above, it is the act of giving an explicit list-a "box" of reals-that
causes the downfall.
        Now in the case of formal systems, it is the act of giving an explicit recipe for
what supposedly characterizes number-theoretical truth that causes the incompleteness.
This is the crux of the problem with TNT+Gω,. Once you insert all the G's in a well-
defined way into TNT, there is seen to be some other G-some unforeseen G-which you
didn't capture in your axiom schema. And in the case of the TC-battle inside the
ContracrostiPunctus, the instant a record player's "architecture" is determined, the record
player becomes capable of being shaken to pieces.
        So what is to be done? There is no end in sight. It appears that TNT, even when
extended ad infinitum, cannot be made complete. TNT is therefore said to suffer from
essential incompleteness because the income-




Jumping out of the System                                                                463

pleteness here is part and parcel of TNT; it is an essential part of the nature of TNT and
cannot be eradicated in any way, whether simpleminded or ingenious. What's more, this
problem will haunt any formal version of number theory, whether it is an extension of
TNT, a modification of TNT, or an alternative to TNT. The fact of the matter is this: the
possibility of constructing, in a given system, an undecidable string via Gödel’s self-
reference method, depends on three basic conditions:

   (1) That the system should be rich enough so that all desired statements about
       numbers, whether true or false, can be expressed in it. (Failure on this count
       means that the system is from the very start too weak to be counted as a rival to
       TNT, because it can't even express number-theoretical notions that TNT can.
       In the metaphor of the Contracrosttpunctus, it is as if one did not have a
       phonograph but a refrigerator or some other kind of object.)
   (2) That all general recursive relations should be represented by formulas in the
       system. (Failure on this count means the system fails to capture in a theorem
       some general recursive truth, which can only be considered a pathetic bellyflop
       if it is attempting to produce all of number theory's truths. In the
       Contracrostipunctus metaphor, this is like having a record player, but one of
       low fidelity.)
   (3) That the axioms and typographical patterns defined by its rules be recognizable
       by some terminating decision procedure. (Failure on this count means that there
       is no method to distinguish valid derivations in the system from invalid ones-
       thus that the "formal system" is not formal after all, and in fact is not even well-
       defined. In the Contracrostipunctus metaphor, it is a phonograph which is still
       on the drawing board, only partially designed.)

Satisfaction of these three conditions guarantees that any consistent system will be
incomplete, because Gödel’s construction is applicable.
        The fascinating thing is that any such system digs its own hole; the system's own
richness brings about its own downfall. The downfall occurs essentially because the
system is powerful enough to have self-referential sentences. In physics, the notion exists
of a "critical mass" of a fissionable substance, such as uranium. A solid lump of the
substance will just sit there, if its mass is less than critical. But beyond the critical mass,
such a lump will undergo a chain reaction, and blow up. It seems that with formal
systems there is an analogous critical point. Below that point, a system is "harmless" and
does not even approach defining arithmetical truth formally; but beyond the critical point,
the system suddenly attains the capacity for self-reference, and thereby dooms itself to
incompleteness. The threshold seems to be roughly when a system attains the three
properties listed above.




Jumping out of the System                                                                  464

Once this ability for self-reference is attained, the system has a hole which is tailor-made
for itself; the hole takes the features of the system into account and uses them against the
system.

                         The Passion According to Lucas
The baffling repeatability of the Gödel argument has been used by various people-notably
J. R. Lucas-as ammunition in the battle to show that there is some elusive and ineffable
quality to human intelligence, which makes it unattainable by "mechanical automata"-that
is, computers. Lucas begins his article "Minds, Machines, and Gödel" with these words:

       Gödel’s theorem seems to me to prove that Mechanism is false, that is, that minds
       cannot be explained as machines.'

         Then he proceeds to give an argument which, paraphrased, runs like this. For a
computer to be considered as intelligent as a person is, it must be able to do every
intellectual task which a person can do. Now Lucas claims that no computer can do
"Gödelization" (one of his amusingly irreverent terms) in the manner that people can.
Why not? Well, think of any particular formal system, such as TNT, or TNT+G, or even
TNT+G.. One can write a computer program rather easily which will systematically
generate theorems of that system, and in such a manner that eventually, any preselected
theorem will be printed out. That is, the theorem-generating program won't skip any
portion of the "space" of all theorems. Such a program would be composed of two major
parts: (1) a subroutine which stamps out axioms, given the "molds" of the axiom schemas
(if there are any), and (2) a subroutine which takes known theorems (including axioms, of
course) and applies rules of inference to produce new theorems. The program would
alternate between running first one of these subroutines, and then the other.
         We can anthropomorphically say that this program "knows" some facts of number
theory-namely, it knows those facts which it prints out. If it fails to print out some true
fact of number theory, then of course it doesn't "know" that fact. Therefore, a computer
program will be inferior to human beings if it can be shown that humans know something
which the program cannot know. Now here is where Lucas starts rolling. He says that we
humans can always do the Gödel trick on any formal system as powerful as TNT-and
hence no matter what the formal system, we know more than it does. Now this may only
sound like an argument about formal systems, but it can also be slightly modified so that
it becomes, seemingly, an invincible argument against the possibility of Artificial
Intelligence ever reproducing the human level of intelligence. Here is the gist of it:

       Rigid internal codes entirely rule computers and robots; ergo ...
       Computers are isomorphic to formal systems. Now .. .
       Any computer which wants to be as smart as we are has got to be able to do
       number theory as well as we can, so….




Jumping out of the System                                                               465

       Among other things, it has to be able to do primitive recursive arithmetic. But for
          this very reason .. .
       It is vulnerable to the Gödelian "hook", which implies that ...
       We, with our human intelligence, can concoct a certain statement of number
          theory which is true, but the computer is blind to that statement's truth (i.e., will
          never print it out), precisely because of Gödel’s boomeranging argument.
       This implies that there is one thing which computers just cannot be programmed
          to do, but which we can do. So we are smarter.

Let us enjoy, with Lucas, a transient moment of anthropocentric glory:

  However complicated a machine we construct, it will, if it is a machine, correspond
  to a formal system, which in turn will be liable to the Gödel procedure for finding a
  formula unprovable-in-that-system. This formula the machine will be unable to
  produce as being true, although a mind can see it is true. And so the machine will
  still not be an adequate model of the mind. We are trying to produce a model of the
  mind which is mechanical-which is essentially "dead"-but the mind, being in fact
  "alive," can always go one better than any formal, ossified, dead system can. Thanks
  to Gödel’s theorem. the mind always has the last word.2

   On first sight, and perhaps even on careful analysis, Lucas' argument appears
compelling. It usually evokes rather polarized reactions. Some ;eize onto it as a nearly
religious proof of the existence of souls, while others laugh it off as being unworthy of
comment. I feel it is wrong, but Fascinatingly so-and therefore quite worthwhile taking
the time to rebut. In fact, it was one of the major early forces driving me to think over the
matters in this book. I shall try to rebut it in one way in this Chapter, and in ether ways in
Chapter XVII.
   We must try to understand more deeply why Lucas says the computer cannot be
programmed to "know" as much as we do. Basically the idea is :hat we are always
outside the system, and from out there we can always perform the "Gödelizing"
operation, which yields something which the program, from within, can't see is true. But
why can't the "Gödelizing operator", as Lucas calls it, be programmed and added to the
program as a third major component, Lucas explains:

   The procedure whereby the Gödelian formula is constructed is a standard
   procedure-only so could we be sure that a Gödelian formula can be constructed for
   every formal system. But if it is a standard procedure, then a machine should be
   able to be programmed to carry it out too.... This would correspond to having a
   system with an additional rule of inference which allowed one to add, as a theorem,
   the Gödelian formula of the rest of' the formal system, and then the Gödelian
   formula of this new, strengthened, formal system, and so on. It would be
   tantamount to adding to the original formal system an infinite sequence of axioms,
   each the Gödelian formula of the system hitherto obtained… We might expect a
   mind, faced with a machine that possessed a Gödelizing operator, to take this into
   account, and



Jumping out of the System                                                                 466

   out-Gödel the new machine, Gödelizing operator and all. This has, in fact, proved
   to be the case. Even if we adjoin to a formal system the infinite set of axioms
   consisting of the successive Gödelian formulae, the resulting system is still
   incomplete, and contains a formula which cannot be proved-in-the system, although
   a rational being can, standing outside the system, see that it is true. We had
   expected this, for even if an infinite set of axioms were added, they would have to
   be specified by some finite rule or specification, and this further rule or
   specification could then be taken into account by a mind considering the enlarged
   formal system. In a sense, just because the mind has the last word, it can always
   pick a hole in any formal system presented to it as a model of its own workings.
   The mechanical model must be, in some sense, finite and definite: and then the
   mind can always go one better.'

                              Jumping Up a Dimension
A visual image provided by M. C. Escher is extremely useful in aiding the intuition here:
his drawing Dragon (Fig. 76). Its most salient feature is, of course, its subject matter-a
dragon biting its tail, with all the Gödelian connotations which that carries. But there is a
deeper theme to this picture. Escher himself wrote the following most interesting
comments. The first comment is about a set of his drawings all of which are concerned
with "the conflict between the flat and the spatial"; the second comment is about Dragon
in particular.

  I      Our three-dimensional space is the only true reality we know. The two-
  dimensional is every bit as fictitious as the four-dimensional, for nothing is flat, not
  even the most finely polished mirror. And yet we stick to the convention that a wall
  or a piece of paper is flat, and curiously enough, we still go on, as we have done
  since time immemorial, producing illusions of space on just such plane surfaces as
  these. Surely it is a bit absurd to draw a few lines and then claim: "This is a house".
  This odd situation is the theme of the next five pictures ( Including Dragon)
  II. However much this dragon tries to be spatial, he remains completely flat. Two
  incisions are made in the paper on which he is printed. Then it is folded in such a
  way as to leave two square openings. But this dragon is an obstinate beast, and in'
  spite of his two dimensions he persists in assuming that he has three; so he sticks his
  head through one of the holes and his tail through the others 5

This second remark especially is a very telling remark. The message is that no matter
how cleverly you try to simulate three dimensions in two, you are always missing some
"essence of three-dimensionality". The dragon tries very hard to fight his two-
dimensionality. He defies the two-dimensionality of the paper on which he thinks he is
drawn, by sticking his head through it; and yet all the while, we outside the drawing can
see the pathetic futility of it all, for the dragon and the holes and the folds are all merely
two-dimensional simulations of those concepts, and not a one of them is real. But the
dragon cannot step out of his two-dimensional space, and cannot




Jumping out of the System                                                                 467

FIGURE 76. Dragon, by M. C. Escher (wood-engraving, 1952).

know it as we do. We could, in fact, carry the Escher picture any number of steps further.
For instance, we could tear it out of the book, fold it, cut holes in it, pass it through itself,
and photograph the whole mess, so that it again becomes two-dimensional. And to that
photograph, we could once again do the same trick. Each time, at the instant that it
becomes two- Matter how 'cleverly we seem to have simulated three dimensions inside
two—it becomes vulnerable to being cut and folded again.




Jumping out of the System                                                                   468

        Now with this wonderful Escherian metaphor, let us return to the program versus
the human. We were talking about trying to encapsulate the "Gödelizing operator" inside
the program itself. Well, even if we had written a program which carried the operation
out, that program would not capture the essence of Gödel’s method. For once again, we,
outside the system, could still "zap" it in a way which it couldn't do. But then are we
arguing with, or against, Lucas

                         The Limits of Intelligent Systems
Against. For the very fact that we cannot write a program to do "Gödelizing" must make
us somewhat suspicious that we ourselves could do it in every case. It is one thing to
make the argument in the abstract that Gödelizing "can be done"; it is another thing to
know how to do it in every particular case. In fact, as the formal systems (or programs)
escalate in complexity, our own ability to "Gödelize" will eventually begin to waver. It
must, since, as we have said above, we do not have any algorithmic way of describing
how to perform it. If we can't tell explicitly what is involved in applying the Gödel
method in all cases, then for each of us there will eventually come some case so
complicated that we simply can't figure out how to apply it.
         Of course, this borderline of one's abilities will be somewhat ill-defined, just as is
the borderline of weights which one can pick up off the ground. While on some days you
may not be able to pick up a 250-pound object, on other days maybe you can.
Nevertheless, there are no days whatsoever on which you can pick up a 250-ton object.
And in this sense, though everyone's Godelization threshold is vague, for each person,
there are systems which lie far beyond his ability to Godelize.
         This notion is illustrated in the Birthday Cantatatata. At first, it seems obvious
that the Tortoise can proceed as far as he wishes in pestering Achilles. But then Achilles
tries to sum up all the answers in a single swoop. This is a move of a different character
than any that has gone before, and is given the new name 'co'. The newness of the name
is quite important. It is the first example where the old naming scheme-which only
included names for all the natural numbers-had to be transcended. Then come some more
extensions, some of whose names seem quite obvious, others of which are rather tricky.
But eventually, we run out of names once again-at the point where the answer-schemas

       Ω, ωω, ωωω …..

are all subsumed into one outrageously complex answer schema. The altogether new
name 'e„' is supplied for this one. And the reason a new name is needed is that some
fundamentally new kind of step has been taken—a sort of irregularity has been
encountered. Thus a new name must be applied ad hoc.




Jumping out of the System                                                                  469

               There Is No Recursive Rule for Naming Ordinals.
Now offhand you might think that these irregularities in the progression >m ordinal to
ordinal (as these names of infinity are called) could be handled by a computer program.
That is, there would be a program to produce new names in a regular way, and when it
ran out of gas, it would invoke the "irregularity handler", which would supply a new
name, and pass control back to the simple one. But this will not work. It turns out that
irregularities themselves happen in irregular ways, and one would need o a second-order
program-that is, a program which makes new programs which make new names. And
even this is not enough. Eventually, a third-order program becomes necessary. And so on,
and so on.
        All of this perhaps ridiculous-seeming complexity stems from a deep °theorem,
due to Alonzo Church and Stephen C. Kleene, about the structure of these "infinite
ordinals", which says:

            There is no recursively related notation-system which gives a
            name to every constructive ordinal.

hat "recursively related notation-systems" are, and what "constructive ordinals" are, we
must leave to the more technical sources, such as Hartley )gets' book, to explain. But the
intuitive idea has been presented. As the ordinals get bigger and bigger, there are
irregularities, and irregularities in e irregularities, and irregularities in the irregularities in
the irregularities, etc. No single scheme, no matter how complex, can name all e ordinals.
And from this, it follows that no algorithmic method can tell w to apply the method of
Gödel to all possible kinds of formal systems. ad unless one is rather mystically inclined,
therefore one must conclude at any human being simply will reach the limits of his own
ability to 5delize at some point. From there on out, formal systems of that complex,
though admittedly incomplete for the Gödel reason, will have as much power as that
human being.

                                 Other Refutations of Lucas

Now this is only one way to argue against Lucas' position. There are others, possibly
more powerful, which we shall present later. But this counterargument has special
interest because it brings up the fascinating concept trying to create a computer program
which can get outside of itself, see itself completely from the outside, and apply the
Gödel zapping-trick to itself. Of course this is just as impossible as for a record player to
be able to ay records which would cause it to break.
But-one should not consider TNT defective for that reason. If there a defect anywhere, it
is not in TNT, but in our expectations of what it should he able to do. Furthermore, it is
helpful to realize that we are equally vulnerable to the word trick which Gödel
transplanted into mathematical formalisms: the Epimenides paradox. This was quite
cleverly pointed out




Jumping out of the System                                                                     470

by C. H. Whitely, when he proposed the sentence "Lucas cannot consistently assert this
sentence." If you think about it, you will see that (1) it is true, and yet (2) Lucas cannot
consistently assert it. So Lucas is also "incomplete" with respect to truths about the
world. The way in which he mirrors the world in his brain structures prevents him from
simultaneously being "consistent" and asserting that true sentence. But Lucas is no more
vulnerable than any of us. He is just on a par with a sophisticated formal system.
        An amusing way to see the incorrectness of Lucas' argument is to translate it into
a battle between men and women ... In his wanderings, Loocus the Thinker one day
comes across an unknown object-a woman. Such a thing he has never seen before, and at
first he is wondrous thrilled at her likeness to himself: but then, slightly scared of her as
well, he cries to all the men about him, "Behold! I can look upon her face, which is
something she cannot do-therefore women can never be like me!" And thus he proves
man's superiority over women, much to his relief, and that of his male companions.
Incidentally, the same argument proves that Loocus is superior to all other males, as well-
but he doesn't point that out to them. The woman argues back: "Yes, you can see my face,
which is something I can't do-but I can see your face, which is something you can't do!
We're even." However, Loocus comes up with an unexpected counter: "I'm sorry, you're
deluded if you think you can see my face. What you women do is not the same as what
we men do-it is, as I have already pointed out, of an inferior caliber, and does not deserve
to be called by the same name. You may call it `womanseeing'. Now the fact that you can
'womansee' my face is of no import, because the situation is not symmetric. You see?" "I
womansee," womanreplies the woman, and womanwalks away .. .
        Well, this is the kind of "heads-in-the-sand" argument which you have to be
willing to stomach if you are bent on seeing men and women running ahead of computers
in these intellectual battles.

                       Self-Transcendence-A Modern Myth
It is still of great interest to ponder whether we humans ever can jump out of ourselves-or
whether computer programs can jump out of themselves. Certainly it is possible for a
program to modify itself-but such modifiability has to be inherent in the program to start
with, so that cannot be counted as an example of "jumping out of the system". No matter
how a program twists and turns to get out of itself, it is still following the rules inherent
in itself. It is no more possible for it to escape than it is for a human being to decide
voluntarily not to obey the laws of physics. Physics is an overriding system, from which
there can be no escape. However, there is a lesser ambition which it is possible to
achieve: that is, one can certainly Jump from a subsystem of one's brain into a wider
subsystem. One can step out of ruts on occasion. This is still due to the interaction of
various subsystems of one’s brain, but it can feel very much like stepping entirely out of
oneself. Similarly, it is entirely conceivable that a partial ability to “step outside of itself”
could be embodied in a computer program.




Jumping out of the System                                                                   471

        However, it is important to see the distinction between perceiving oneself, and
transcending oneself. You can gain visions of yourself in all sorts of rays-in a mirror, in
photos or movies, on tape, through the descriptions if others, by getting psychoanalyzed,
and so on. But you cannot quite break out of your own skin and be on the outside of
yourself (modern occult movements, pop psychology fads, etc. notwithstanding). TNT
can talk about itself, but it cannot jump out of itself. A computer program can modify
itself but it cannot violate its own instructions-it can at best change some parts of itself by
obeying its own instructions. This is reminiscent of the numerous paradoxical question,
"Can God make a stone so heavy that he can’t lift it?"

                       Advertisement and Framing Devices
[his drive to jump out of the system is a pervasive one, and lies behind all progress in art,
music, and other human endeavors. It also lies behind such trivial undertakings as the
making of radio and television commercials. [his insidious trend has been beautifully
perceived and described by Irving Goffman in his book Frame Analysis:

   For example, an obviously professional actor completes a commercial pitch and,
   with the camera still on him, turns in obvious relief from his task, now to take real
   pleasure in consuming the product he had been advertising.
        This is, of course, but one example of the way in which TV and radio
   commercials are coming to exploit framing devices to give an appearance of
   naturalness that (it is hoped) will override the reserve auditors have developed.
   Thus, use is currently being made of children's voices, presumably because these
   seem unschooled; street noises, and other effects to give the impression of
   interviews with unpaid respondents; false starts, filled pauses, byplays, and
   overlapping speech to simulate actual conversation; and, following Welles, the
   interception of a firm's jingle commercials to give news of its new product,
   alternating occasionally with interception by a public interest spot, this presumably
   keeping the faith of the auditor alive.
        The more that auditors withdraw to minor expressive details as a test of
   genuineness, the more that advertisers chase after them. What results is a sort of
   interaction pollution, a disorder that is also spread by the public relations
   consultants of political figures, and, more modestly, by micro-sociology.'


Here we have yet another example of an escalating "TC-battle"-the antagonists this time
being Truth and Commercials.


                    Simplicio, Salviati, Sagredo: Why Three?
There is a fascinating connection between the problem of jumping out of ie system and
the quest for complete objectivity. When I read Jauch's four dialogues in Are Quanta
Real? based on Galileo's four Dialogues Concerning Two New Sciences, I found myself


Jumping out of the System                                                                  472

wondering why there were three characters participating. Simplico, Salviati and Sagredo.
Why wouldn’t two have




Jumping out of the System                                                           473

sufficed: Simplicio, the educated simpleton, and Salviati, the knowledgeable thinker?
What function does Sagredo have? Well, he is supposed to be a sort of neutral third party,
dispassionately weighing the two sides and coming out with a "fair" and "impartial"
judgment. It sounds very balanced, and yet there is a problem: Sagredo is always
agreeing with Salviati, not with Simplicio. How come Objectivity Personified is playing
favorites? One answer, of course, is that Salviati is enunciating correct views, so Sagredo
has no choice. But what, then, of fairness or "equal time"?
        By adding Sagredo, Galileo (and Jauch) stacked the deck more against Simplicio,
rather than less. Perhaps there should be added a yet higher level Sagredo-someone who
will be objective about this whole situation ... You can see where it is going. We are
getting into a never-ending series of "escalations in objectivity", which have the curious
property of never getting any more objective than at the first level: where Salviati is
simply right, and Simplicio wrong. So the puzzle remains: why add Sagredo at all? And
the answer is, it gives the illusion of stepping out of the system, in some intuitively
appealing sense.

                             Zen and "Stepping Out"
In Zen, too, we can see this preoccupation with the concept of transcending the system.
For instance, the koan in which Tozan tells his monks that "the higher Buddhism is not
Buddha". Perhaps, self-transcendence is even the central theme of Zen. A Zen person is
always trying to understand more deeply what he is, by stepping more and more out of
what he sees himself to be, by breaking every rule and convention which he perceives
himself to be chained by-needless to say, including those of Zen itself. Somewhere along
this elusive path may come enlightenment. In any case (as I see it), the hope is that by
gradually deepening one's self-awareness, by gradually widening the scope of "the
system", one will in the end come to a feeling of being at one with the entire universe.




Jumping out of the System                                                              474

                            Edifying Thoughts
                           of a Tobacco Smoker
                     Achilles has been invited to the Crab's home.

Achilles: I see you have made a few additions since I was last here, Mr. Crab. Your new
  paintings are especially striking.
Crab: Thank you. I am quite fond of certain painters-especially Rene Magritte. Most of
  the paintings I have are by him. He's my favorite artist.
Achilles: They are very intriguing images, I must say. In some ways, these paintings by
  Magritte remind me of works by MY favorite artist, M. C. Escher.
Crab: I can see that. Both Magritte and Escher use great realism in exploring the worlds
  of paradox and illusion; both have a sure sense for the evocative power of certain
  visual symbols, and-something which even their admirers often fail to point out-both
  of them have a sense of the graceful line.
Achilles: Nevertheless, there is something quite different about them. I wonder how one
  could characterize that difference.
Crab: It would be fascinating to compare the two in detail.
Achilles: I must say, Magritte's command of realism is astonishing. For instance, I was
  quite taken in by that painting over there of a tree with a giant pipe behind it.




                 FIGURE 77. The Shadows, by Rene Magritte (1966).




Edifying Thoughts of a Tobacco Smoker                                               475

Crab: You mean a normal pipe with a tiny tree in front of it!
Achilles: Oh, is that what it is? Well, in any case, when I first spotted it, I was convinced
  I was smelling pipe smoke! Can you imagine how silly I felt?
Crab: I quite understand. My guests are often taken in by that one.

   (So saying, he reaches up, removes the pipe from behind the tree in the painting,
   turns over and taps it against the table, and the room begins to reek of pipe tobacco.
   He begins packing in a new wad of tobacco.)

  This is a fine old pipe, Achilles. Believe it or not, the bowl has a copper lining,
  which makes it age wonderfully.
Achilles: A copper lining! You don't say!
Crab (pulls out a box of matches, and lights his pipe): Would you care for a smoke,
  Achilles?
Achilles: No, thank you. I only smoke cigars now and then.
Crab: No problem! I have one right here! (Reaches out towards another Magritte
  painting, featuring a bicycle mounted upon a lit cigar.)
Achilles: Uhh-no thank you, not now.
Crab: As you will. I myself am an incurable tobacco smoker. Which reminds me-you
  undoubtedly know of Old Bach's predilection for pipe smoking?
Achilles: I don't recall exactly.
Crab: Old Bach was fond of versifying, philosophizing, pipe smoking, and

                  FIGURE 78. State of Grace, by Rene Magritte (1959).




Edifying Thoughts of a Tobacco Smoker                                                    476

music making (not necessarily in that order). He combined all four into a droll poem
which he set to music. It can be found in the famous musical notebook he kept for his
wife, Anna Magdalena, and it is called

                       Edifying Thoughts of a Tobacco Smoker'

            Whene'er I take my pipe and stuff it
            And smoke to pass the time away,
            My thoughts, as I sit there and puff it,
            Dwell on a picture sad and gray:
            It teaches me that very like
            Am I myself unto my pipe.

            Like me, this pipe so fragrant burning
            Is made of naught but earth and clay;
            To earth I too shall be returning.
            It falls and, ere I'd think to say,
            It breaks in two before my eyes;
            In store for me a like fate lies.

            No stain the pipe's hue yet cloth darken;
            It remains white. Thus do I know
            That when to death's call I must harken
            My body, too, all pale will grow.
            To black beneath the sod 'twill turn,
            Likewise the pipe, if oft it burn.

            Or when the pipe is fairly glowing,
            Behold then, instantaneously,
            The smoke off into thin air going,
            Till naught but ash is left to see.
            Man's fame likewise away will burn
            And unto dust his body turn.

            How oft it happens when one's smoking:
            The stopper's missing from its shelf,
            And one goes with one's finger poking
            Into the bowl and burns oneself.
            If in the pipe such pain cloth dwell,
            How hot must be the pains of hell.

            Thus o'er my pipe, in contemplation
            Of such things, I can constantly
            Indulge in fruitful meditation,
            And so, puffing contentedly,
            On land, on sea, at home, abroad
            I smoke my pipe and worship God.

Edifying Thoughts of a Tobacco Smoker                                            477

  A charming philosophy, is it not?
Achilles: Indeed. Old Bach was a turner of phrases quite pleasin'.
Crab: You took the very words from my mouth. You know, in my time I have tried to
  write clever verses. But I fear mine don't measure up to much. I don't have such a way
  with words.
Achilles: Oh, come now, Mr. Crab. You have-how to put it?-quite a penchant for trick'ry
  and teasin'. I'd be honored if you'd sing me one of your songs, Mr. C.
Crab: I'm most flattered. How about if I play you a record of myself singing one of my
  efforts? I don't remember when it dates from. Its title is "A Song Without Time or
  Season".
Achilles: How poetic!

  (The Crab pulls a record from his shelves, and walks over to a huge, complex piece
  of apparatus. He opens it up, and inserts the record into an ominous-looking
  mechanical mouth. Suddenly a bright flash of greenish light sweeps over the
  surface of the record, and after a moment, the record is silently whisked into some
  hidden belly of the fantastic machine. A moment passes, and then the strains of the
  Crab's voice ring out.)

                          A turner of phrases quite pleasin',
                        Had a penchant for trick'ry and teasin'.
                               In his songs, the last line
                               Might seem sans design;
                       What I mean is, without why or wherefore.

Achilles: Lovely! Only, I'm puzzled by one thing. It seems to me your song, the last line
  is
Crab: Sans design?
Achilles: No ... What I mean is, without rhyme or reason. Crab: You could be right.
Achilles: Other than that, it's a very nice song, but I must say I am even more intrigued
  by this monstrously complex contraption. Is it merely an oversized record player?
Crab: Oh, no, it's much more than that. This is my Tortoise-chomping record player.
  Achilles: Good grief!
Crab: Well, I don't mean that it chomps up Tortoises. But it chomps up records produced
  by Mr. Tortoise.
Achilles: Whew! That's a little milder. Is this part of that weird musical battle that
  evolved between you and Mr. T some time ago?
Crab: In a way. Let me explain a little more fully. You see, Mr. Tortoise's sophistication
  had reached the point where he seemed to be able to destroy almost any record player
  I would obtain.
Achilles: But when I heard about your rivalry, it seemed to me you had at last come into
  possession of an invincible phonograph—one with a




Edifying Thoughts of a Tobacco Smoker                                                 478

built-in TV camera, minicomputer and so on, which could take itself apart and rebuild
   itself in such a way that it would not be destroyed.
Crab: Alack and alas! My plan was foiled. For Mr. Tortoise took advantage of one small
   detail which I had overlooked: the subunit which directed the disassembly and
   reassembly processes was itself stable during the entire process. That is, for obvious
   reasons, it could not take itself apart and rebuild itself, so it stayed intact.
Achilles: Yes, but what consequences did that have.
Crab: Oh, the direst ones! For you see, Mr. T focused his method down onto that subunit
   entirely.
Achilles: How is that=
Crab: He simply made a record which would induce fatal vibrations in the one structure
   he knew would never change-the disassembly reassembly subunit.
Achilles: Oh, I see ... Very sneaky.
Crab: Yes, so I thought, too. And his strategy worked. Not the first time, mind you. I
   thought I had outwitted him when my phonograph survived his first onslaught. I
   laughed gleefully. But the next time, he returned with a steely glint in his eye, and I
   knew he meant business. I placed his new record on my turntable. Then, both of us
   eagerly watched the computer-directed subunit carefully scan the grooves. then
   dismount the record, disassemble the record player, reassemble it in an astonishingly
   different way, remount the record-and then slowly lower the needle into the outermost
   groove.
Achilles: Golly!
Crab: No sooner had the first strains of sound issued forth than a loud SMASH! filled the
   room. The whole thing fell apart, but particularly badly destroyed was the assembler-
   disassembler. In that painful instant I finally realized, to my chagrin, that the Tortoise
   would ALWAYS be able to focus down upon-if you'll pardon the phrase-the Achilles'
   heel of the system.
Achilles: Upon my soul! You must have felt devastated.
Crab: Yes, I felt rather forlorn for a while. But, happily, that was not the end of the story.
   There is a sequel to the tale, which taught me a valuable lesson, which I may pass on
   to you. On the Tortoise's recommendation, I was browsing through a curious book
   filled with strange Dialogues about many subjects, including molecular biology,
   fugues, Zen Buddhism, and heaven knows what else.
Achilles: Probably some crackpot wrote it. What is the book called:'
Crab: If I recall correctly, it was called Copper, Silver, Gold: an Indestructible Metallic
   Alloy.
Achilles: Oh, Mr. Tortoise told me about it, too. It's by a friend of his, who, it appears, is
   quite taken with metal-logic.
Crab- I wonder which friend it is ... Anyway\_ in one of the Dialogues, I encountered
   some Edifying Thoughts on the Tobacco Mosaic Virus, ribosomes, and other strange
   things I had never heard of.




Edifying Thoughts of a Tobacco Smoker                                                     479

                       FIGURE 79. Tobacco Mosaic Virus.
        From A. Lehninger, Biochemistry (New York: Worth Publishers, 1976).


Achilles: What is the Tobacco Mosaic Virus? What are ribosomes?
Crab: I can't quite say, for I'm a total dunce when it comes to biology. All I know is what
  I gathered from that Dialogue. There, it said that Tobacco Mosaic Viruses are tiny
  cigarette-like objects that cause a disease in tobacco plants.
Achilles: Cancer?
Crab: No, not exactly, but
Achilles: What next? A tobacco plant smoking and getting cancer! Serves it right!
Crab: I believe you've jumped to a hasty conclusion, Achilles. Tobacco plants don't
  SMOKE these "cigarettes". The nasty little "cigarettes" just come and attack them,
  uninvited.
Achilles: I see. Well, now that I know all about Tobacco Mosaic Viruses, tell me what a
  ribosome is.
Crab: Ribosomes are apparently some sort of sub cellular entities which take a message
  in one form and convert it into a message in another form.
Achilles: Something like a teeny tape recorder or phonograph?
Crab: Metaphorically, I suppose so. Now the thing which caught my eye was a line where
  this one exceedingly droll character mentions the fact that ribosomes-as well as
  Tobacco Mosaic Viruses and certain other bizarre biological structures-possess "the
  baffling ability to spontaneously self-assemble, Those were his exact words.
Achilles: That was one of his droller lines, I take it.




Edifying Thoughts of a Tobacco Smoker                                                  480

Crab: That's just what the other character in the Dialogue thought. But that's a
  preposterous interpretation of the statement. (The Crab draws deeply from his pipe,
  and puffs several billows of smoke into the air.)
Achilles: Well, what does "spontaneous self-assembly" mean, then?
Crab: The idea is that when some biological units inside a cell are taken apart, they can
  spontaneously reassemble themselves-without being directed by any other unit. The
  pieces just come together, and presto!-they stick.
Achilles: That sounds like magic. Wouldn't it be wonderful if a full-sized record player
  could have that property? I mean, if a miniature "record player" such as a ribosome
  can do it, why not a big one? That would allow you to create an indestructible
  phonograph, right? Any time it was broken, it would just put itself together again.
Crab: Exactly my thought. I breathlessly rushed a letter off to my manufacturer
  explaining the concept of self-assembly, and asked him if he could build me a record
  player which could take itself apart and spontaneously self-assemble in another form.
Achilles: A hefty bill to fill.
Crab: True; but after several months, he wrote to me that he had succeeded, at long last-
  and indeed he sent me quite a hefty bill. One fine day, ho! My Grand Self-assembling
  Record Player arrived in the mail, and it was with great confidence that I telephoned
  Mr. Tortoise, and invited him over for the purpose of testing my ultimate record
  player.
Achilles: So this magnificent object before us must be the very machine of which you
  speak.
Crab: I'm afraid not, Achilles.
Achilles: Don't tell me that once again ...
Crab: What you suspect, my dear friend is unfortunately the case. I don't pretend to
  understand the reasons why. The whole thing is too painful to recount. To see all those
  springs and wires chaotically strewn about on the floor, and puffs of smoke here and
  there-oh, me ...
Achilles: There, there, Mr. Crab, don't take it too badly.
Crab: I'm quite all right; I just have these spells every so often. Well, to go on, after Mr.
  Tortoise's initial gloating, he at last realized how sorrowful I was feeling, and took
  pity. He tried to comfort me by explaining that it couldn't be helped-it all had to do
  with somebody-or-other's "Theorem", but I couldn't follow a word of it. It sounded
  like "Turtle's Theorem".
Achilles: I wonder if it was that "Gödel’s Theorem" which he spoke of once before to me
  ... It has a rather sinister ring to it. Crab: It could be. I don't recall.
Achilles: I can assure you, Mr. Crab, that I have followed this tale with the utmost
  empathy for your position. It is truly sad. But, you mentioned that there was a silver
  lining. Pray tell, what was that?
Crab: Oh, yes—the silver lining. Well eventually, I abandoned my quest after
  “Perfection” in phonographs, and decided that I might do better




Edifying Thoughts of a Tobacco Smoker                                                    481

  to tighten up my defenses against the Tortoise's records. I concluded that a more
  modest aim than a record player which can play anything is simply a record player that
  can SURVIVE: one that will avoid getting destroyed-even if that means that it can
  only play a few particular records.
Achilles: So you decided you would develop sophisticated anti-Tortoise mechanisms at
  the sacrifice of being able to reproduce every possible sound, eh?
Crab: Well ... I wouldn't exactly say I "decided" it. More accurate would be to say that I
  was FORCED into that position.
Achilles: Yes, I can see what you mean.
Crab: My new idea was to prevent all "alien" records from being played on my
  phonograph. I knew my own records are harmless, and so if I prevented anyone else
  from infiltrating THEIR records, that would protect my record player, and still allow
  me to enjoy my recorded music.
Achilles: An excellent strategy for your new goal. Now does this giant thing before us
  represent your accomplishments to date along those lines?
Crab: That it does. Mr. Tortoise, of course, has realized that he must change HIS strategy,
  as well. His main goal is now to devise a record which can slip past my censors-a new
  type of challenge.
Achilles: For your part, how are you planning to keep his and other "alien" records out?
Crab: You promise you won't reveal my strategy to Mr. T, now?
Achilles: Tortoise's honor.
Crab: What!?
Achilles: Oh-it's just a phrase I've picked up from Mr. T. Don't worry-I swear your secret
  will remain secret with me.
Crab: All right, then. My basic plan is to use a LABELING technique. To each and every
  one of my records will be attached a secret label. Now the phonograph before you
  contains, as did its predecessors, a television camera for scanning the records, and a
  computer for processing the data obtained in the scan and controlling subsequent
  operations. My idea is simply to chomp all records which do not bear the proper label!
Achilles: Ah, sweet revenge! But it seems to me that your plan will be easy to foil. All
  Mr. T needs to do is to get a hold of one of your records, and copy its label!
Crab: Not so simple, Achilles. What makes you think he will be able to tell the label from
  the rest of the record? It may be better integrated than you suspect.
Achilles: Do you mean that it could be mixed up somehow with the actual music?
Crab: Precisely. But there is a way to disentangle the two. It requires sucking the data off
  the record visually and then--




Edifying Thoughts of a Tobacco Smoker                                                   482

Achilles: Is that what that bright green Hash was for?
Crab: That's right. That was the TV camera scanning the grooves. The groove-patterns
  were sent to the minicomputer, which analyzed the musical style of the piece I had put
  on-all in silence. Nothing had been played yet.
Achilles: Then is there a screening process, which eliminates pieces which aren't in the
  proper styles?
Crab: You've got it, Achilles. The only records which can pass this second test are
  records of pieces in my own style-and it will be hopelessly difficult for Mr. T to
  imitate that. So you see, I am convinced I will win this new musical battle. However, I
  should mention that Mr. T is equally convinced that somehow, he will manage to slip
  a record past my censors.
Achilles: And smash your marvelous machine to smithereens?
Crab: Oh, no-he has proved his point on that. Now he just wants to prove to me that he
  can slip a record-an innocuous one-by me, no matter what measures I take to prevent
  it. He keeps on muttering things about songs with strange titles, such as "I Can Be
  Played on Record Player X". But he can't scare MtE! The only thing that worries me a
  little is that, as before, he seems to have some murky arguments which ... which ... (He
  trails off into silence. Then, looking quite pensive, he takes a few puffs on his pipe.)
Achilles: Hmm ... I'd say Mr. Tortoise has an impossible task on his hands. He's met his
  match, at long last!
Crab: Curious that you should think so ... I don't suppose that you know Henkin's
  Theorem forwards and backwards, do you?
Achilles: Know WHOSE Theorem forwards and backwards? I've never heard of anything
  that sounds like that. I'm sure it's fascinating, but I'd rather hear more about "music to
  infiltrate phonographs by". It's an amusing little story. Actually, I guess I can fill in the
  end. Obviously, Mr. T will find out that there is no point in going on, and so he will
  sheepishly admit defeat, and that will be that. Isn't that exactly it?
Crab: That's what I'm hoping, at least. Would you like to see a little bit of the inner
  workings of my defensive phonograph?
Achilles: Gladly. I've always wanted to see a working television camera.
Crab: No sooner said than done, my friend. (Reaches into the gaping"mouth" of the large
  phonograph, undoes a couple of snaps, and pulls out a neatly packaged instrument.)
  You see, the whole thing is built of independent modules, which can be detached and
  used independently. This TV camera, for instance, works very well by itself. Watch
  the screen over there, beneath the painting with the flaming tuba. (He points the
  camera at Achilles, whose face instantly appears on the large screen.)
Achilles: Terrific! May I try it out?
Crab: Certainly.
   Achilles: (pointing the camera at the Crab. There YOU are, Mt Crab, on the screen.




Edifying Thoughts of a Tobacco Smoker                                                     483

                FIGURE 80. The Fair Captive, by Rene Magritte (1947).

Crab: So I am.
Achilles: Suppose I point the camera at the painting with the burning tuba. Now it is on
  the screen, too!
Crab: The camera can zoom in and out, Achilles. You ought to try it. Achilles: Fabulous!
  Let me just focus down onto the tip of those flames, where they meet the picture frame
  ... It's such a funny feeling to be able to instantaneously "copy" anything in the room-
  anything I want-onto that screen. I merely need to point the camera at it, and it pops
  like magic onto the screen.
Crab: ANYTHING in the room, Achilles? Achilles: Anything in sight, yes. That's
  obvious.
Crab: What happens, then, if you point the camera at the flames on the TV screen?

  (Achilles shifts the camera so that it points directly at that part of the television
  screen on which the flames are-or were-displayed.)

Achilles: Hey, that's funny! That very act makes the flames DISAPPEAR from the
  screen! Where did they go?
Crab: You can't keep an image still on the screen and move the camera at the same time.
Achilles: So I see… But I don’t understand what’s on the screen now—not at all! It
  seems to be a strange long corridor. Yet I’m certainly not




Edifying Thoughts of a Tobacco Smoker                                                 484

FIGURE 81. Twelve self-engulfing TV screens. I would have included one more, had 13
  not been prime



Edifying Thoughts of a Tobacco Smoker                                          485

Edifying Thoughts of a Tobacco Smoker   486

  pointing the camera down any corridor. I'm merely pointing it at an ordinary TV
  screen.
Crab: Look more carefully, Achilles. Do you really see a corridor?
Achilles: Ahhh, now I see. It's a set of nested copies of the TV screen itself, getting
  smaller and smaller and smaller ... Of course! The image of the flames HAD to go
  away, because it came from my- pointing the camera at the PAINTING. When I point
  the camera at the SCREEN, then the screen itself appears, with whatever is on the
  screen at the time which is the screen itself, with whatever is on the screen at the time
  which is the screen itself, with
Crab: I believe I can fill in the rest, Achilles. Why- don't you try rotating the camera?
Achilles: Oh! I get a beautiful spiraling corridor! Each screen is rotated inside its framing
  screen, so that the littler they get, the more rotated they are, with respect. to the
  outermost screen. This idea of having a TV screen "engulf itself" is weird.
Crab: What do you mean by "self-engulfing", Achilles?
Achilles: I mean, when I point the camera at the screen-or at part of the screen. THAT'S
  self-engulfing.
Crab: Do you mind if I pursue that a little further? I'm intrigued by this new notion.
Achilles: So am I.
Crab: Very well, then. If you point the camera at a CORNER of the screen, is that still
  what you mean by "self-engulfing"?
Achilles: Let me try it. Hmm-the "corridor" of screens seems to go off the edge, so there
  isn't an infinite nesting any more. It's pretty, but it doesn't seem to me to have the spirit
  of self-engulfing. It's a "failed self-engulfing".
Crab: If you were to swing the TV camera back towards the center of the screen, maybe
  you could fix it up again ...
Achilles (slowly and cautiously turning the camera): Yes! The corridor is getting longer
  and longer ... There it is! Now it's all back. I can look down it so far that it vanishes in
  the distance. The corridor became infinite again precisely at the moment when the
  camera took in the WHOLE screen. Hmm-that reminds me of something Mr. Tortoise
  was saying a while back, about self-reference only occurring when a sentence talks
  about ALL of itself ...
Crab: Pardon me?
Achilles: Oh, nothing just muttering to myself.

   (As Achilles plays with the lens and other controls on the camera, a profusion of new
   kinds of self-engulfing images appear: swirling spirals that resemble galaxies,
   kaleidoscopic flower-like shapes, and other assorted patterns ...)

Crab: You seem to be having a grand time.
Achilles: (turns away from the camera); I’ll say! What a wealth of images this simple
  idea can produce! (He glances back at the screen, and a look of




Edifying Thoughts of a Tobacco Smoker                                                      487

  astonishment crosses his face.) Good grief, Mr. Crab! There's a pulsating petal-pattern
  on the screen! Where do the pulsations come from? The TV is still, and so is the
  camera.
Crab: You can occasionally set up patterns which change in time. This is because there is
  a slight delay in the circuitry between the moment the camera "sees" something, and
  the moment it appears on the screen around a hundredth of a second. So if you have a
  nesting of depth fifty or so, roughly a half-second delay will result. If somehow a
  moving image gets onto the screen-for example, by you putting your finger in front of
  the camera-then it takes a while for the more deeply nested screens to "find out" about
  it. This delay then reverberates through the whole system, like a visual echo. And if
  things are set up so the echo doesn't die away, then you can get pulsating patterns.
Achilles: Amazing! Say-what if we tried to make a TOTAL self-engulfing?
Crab: What precisely do you mean by that?
Achilles: Well, it seems to me that this stuff with screens within screens is interesting, but
  I'd like to get a picture of the TV camera AND the screen, ON the screen. Only then
  would I really have made the system engulf itself. For the screen is only PART of the
  total system.
Crab: I see what you mean. Perhaps with this mirror, you can achieve the effect you
  want.

   (The Crab hands him a mirror, and Achilles maneuvers the mirror and camera in such
   a way that the camera and the screen are both pictured on the screen.)

Achilles: There! I've created a TOTAL self-engulfing!
Crab: It seems to me you only have the front of the mirror-what about its back? If it
  weren't for the back of the mirror, it wouldn't be reflective-and you wouldn't have the
  camera in the picture.
Achilles: You're right. But to show both the front and back of this mirror, I need a second
  mirror.
Crab: But then you'll need to show the back of that mirror, too. And what about including
  the back of the television, as well as its front? And then there's the electric cord, and
  the inside of the television, and
Achilles: Whoa, whoa! My head's beginning to spin! I can see that this "total self-
  engulfing project" is going to pose a wee bit of a problem. I'm feeling a little dizzy.
Crab: I know exactly how you feel. Why don't you sit down here and take your mind off
  all this self-engulfing? Relax! Look at my paintings, and you'll calm down.
(Achilles lies down, and sighs.)
Oh-perhaps my pipe smoke is bothering you? Here, I'll put my pipe away. (Takes the pipe
  from his mouth, and carefully places it above some written words in another .Magritte
  painting.) There! Feeling any better?
Achilles: I’m still a little woozy, (Points at the Magritte.) That’s an interesting painting. I
  like the way it’s framed, especially the shiny inlay inside the wooden frame.




Edifying Thoughts of a Tobacco Smoker                                                      488

              FIGURE 82. The Air and the Song, by Rene Magritte (1964).

Crab: Thank you. I had it specially done-it's a gold lining.
Achilles: A gold lining? What next? What are those words below the pipe? They aren't in
  English, are they?
Crab: No, they are in French. They say, "Ceci n'est pas une pipe." That means, "This is
  not a pipe". Which is perfectly true. 4chilles: But it is a pipe! You were just smoking
  it!
Crab: Oh, you misunderstand the phrase, I believe. The word "ceci" refers to the painting,
  not to the pipe. Of course the pipe is a pipe. But a painting is not a pipe.
Achilles: I wonder if that "ceci" inside the painting refers to the WHOLE painting, or just
  to the pipe inside the painting. Oh, my gracious! That would be ANOTHER self-
  engulfing! I'm not feeling at all well, Mr. Crab. I think I'm going to be sick ...




Edifying Thoughts of a Tobacco Smoker                                                  489


                                  Self-Ref and Self-Rep
IN THIS CHAPTER, we will look at some of the mechanisms which create self-reference in
various contexts, and compare them to the mechanisms which allow some kinds of systems
to reproduce themselves. Some remarkable and beautiful parallels between these mechanisms
will come to light.

                    Implicitly and Explicitly Self-Referential Sentences

To begin with, let us look at sentences which, at first glance, may seem to provide the
simplest examples of self-reference. Some such sentences are these:

       (1) This sentence contains five words.
       (2) This sentence is meaningless because it is self-referential.
       (3) This sentence no verb.
       (4) This sentence is false. (Epimenides paradox)
       (5) The sentence I am now writing is the sentence you are now reading.

All but the last one (which is an anomaly) involve the simple-seeming mechanism contained
in the phrase "this sentence". But that mechanism is in reality far from simple. All of these
sentences are "floating" in the context of the English language. They can be compared to
icebergs, whose tips only are visible. The word sequences are the tips of the icebergs, and the
processing which must be done to understand them is the hidden part. In this sense their
meaning is implicit, not explicit. Of course, no sentence's meaning is completely explicit, but
the more explicit the self-reference is, the more exposed will be the mechanisms underlying
it. In this case, for the self-reference of the sentences above to be recognized, not only has
one to be comfortable with a language such as English which can deal with linguistic subject
matter, but also one has to be able to figure out the referent of the phrase "this sentence". It
seems simple, but it depends on our very complex yet totally assimilated ability to handle
English. What is especially important here is the ability to figure out the referent of a noun
phrase with a demonstrative adjective in it. This ability is built up slowly, and should by no
means be considered trivial. The difficulty is perhaps underlined when a sentence such as
number 4 is presented to someone naive about paradoxes and linguistic tricks, such as a
child. They may say, "What sentence is false and it may take a bit of persistence to get across
the idea that the sentence is talking about itself. The whole idea is a little mind




Self-Rep and Self-Rep                                                                       490

boggling at first. A couple of pictures may help (Figs. 83, 84). Figure 83 is a picture which
can be interpreted on two levels. On one level, it is a sentence pointing at itself; on the other
level, it is a picture of Epimenides executing his own death sentence.




.

  Figure 84, showing visible and invisible portions of the iceberg, suggests the relative
proportion of sentence to processing required for the recognition of self-reference:




        It is amusing to try to create a self-referring sentence without using the trick of saving
this sentence". One could try to quote a sentence inside itself. Here is an attempt:

The sentence "The sentence contains five words" contains five words.

But such an attempt must fail, for any sentence that could be quoted entirely inside itself
would have to be shorter than itself. This is actually possible, but only if you are willing to
entertain infinitely long sentences, such as:



Self-Rep and Self-Rep                                                                         491

The sentence
       "The sentence
              "The sentence
                     "The sentence
                            …
                                            ….
                                                                etc,.,etc.
                                            …
                                     ….
                      is infinitely long"
             is infinitely long"
         is infinitely long'
is infinitely long.

But this cannot work for finite sentences. For the same reason, Godel's string G could not
contain the explicit numeral for its Godel number: it would not fit. No string of TNT can
contain the TNT-numeral for its own Godel number, for that numeral always contains more
symbols than the string itself does. But you can get around this by having G contain a
description of its own Godel number, by means of the notions of "sub" and
"arithmoquinification".
         One way of achieving self-reference in an English sentence by means of description
instead of by self-quoting or using the phrase "this sentence" is the Quine method, illustrated
in the dialogue Air on G's String. The understanding of the Quine sentence requires less
subtle mental processing than the four examples cited earlier. Although it may appear at first
to be trickier, it is in some ways more explicit. The Quine construction is quite like the Godel
construction, in the way that it creates self-reference by describing another typographical
entity which, as it turns out, is isomorphic to the Quine sentence itself. The description of the
new typographical entity is carried out by two parts of the Quine sentence. One part is a set
of instructions telling how to build a certain phrase, while the other part contains the
construction materials to be used; that is, the other part is a template. This resembles a
floating cake of soap more than it resembles an iceberg (See Fig. 85).




Self-Rep and Self-Rep                                                                        492

The self-reference of this sentence is achieved in a more direct way than in the Epimenides
paradox; less hidden processing is needed. By the way, it is interesting to point out that the
phrase "this sentence" appears in the previous sentence; yet it is not there to cause self-
reference: you probably understood that its referent was the Quine sentence, rather than the
sentence in which it occurs. This just goes to show how pointer phrases such as "this
sentence" are interpreted according to context, and helps to show that the processing of such
phrases is indeed quite involved.

                                A Self-Reproducing Program

The notion of quining, and its usage in creating self-reference, have already been explained
inside the Dialogue itself, so we need not dwell on such matters here. Let us instead show
how a computer program can use precisely the same technique to reproduce itself. The
following selfreproducing program is written in a BlooP-like language and is based on
following a phrase by its own quotation (the opposite order from quining, so I reverse the
name "quine" to make "eniuq"):

   DEFINE PROCEDURE "ENIUQ" [TEMPLATE]: PRINT [TEMPLATE, LEFT-
   BRACKET, QUOTE-MARK, TEMPLATE, QUOTE-MARK, RIGHT-BRACKET,
   PERIOD].

   ENIUQ
     ['DEFINE PROCEDURE "ENIUQ" [TEMPLATE]: PRINT [TEMPLATE, LEFT-
     BRACKET, QUOTE-MARK, TEMPLATE, QUOTE-MARK, RIGHT-BRACKET,
     PERIOD]. ENIUQ'].

ENIUQ is a procedure defined in the first two lines, and its input is called "TEMPLATE". It
is understood that when the procedure is called, TEMPLATE's value will be some string of
typographical characters. The effect of ENIUQ is to carry out a printing operation, in which
TEMPLATE gets printed twice: the first time just plain; the second time wrapped in (single)
quotes and brackets, and garnished with a final period. Thus, if TEMPLATE's value were
the string DOUBLE-BUBBLE, then performing ENIUQ on it would yield:

DOUBLE-BUBBLE ['DOUBLE-BUBBLE'].

Now in the last four lines of the program above, the procedure ENIUQ is called with a
specific value of TEMPLATE-namely the long string inside the single quotes: DEFINE ...
ENIUQ. That value has been carefully chosen; it consists of the definition of ENIUQ,
followed by the word ENIUQ. This makes the program itself-or, if you prefer, a perfect copy
of it-get printed out. It is very similar to Quine's version of the Epimenides sentence:

                      "yields falsehood when preceded by its quotation"
                       yields falsehood when preceded by its quotation.

It is very important to realize that the character string which appears n quotes in the last three
lines of the program above-that is, the value of




Self-Rep and Self-Rep                                                                         493

TEMPLATE-is never interpreted as a sequence of instructions. That it happens to be one is,
in a sense. just an accident. As was pointed out above, it could just as well have been
DOUBLE-BUBBLE or any other string of characters. The beauty of the scheme is that when
the same string appears in the top two lines of this program, it is treated as a program
(because it is not in quotes). Thus in this program, one string functions in two ways: first as
program, and second as data. This is the secret of self-reproducing programs, and, as we shall
see, of self-reproducing molecules. It is useful, incidentally, to call any kind of self-
•reproducing object or entity a self-rep; and likewise to call any self-referring object or entity
a self-ref. I will use those terms occasionally from here on.
         The preceding program is an elegant example of a self-reproducing program written
in a language which was not designed to make the writing of self-reps particularly easy.
Thus, the task had to be carried out using those notions and operations which were assumed
to be part of the language-such as the word QUOTE-MARK, and the command PRINT. But
suppose a language were designed expressly for making self-reps easy to write. Then one
could write much shorter self-reps. For example, suppose that the operation of eniuq-ing
were a built-in feature of the language, needing no explicit definition (as we assumed PRINT
was). Then a teeny self-rep would be this:

                                      ENIUQ ['ENIUQ'].

It is very similar to the Tortoise's version of Quine's version of the Epimenides self-ref,
where the verb "to quine" is assumed to be known:

                "yields falsehood when quined" yields falsehood when quined

        But self-reps can be even shorter. For instance, in some computer language it might
be a convention that any program whose first symbol is an asterisk is to be copied before
being executed normally. Then the program consisting of merely one asterisk is a self-rep!
You may complain that this is silly and depends on a totally arbitrary convention. In doing
so, you are echoing my earlier point that it is almost cheating to use the phrase "this
sentence" to achieve self-reference-it relies too much on the processor, and not enough on
explicit directions for self-reference. Using an asterisk as an example of a self-rep is like
using the word "I" as an example of a self-ref: both conceal all the interesting aspects of their
respective problems.
        This is reminiscent of another curious type of self-reproduction: via photocopy
machine. It might be claimed that any written document is a self-rep because it can cause a
copy of itself to be printed when it is placed in a photocopy machine and the appropriate
button is pushed. But somehow this violates our notion of self-reproduction; the piece of
paper is not consulted at all, and is therefore not directing its own reproduction. Again,
everything is in the processor. Before we call something a self-rep, we want to have the
feeling that, to the maximum extent possible, it explicitly contains the directions for copying
itself.




Self-Rep and Self-Rep                                                                         494

To be sure, explicitness is a matter of degree; nonetheless there is an intuitive borderline on
one side of which we perceive true self-directed self-reproduction, and on the other side of
which we merely see copying being carried out by an inflexible and autonomous copying
machine.

                                      What Is a Copy?

Now in any discussion of self-refs and self-reps, one must sooner or later come to grips with
the essential issue: what is a copy? We already dealt with that question quite seriously in
Chapters V and VI; and now we come back to it. To give the flavor of the issue, let us
describe some highly fanciful, yet plausible, examples of self-reps.

                                  A Self-Reproducing Song

Imagine that there is a nickelodeon in the local bar which, if you press buttons 11-U, will
play a song whose lyrics go this way:

                          Put another nickel in, in the nickelodeon,
                         All I want is 11-U, and music, music, music.

We could make a little diagram of what happens one evening (Fig. 86).




FIGURE 86. A self-reproducing song.

Although the effect is that the song reproduces itself, it would feel strange to call the song a
self-rep, because of the fact that when it passes through the 11-U stage, not all of the
information is there. The information only gets put back by virtue of the fact that it is fully
stored in the nickelodeon that is, in one of the arrows in the diagram, not in one of the ovals.
It is questionable whether this song contains a complete description of how to get itself
played again, because the symbol pair "1 1-U" is only a trigger, not a copy.

                                     A "Crab" Program

Consider next a computer program which prints itself out backwards. (Some readers might
enjoy thinking about how to write such a program in




Self-Rep and Self-Rep                                                                       495

the blooP-like language above, using the given sell-rep as a inouel.) vvouiu this funny
program count as a self-rep. Yes, in a way, because a trivial transformation performed on its
output will restore the original program. It seems fair to say that the output contains the same
information as the program itself, just recast in a simple way. Yet it is clear that someone
might look at the output and not recognize it as a program printed backwards. To recall
terminology from Chapter VI, we could say that the "inner messages" of the output and the
program itself are the same, but they have different "outer messages"-that -is, they must be
read by using different decoding mechanisms. Now if one counts the outer message as part of
the information-which seems quite reasonable-then the total information is not the same after
all, so the program can't be counted as a self-rep.
         However, this is a disquieting conclusion, because we are accustomed to considering
something and its mirror image as containing the same information. But recall that in Chapter
VI, we made the concept of "intrinsic meaning" dependent on a hypothesized universal
notion of intelligence. The idea was that, in determining the intrinsic meaning of an object,
we could disregard some types of outer message-those which would be universally
understood. That is, if the decoding mechanism seems fundamental enough, in some still ill-
defined sense, then the inner message which it lets be revealed is the only meaning that
counts. In this example, it seems reasonably safe to guess that a "standard intelligence" would
consider two mirror images to contain the same information as each other; that is, it would
consider the isomorphism between the two to be so trivial as to be ignorable. And thus our
intuition that the program is in some sense a fair self-rep, is allowed to stand.

                               Epimenides Straddles the Channel

Now another far-fetched example of a self-rep would be a program which prints itself our,
but translated into a different computer language. One might liken this to the following
curious version of the Quine version of the Epimenides self-ref:

     'lest une expression qui, quand elle est precedee de sa traduction, mise entre
     guillemets, clans la langue provenant de l'autre tote de la Manche. tree une faussete"
     is an expression which, when it is preceded by its translation, placed in quotation
     marks, into the language originating on the other side of the Channel, yields a
     falsehood.

You might try to write down the sentence which is described by this weird concoction. (Hint:
It is not itself-or at least it is not if "itself" is taken in a naive sense.) If the notion of "self-rep
by retrograde motion" (i.e., a program which writes itself out backwards) is reminiscent of a
crab canon, the notion of "self-rep by translation" is no less reminiscent of "a canon which
involves a transposition of the theme into another key.




Self-Rep and Self-Rep                                                                               496

                    A Program That Prints Out Its Own Godel Number

The idea of printing out a translation instead of an exact copy of the original program may
seem pointless. However, if you wanted to write a self-rep program in BlooP or FlooP, you
would have to resort to some such device, for in those languages, OUTPUT is always a
number, rather than a typographical string. Therefore, you would have to make the program
print out its own Godel number: a very huge integer whose decimal expansion codes for the
program, character by character, by using three digit codons. The program is coming as close
as it can to printing itself, within the means available to it: it prints out a copy of itself in
another "space", and it is easy to switch back and forth between the space of integers and the
space of strings. Thus, the value of OUTPUT is not a mere trigger, like "11-12". Instead, all
the information of the original program lies "close to the surface" of the output.

                                   Godelian Self-Reference

This comes very close to describing the mechanism of Godel's self-ref G. After all, that string
of TNT contains a description not of itself, but of an integer (the arithmoquinification of u). It
just so happens that that integer is an exact "image" of the string G, in the space of natural
numbers. Thus, G refers to a translation of itself into another space. We still feel comfortable
in calling G a self-referential string, because the isomorphism between the two spaces is so
tight that we can consider them to be identical.
         This isomorphism that mirrors TNT inside the abstract realm of natural numbers can
be likened to the quasi-isomorphism that mirrors the real world inside our brains, by means
of symbols. The symbols play quasi-isomorphic roles to the objects, and it is thanks to them
that we can think. Likewise, the Godel numbers play isomorphic roles to strings, and it is
thanks to them that we can find metamathematical meanings in statements about natural
numbers. The amazing, nearly magical, thing about G is that it manages to achieve self-
reference despite the fact that the language in which it is written, TNT, seems to offer no
hope of referring to its own structures, unlike English, in which it is the easiest thing in the
world to discuss the English language.
         So G is an outstanding example of a self-ref via translation-hardly the most
straightforward case. One might also think back to some of the Dialogues, for some of them,
too, are self-refs via translation. For instance, take the Sonata for Unaccompanied Achilles.
In that Dialogue, several references are made to the Bach Sonatas for unaccompanied violin,
and the Tortoise's suggestion of imagining harpsichord accompaniments is particularly
interesting. After all, if one applies this idea to the Dialogue itself, one invents lines which
the Tortoise is saying; but if one assumes that Achilles' part stands alone (as does the violin),
then it is quite wrong to attribute any lines at all to the Tortoise. In any case, here again is a
self-ref by means of a mapping which maps Dialogues onto pieces by Bach. And this
mapping is




Self-Rep and Self-Rep                                                                         497

left, of course, for the reader to notice. Yet even if the reader does not notice it, the mapping
is still there, and the Dialogue is still a self-ref.

                                A Self-Rep by Augmentation

We have been likening self-reps to canons. What, then, would be a fair analogue to a canon
by augmentation? Here is a possibility: consider a program which contains a dummy loop
whose only purpose is to slow up the program. A parameter might tell how often to repeat the
loop. A self-rep could be made which prints out a copy of itself, but with the parameter
changed, so that when that copy is run, it will run at half the speed of its parent program; and
its "daughter" will in turn run at half again the speed, and so on . . . None of these programs
prints itself out precisely; yet all clearly belong to a single "family".
         This is reminiscent of the self-reproduction of living organisms. Clearly, an
individual is never identical to either of its parents; why, then, is the act of making young
called "self-reproduction'? The answer is that there is a coarse-grained isomorphism between
parent and child; it is an isomorphism which preserves the information about species. Thus,
what is reproduced is the class, rather than the instance. This is also the case in the recursive
picture Gplot, in Chapter V: that is, the mapping between "magnetic butterflies" of various
sizes and shapes is coarse-grained; no two are identical, but they all belong to a single
"species", and the mapping preserves precisely that fact. In terms of self-replicating
programs, this would correspond to a family of programs, all written in "dialects" of a single
computer language; each one can write itself out, but slightly modified, so that it comes out
in a dialect of its original language.

                                     A Kimian Self-Rep

 Perhaps the sneakiest example of a self-rep is the following: instead of writing a legal
expression in the compiler language, you type one of the compiler's own error messages.
When the compiler looks at your "program", the first thing it does is get confused, because
your "program" is ungrammatical; hence the compiler prints out an error message. All you
need to do is arrange that the one it prints out will be the one you typed in. This kind of self-
rep, suggested to me by Scott Kim, exploits a different level of the system from the one you
would normally approach. Although it may seem frivolous, it may have counterparts in
complex systems where self-reps vie against each other for survival, as we shall soon discuss.

                                    What Is the Original?

Besides the question "What constitutes a copy?", there is another fundamental philosophical
question concerning self-reps. That is the obverse




Self-Rep and Self-Rep                                                                        498

side of the coin: "What is the original?" This can best be explained by referring to some
examples:

      (1) a program which, when interpreted by some interpreter running on some
     computer, prints itself out;

     (2) a program which, when interpreted by some interpreter running on some
     computer. prints itself out along with a complete copy of the interpreter (which,
     after all, is also a program);

     (3) a program which, when interpreted by some interpreter running on some
     computer, not only prints itself out along with a complete copy of the interpreter,
     but also directs a mechanical assembly process in which a second computer,
     identical to the one on which the interpreter and program are running, is put
     together.

It is clear that in (1), the program is the self-rep. But in (3), is it the program which is the self-
rep, or the compound system of program plus interpreter, or the union of program,
interpreter, and processor?
         Clearly, a self-rep can involve more than just printing itself out. In fact, most of the
rest of this Chapter is a discussion of self-reps in which data, program, interpreter, and
processor are all extremely intertwined, and in which self-replication involves replicating all
of them at once.

                                           Typogenetics

We are now about to broach one of the most fascinating and profound topics of the twentieth
century: the study of "the molecular logic of the living state", to borrow Albert Lehninger's
richly evocative phrase. And logic it is, too but of 'a sort more complex and beautiful than
any a human mind ever imagined. We will come at it from a slightly novel angle: via an
artificial solitaire game which I call Typogenetics-short for "Typographical Genetics". In
Typogenetics I have tried to capture some ideas of molecular genetics in a typographical
system which, on first sight, resembles very much the formal systems exemplified by the
MIU-system. Of course, Typogenetics involves many simplifications, and therefore is useful
primarily for didactic purposes.
         I should explain immediately that the field of molecular biology is a field in which
phenomena on several levels interact, and that Typogenetics is only trying to illustrate
phenomena from one or two levels. In particular, purely chemical aspects have been
completely avoided-they belong to a level lower than is here dealt with; similarly, all aspects
of classical genetics (viz., nonmolecular genetics) have also been avoided-they belong to a
level higher than is here dealt with. I have intended in Typogenetics only to give an intuition
for those processes centered on the celebrated Central Dogma of




Self-Rep and Self-Rep                                                                            499

Molecular Biology, enunciated by Francis Crick (one of the co-discoverers of the double-
helix structure of DNA):

                                DNA => RNA => proteins.

It is my hope that with this very skeletal model I have constructed the reader will perceive
some simple unifying principles of the field principles which might otherwise be obscured by
the enormously intricate interplay of phenomena at many different levels. What is sacrificed
is, of course, strict accuracy; what is gained is, I hope, a little insight.

                                  Strands, Bases, Enzymes

The game of Typogenetics involves typographical manipulation on sequences of letters.
There are four letters involved:

                                          A C G T.

Arbitrary sequences of them are called strands. Thus, some strands are:

                                        GGGG
                                      ATTACCA
                                    CATCATCATCAT

Incidentally, "STRAND" spelled backwards begins with "DNA". This is appropriate since
strands, in Typogenetics, play the role of pieces of DNA (which, in real genetics, are often
called "strands"). Not only this, but "STRAND" fully spelled out backwards is "DNA RTS",
which may be taken as an acronym for "DNA Rapid Transit Service". This, too, is
appropriate, for the function of "messenger RNA"-which in Typogenetics is represented by
strands as well-is quite well characterized by the phrase "Rapid Transit Service" for DNA, as
we shall see later.
        I will sometimes refer to the letters A, C, G, T as bases, and to the positions which
they occupy as units. Thus, in the middle strand, there are seven units, in the fourth of which
is found the base A.
        If you have a strand, you can operate on it and change it in various ways. You can
also produce additional strands, either by copying, or by cutting a strand in two. Some
operations lengthen strands, some shorten them, and some leave their length alone.
        Operations come in packets-that is, several to be performed together, in order. Such a
packet of operations is a little like a programmed machine which moves up and down the
strand doing things to it. These mobile machines are called "typographical enzymes"-
enzymes for short. Enzymes operate on strands one unit at a time, and are said to be "bound"
to the unit they are operating on at any given moment.
        I will show how some sample enzymes act on particular strings. The first thing to
know is that each enzyme likes to start out bound to a particular letter. Thus, there are four
kinds of enzyme-those which prefer




Self-Rep and Self-Rep                                                                      500

A, those which prefer C, etc. Given the sequence of operations which an enzyme performs,
you can figure out which letter it prefers, but for now I'll just give them without explanation.
Here's a sample enzyme, consisting of three operations:

   (1) Delete the unit to which the enzyme is bound (and then bind to the next unit to the
   right).
   (2) Move one unit to the right.
   (3) Insert a T (to the immediate right of this unit).

This enzyme happens to like to bind to A initially. And here's a sample strand:

                                               ACA

What happens if our enzyme binds to the left A and begins acting? Step I deletes the A, so we
are left with CA-and the enzyme is now bound to the C. Step 2 slides the enzyme rightwards,
to the A, and Step 3 appends a T onto the end to form the strand CAT. And the enzyme has
done its complete duty: it has transformed ACA into CAT.
         What if it had bound itself to the right A of ACA? It would have deleted that A and
moved off the end of the strand. Whenever this happens, the enzyme quits (this is a general
principle). So the entire effect would just be to lop off one symbol.
Let's see some more examples. Here is another enzyme:

   (1) Search for the nearest pyrimidine to the right of this unit.
   (2) Go into Copy mode.
   (3) Search for the nearest purine to the right of this unit.
   (4) Cut the strand here (viz., to the right of the present unit).

Now this contains the terms "pyrimidine" and "purine". They are easy terms. A and G are
called purines, and C and T are called pyrimidines. So searching for a pyrimidine merely
means searching for the nearest C or T.

                               Copy Mode and Double Strands

The other new term is Copy mode. Any strand can be "copied" onto another strand, but in a
funny way. Instead of copying A onto A, you copy it onto T, and vice versa. And instead of
copying C onto C, you copy it onto G, and vice versa. Note that a purine copies onto a
pyrimidine, and vice versa. This is called complementary base pairing. The complements are
shown below

               Complement
Purinas        A <====> T pyrimidines
               G <====> C




Self-Rep and Self-Rep                                                                       501

You can perhaps remember this molecular pairing scheme by recalling that Achilles is paired
with the Tortoise, and the Crab with his Genes.
        When "copying" a strand, therefore, you don't actually copy it, but you manufacture
its complementary strand. And this one will be written upside down above the original strand.
Let's see this in concrete terms. Let the previous enzyme act on the following strand (and that
enzyme also happens to like to start at A):

                              CAAAGAGAATCCTCTTTGAT

There are many places it could start. Let's take the second A, for example. The enzyme binds
to it, then executes step 1: Search for the nearest pyrimidine to the right. Well, this means a C
or a T. The first one is a T somewhere near the middle of the strand, so that's where we go.
Now step 2: Copy mode. Well, we just put an upside-down A above our T. But that's not all,
for Copy mode remains in effect until it is shut off-or until the enzyme is done, whichever
comes first. This means that every base which is passed through by the enzyme while Copy
mode is on will get a complementary base put above it. Step 3 says to look for a purine to the
right of our T. That is the G two symbols in from the right-hand end. Now as we move up to
that G, we must "copy"-that is, create a complementary strand. Here's what that gives:
(editor’ s note, I can’t print upside down ie it is too much hard work so V = A and D = G
upside down)




The last step is to cut the strand. This will yield two pieces:

                                        VDDVDVVVJ
                                CAAAGAGAATCCTCTTTG

                                             and AT.

And the instruction packet is done. We are left with a double strand, however. Whenever this
happens, we separate the two complementary strands from each other (general principle); so
in fact our end product is a set of three strands:

                AT, CAAAGAGGA,             and CAAAGAGAATCCTCTTTG

Notice that the upside-down strand has been turned right side up, and thereby right and left
have been reversed.
         Now you have seen most of the typographical operations which can be carried out on
strands. There are two other instructions which should be mentioned. One shuts off Copy
mode; the other switches the enzyme from a strand to the upside-down strand above it. When
this happens, if you keep the paper right side up, then you must switch "left" and "right" in all
the instructions. Or better, you can keep the wording and just turn the paper around so the top
strand becomes legible. If the "switch" command is




Self-Rep and Self-Rep                                                                        502

given, but there is no complementary base where the enzyme is bound at that instant, then the
enzyme just detaches itself from the strand, and its job is done.
        It should be mentioned that when a "cut" instruction is encountered, this pertains to
both strands (if there are two): however, "delete" pertains only to the strand on which the
enzyme is working. If Copy mode is on, then the "insert" command pertains to both strands-
the base itself into the strand the enzyme is working on, and its complement into the other
strand.
If Copy mode is off, then the "insert" command pertains only to the one strand, so a blank
space must he inserted into the complementary strand.
        And, whenever Copy mode is on, "move" and "search" commands require that one
manufacture complementary bases to all bases which the sliding enzyme touches.
Incidentally, Copy mode is always off when an enzyme starts to work. If Copy mode is off,
and the command "Shut off copy mode" is encountered, nothing happens. Likewise, If Copy
mode is already on, and the command "Turn copy mode on" is encountered, then nothing
happens.

                                           Amino Acids

There are fifteen types of command, listed below:

    Cut         cut strand(s)
    del         delete a base from strand
    swi         switch enzyme to other strand
    mvr         move one unit to the right
    mvl         move one unit to the left
    cop         turn on Copy mode
    off         turn off Copy mode
    ina         insert A to the right of this unit
    inc         insert C to the right of this unit
    ing         insert G to the right of this unit
    int         insert T to the right of this unit
    rpy         search for the nearest pyrimidine to the right
    rpu         search for the nearest purine to the right
    Ipy         search for the nearest pyrimidine to the left
    lpu         search for the nearest purine to the left

Each one has a three-letter abbreviation. We shall refer to the three-letter abbreviations of
commands as amino acids. Thus, every enzyme is made up of a sequence of amino acids. Let
us write down an arbitrary enzyme:

       rpu - inc - cop - myr - tnyl - swi - Ipu - int

and an arbitrary strand:

                                  TAGATCCAGTCCATCGA




Self-Rep and Self-Rep                                                                    503

and see how the enzyme acts on the strand. It so happens that the enzyme binds to G only.
Let us bind to the middle G and begin. Search rightwards for a purine (viz., A or G). We (the
enzyme) skip over TCC and land on A.
Insert a C. Now we have

                                 TAGATCCAGTCCACTCGA

where the arrow points to the unit to which the enzyme is bound. Set Copy mode. This puts
an upside-down G above the C. Move right, move left, then switch to the other strand. Here's
what we have so far:

                                         7V
                                 TAGATCCAGTCCACTCGA

Let's turn it upside down, 'so that the enzyme is attached to the lower strand:

                                  VDJIDV»1DVDD1VDV1
                                          AG

Now we search leftwards for a purine and find A. Copy mode is on, but the complementary
bases are already there, so nothing is added. Finally, we insert a T (in Copy mode), and quit:

                                  VD)IVJVJJl7V)DIVDV1
                                         ATG

Our final product is thus two strands:

        ATG, and TAGATCCAGTCCACATCGA

The old one is of course gone.

                           Translation and the Typogenetic Code

Now you might be wondering where the enzymes and strands come from, and how to tell the
initial binding-preference of a given enzyme. One way might be just to throw some random
strands and some random enzymes together, and see what happens when those enzymes act
on those strands and their progeny. This has a similar flavor to the MU-puzzle, where there
were some given rules of inference and an axiom, and you just began. The only difference is
that here, every time a strand is acted on, its original form is gone forever. In the MU-puzzle,
acting on MI to make MIU didn't destroy MI
         But in Typogenetics, as in real genetics, the scheme is quite a bit trickier. We do
begin with some arbitrary strand, somewhat like an axiom in a formal system. But we have,
initially, no "rules of inference"-that is, no enzymes. However, we can translate each strand
into one or more enzymes! Thus, the strands themselves will dictate the operations which
will be performed upon them, and those operations will in turn produce



Self-Rep and Self-Rep                                                                       504

new strands which will dictate further enzymes, etc. etc.! This is mixing levels with a
vengeance! Think, for the sake of comparison, how different the MU-puzzle would have
been if each new theorem produced could have been turned into a new rule of inference by
means of some code.
  How is this "translation" done? It involves a Typogenetic Code by which adjacent pairs of
bases-called "duplets"-in a single strand represent different amino acids. There are sixteen
possible duplets: AA, AC, AG, AT, CA, CC, etc. And there are fifteen amino acids. The
Typogenetic Code is shown in Figure 87.




According to the table, the translation of the duplet GC is "inc" ("insert a C"); that of AT is
"swi" ("switch strands"); and so on. Therefore it becomes clear that a strand can dictate an
enzyme very straightforwardly. For example, the strand

                               TAGATCCAGTCCACATCGA

breaks up into duplets as follows:

                            TA GA TC CA GT CC AC AT CG A

with the A left over at the end. Its translation into an enzyme is:

      rpy - ina - rpu - mvr - int - mvl - cut - swi - cop.

(Note that the leftover A contributes nothing.)

                                Tertiary Structure of Enzymes

What about the little letters 's', 'l', and 'r' in the lower righthand corner of each box\% They are
crucial in determining the enzyme's binding-preference, and in a peculiar way. In order to
figure out what letter an enzyme likes to bind to, you have to figure out the enzyme's "tertiary
structure", which is itself determined by the enzyme's "primary structure". By its




Self-Rep and Self-Rep                                                                          505

primary structure is meant its amino acid sequence. By its tertiary structure is meant the way
it likes to "fold up". The point is that enzymes don't like being in straight lines, as we have so
far exhibited them. At each internal amino acid (all but the two ends), there is a possibility of
a "kink", which is dictated by the letters in the corners. In particular, '1' and 'r' stand for "left"
and "right", and 's' stands for "straight". So let us take our most recent sample enzyme, and let
it fold itself up to show its tertiary structure. We will start with the enzyme's primary
structure, and move along it from left to right. At each amino acid whose corner-letter is '1'
we'll put a left turn, for those with 'r', we'll put a right turn, and at 's' we'll put no turn. In
Figure 88 is shown the two-dimensional conformation for our enzyme.

Cop

        swi <== cut <== mvl <== int
                               mvr
              rpy -==> ima ==> rpu

FIGURE 88. The tertiary structure of a typoenzyme.

Note the left-kink at "rpu", the right-kink at "swi", and so on. Notice also that the first
segment ("rpy z> ina") and the last segment ("swi => cop") are perpendicular. This is the key
to the binding-preference. In fact, the relative orientation of the first and last segments of an
enzyme's tertiary structure determines the binding-preference of the enzyme. We can always
orient the enzyme so that its first segment points to the right. If we do so, then the last
segment determines the binding-preference, as shown in Figure 89.




Self-Rep and Self-Rep                                                                            506

So in our case, we have an enzyme which likes the letter C. If, in folding up, an enzyme
happens to cross itself, that's okay-just think of it as going under or over itself. Notice that all
its amino acids play a role in the determination of an enzyme's tertiary structure.

                             Punctuation, Genes, and Ribosomes

Now one thing remains to he explained. Why is there a blank in box AA of the typogenetic
Code' The answer is that the duplet AA acts as a punctuation mark inside a strand, and it
signals the end of the code for an enzyme.
That is to say, one strand may code for two or more enzymes if it has one or more duplets
AA in it. For example, the strand

                                 CG GA TA CT AA AC CG A

Codes for two enzymes

                cop - ina - rpy - off
                        and
                    cut – cop

with the AA serving to divide the strand up into two "genes". The definition of gene is: that
portion of a strand which codes for a single enzyme. Note that the mere presence of AA
inside a strand does not mean that the strand codes for two enzymes. For instance, CAAG
codes for "mvr - del". The AA begins on an even-numbered unit and therefore is not read as a
duplet!
        The mechanism which reads strands and produces the enzymes which are coded
inside them is called a ribosome. (In Typogenetics, the player of the game does the work of'
the ribosomes.) Ribosomes are not in any way responsible for the tertiary structure of
enzymes, for that is entirely determined once the primary structure is created. Incidentally,
the process of translation always goes from strands to enzymes, and never in the reverse
direction.

                               Puzzle: A Typogenetical Self-Rep

Now that the rules of Typogenetics have been fully set out, you may find it interesting to
experiment with the game. In particular, it would he most interesting to devise a self-
replicating strand. This would mean something along the following lines. A single strand is
written down. A ribosome acts on it, to produce any or all of the enzymes which are coded
for in the strand. Then those enzymes are brought into contact with the original strand, and
allowed to work on it. This yields a set of "daughter strands". The daughter strands
themselves pass through the rihosomes, to yield a second generation of enzymes, which act
on the daughter strands; and the




Self-Rep and Self-Rep                                                                          507

cycle goes on and on. This can go on for any number of stages; the hope is that eventually,
among the strands which are present at some point, there
will be found two copies of the original strand (one of the copies may be, in fact, the original
strand).

                            The Central Dogma of Typogenetics

Typogenetical processes can be represented in skeletal form in a diagram (Fig. 90).




This diagram illustrates the Central Dogma of Typogenetics. It shows how strands define
enzymes (via the Typogenetic Code); and how in turn, enzymes act back on the strands
which gave rise to them, yielding new strands. Therefore, the line on the left portrays how
old information flows upwards, in the sense that an enzyme is a translation of a strand, and
contains therefore the same information as the strand, only in a different form-in particular;
in an active form. The line on the right, however, does not show information flowing
downwards; instead, it shows how new information gets created: by the shunting of symbols
in strands.
        An enzyme in Typogenetics, like a rule of inference in a formal system, blindly
shunts symbols in strands without regard to any "meaning" which may lurk in those symbols.
So there is a curious mixture of levels here. On the one hand, strands are acted upon, and
therefore play the role of data (as is indicated by the arrow on the right); on the other hand,
they also dictate the actions which are to be performed on the data, and therefore they play
the role of programs (as is indicated by the arrow on the left). It is the player of Typogenetics
who acts as interpreter and processor, of course. The two-way street which links "upper" and
"lower" levels of Typogenetics shows that, in fact, neither strands nor enzymes can be
thought of as being on a higher level than the other. By contrast, a picture of the Central
Dogma of the MIU-system looks this way:




In the MIU-system, there is a clear distinction of levels: rules of inference simply belong to a
higher level than strings. Similarly for TNT, and all formal systems.




Self-Rep and Self-Rep                                                                        508

                          Strange Loops, TNT, and Real Genetics

However, we have seen that in TNT, levels are mixed, in another sense. In fact, the
distinction between language and metalanguage breaks down: statements about the system
get mirrored inside the system. It turns out that if we make a diagram showing the
relationship between TNT and its metalanguage, we will produce something which resembles
in a remarkable way the diagram which represents the Central Dogma of Molecular Biology.
In fact, it is our goal to make this comparison in detail; but to do so, we need to indicate the
places where Typogenetics and true genetics coincide, and where they differ. Of course, real
genetics is far more complex than Typogenetics-but the "conceptual skeleton" which the
reader has acquired in understanding Typogenetics will be very useful as a guide in the
labyrinth of true genetics.

                                    DNA and Nucleotides

We begin by discussing the relationship between "strands", and DNA. The initials "DNA"
stand for "deoxyribonucleic acid". The DNA of most cells resides in the cell's nucleus, which
is a small area protected by a membrane. Gunther Stent has characterized the nucleus as the
"throne room" of the cell, with DNA acting as the ruler. DNA consists of long chains of
relatively simple molecules called nucleotides. Each nucleotide is made up of three parts: (1)
a phosphate group stripped of one special oxygen atom, whence the prefix "deoxy"; (2) a
sugar called "ribose", and (3) a base. It is the base alone which distinguishes one nucleotide
from another; thus it suffices to specify its base to identify a nucleotide. The four types of
bases which occur in DNA nucleotides are:




(Also see Fig. 91.) It is easy to remember which ones are pyrimidines because the first vowel
in "cytosine", "thymine", and "pyrimidine" is 'y'. Later, when we talk about RNA, "uracil"-
also a pyrimidine-will come in and wreck the pattern, unfortunately. (Note: Letters
representing nucleotides in real genetics will not be in the Quadrata font, as they were in
Typogenetics.)
        A single strand of DNA thus consists of many nucleotides strung together like a chain
of beads. The chemical bond which links a nucleotide to its two neighbors is very strong;
such bonds are called covalent bonds, and the "chain of beads" is often called the covalent
backbone of DNA.
  Now DNA usually comes in double strands-that is, two single strands which are paired up,
nucleotide by nucleotide (see Fig. 92). It is the bases


Self-Rep and Self-Rep                                                                       509

FIGURE 91. The four constituent bases of DNA: Adenine, Guanine, Thymine, Cytosine.
[From Hanawalt and Haynes, The Chemical Basis of Life (San Francisco: W. H. Freeman,
1973), p. 142.J


FIGURE 92. DNA structure resembles a ladder in which the side pieces consist of alternating
units of deoxyrihose and phosphate. The rungs are formed by the bases paired in a special
way, A with T and G with C, and held together respectively by two and three hydrogen
bonds. [From Hanawalt and Haynes, The Chemical Basis of Life, p. 142.




Self-Rep and Self-Rep                                                                  510

which are responsible for the peculiar kind of pairing which takes place between strands.
Each base in one strand faces a complementary base in the other strand, and binds to it. The
complements are as in Typogenetics: A pairs up with T, and C with G. Always one purine
pairs up with a pyrimidine.
         Compared to the strong covalent bonds along the backbone, the interstrand bonds are
quite weak. They are not covalent bonds, but hydrogen bonds. A hydrogen bond arises when
two molecular complexes are aligned in such a way that a hydrogen atom which originally
belonged to one of them becomes "confused" about which one it belongs to, and it hovers
between the two complexes, vacillating as to which one to join. Because the two halves of
double-stranded DNA are held together only by hydrogen bonds, they may come apart or be
put together relatively easily; and this fact is of great import for the workings of the cell.
         When DNA forms double strands, the two strands curl around each other like
twisting vines (Fig. 93). There are exactly ten nucleotide pairs per revolution; in other words,
at each nucleotide, the "twist" is 36 degrees. Single-stranded DNA does not exhibit this kind
of coiling, for it is a consequence of the base-pairing.




FIGURE 93. Molecular model of the DNA double helix. [From Vernon M. Ingram,
Biosynthesis (Menlo Park, Calif.: W. A. Benjamin, 1972)




Self-Rep and Self-Rep                                                                       511

                              Messenger KNA and Ribosomes

As was mentioned above, in many cells, DNA, the ruler of the cell, dwells in its private
"throne room": the nucleus of the cell. But most of the "living" in a cell goes on outside of
the nucleus, namely in the cytoplasm-the "ground" to the nucleus' "figure". In particular,
enzymes, which make practically every life process go, are manufactured by ribosomes in the
cytoplasm, and they do most of their work in the cytoplasm. And just as in Typogenetics, the
blueprints for all enzymes are stored inside the strands-that is, inside the DNA, which
remains protected in its little nuclear home. So how does the information about enzyme
structure get from the nucleus to the ribosomes'
         Here is where messenger RNA-mRNA-comes in. Earlier, mRNA strands were
humorously said to constitute a kind of DNA Rapid Transit Service; by this is meant not that
mRNA physically carries DNA anywhere, but rather that it serves to carry the information, or
message, stored in the DNA in its nuclear chambers, out to the ribosomes in the cytoplasm.
How is this done? The idea is easy: a special kind of enzyme inside the nucleus faithfully
copies long stretches of the DNA's base sequence onto a new strand-a strand of messenger
RNA. This mRNA then departs from the nucleus and wanders out into the cytoplasm, where
it runs into many ribosomes which begin doing their enzyme-creating work on it.
    The process by which DNA gets copied onto mRNA inside the nucleus is called
transcription; in it, the double-stranded DNA must be temporarily separated into two single
strands, one of which serves as a template for the mRNA. Incidentally, "RNA" stands for
"ribonucleic acid", and it is very much like DNA except that all of its nucleotides possess that
special oxygen atom in the phosphate group which DNA's nucleotides lack. Therefore the
"deoxy" prefix is dropped. Also, instead of thymine, RNA uses the' base uracil, so the
information in strands of RNA can be represented by arbitrary sequences of the four letters
'A', 'C', 'G', 'U'. Now when mRNA is transcribed off of DNA, the transcription process
operates via the usual base-pairing (except with U instead of T), so that a DNA-template and
its mRNA-mate might look something like this:

DNA:             GGTAAATCAAGTCA                        (template)
mRNA:            GGCAUUUAGUCAGU                        (copy")

RNA does not generally form long double strands with itself, although it can. Therefore it is
prevalently found not in the helical form which so characterizes DNA, but rather in long,
somewhat randomly curving strands.
        Once a strand of mRNA has escaped the nucleus, it encounters those strange
subcellular creatures called "ribosomes"-but before we go on to explain how a ribosome uses
mRNA, I want to make some comments about enzymes and proteins. Enzymes belong to the
general category of biomolecules called proteins, and the job of ribosomes is to make all pro




Self-Rep and Self-Rep                                                                       512

teins, not just enzymes. Proteins which are not enzymes are much more passive kinds of
beings: many of them, for instance, are structural molecules, which means that they are like
girders and beams and so forth in buildings: they hold the cell's parts together. There are
other kinds of proteins, but for our purposes. the principal proteins are enzymes, and I will
henceforth not make a sharp distinction.

                                        Amino Acids

Proteins are composed of sequences of amino acids, which come in twenty primary varieties,
each with a three-letter abbreviation:

                 ala    alanine
                 arg    arginine
                 asn    asparagines
                 asp    aspartic acid
                 cys    cysteine
                 gln    glutamine
                 glu    glutamic acid
                 gly    glycine
                 his    histidine
                 He     isoleucine
                 leu    leucine
                 lys    lysine
                 met    methionine
                 phe    phenylalanine
                 pro    praline
                 ser    serine
                 thr    threonine
                 trp    tryptophan
                 tyr    tyrosine
                 val    valine

Notice the slight numerical discrepancy with Typogenetics, where we had only fifteen
"amino acids" composing enzymes. An amino acid is a small molecule of roughly the same
complexity as a nucleotide; hence the building blocks of proteins and of nucleic acids (DNA,
RNA) are roughly of the same size. However, proteins are composed of much shorter
sequences of components: typically, about three hundred amino acids make a complete
protein, whereas a strand of DNA can consist of hundreds of thousands or millions of
nucleotides.

                              Ribosomes and Tape Recorders

Now when a strand of mRNA, after its escape into the cytoplasm, encounters a ribosome, a
very intricate and beautiful process called translation takes place. It could be said that this
process of translation is at the very heart of




Self-Rep and Self-Rep                                                                      513

all of life, and there are many mysteries connected with it. But in essence it is easy to
describe. Let us first give a picturesque image, and then render it more precise. Imagine the
mRNA to be like a long piece of magnetic recording tape, and the ribosome to be like a tape
recorder. As the tape passes through the playing head of the recorder, it is "read" and
converted into music, or other sounds. Thus magnetic markings are "translated" into notes.
Similarly, when a "tape" of mRNA passes through the "playing head" of a ribosome, the
"notes" which are produced are amine acids, and the "pieces of music" which they make up
are proteins. This is what translation is all about: it is shown in Figure 96.


                                      The Genetic Code

But how can a ribosome produce a chain of amino acids when it is reading a chain of
nucleotides This mystery was solved in the early 1960's by the efforts of a large number
of people, and at the core of the answer lies the Genetic Code-a mapping from triplets of
nucleotides into amino acids (see Fig. 94). This is in spirit extremely similar to the
Typogenetic     Code,      except     that     here,    three consecutive     bases     (or
nucleotides)                                                             form a codon,
whereas there,                        CUA GAU                            only two were
needed. Thus                         ,Cu Ag Au                           there must be
4x4x4 (equals                                                            64)      different
entries in the                 A typical segment   of mRNA               table, instead of
sixteen.     A                   read first as two triplets              ribosome clicks
down a strand                  (above), and second as three              of RNA three
nucleotides at                 duplets (below): an example               a time-which is
to say, one.                   of hemiolia in biochemistry               codon at a time
-and      each                                                           time it does so,
it appends a single new amino acid to the protein it is presently manufacturing. Thus, a
protein comes out of the ribosome amino acid by amino acid.

                                      Tertiary Structure

 However, as a protein emerges from a ribosome, it is not only getting longer and longer, but
it is also continually folding itself up into an extraordinary three-dimensional shape, very
much in the way that those funny little Fourth-of-July fireworks called "snakes"
simultaneously grow longer and curl up, when they are lit. This fancy shape is called the
protein's tertiary structure (Fig. 95), while the amino acid sequence per se is called the
primary structure of the protein. The tertiary structure is implicit in the primary structure, just
as in Typogenetics. However, the recipe for deriving the tertiary structure, if you know only
the primary structure, is by far more complex than that given in Typogenetics. In fact, it is
one of the outstanding problems of contemporary molecular biology to figure out some rules
by which the tertiary structure of a protein can be predicted if only its primary structure is
known.




Self-Rep and Self-Rep                                                                          514

The Genetic Code.


                              U   C       A        G
                      phe         ser   tyr     cys          U

                      phe         ser   tyr                  C
                                                cys
              U         leu       ser   punt.   punt.        A
                        leu       ser   punc.                G
                                                 trp


                        leu       pro   his      arg         U
                        leu       pro   his      arg         C
              C         leu       pro   gin      arg         A

                        Ieu       pro            arg         G
                                        gln


                        ile       thr   asn      ser         U
                        ile       thr   asn      ser         C
              A         ile       thr   lys      arg         A
                      met         thr   lys      arg         G


              G         val       ala   asp     gly          U
                        val       ala   asp     gly          C
                        val       ala   glu     gly          A

                        val       ala   glu     gly          G




FIGURE 94. The Genetic Code, by which each triplet in a strand of messenger RNA codes
for one of twenty amino acids (or a punctuation mark).

                     Reductionistic Explanation of Protein Function

Another discrepancy between Typogenetics and true genetics-and this is probably the most
serious one of all-is this: whereas in Typogenetics, each component amino acid of an enzyme
is responsible for some specific "piece of the action", in real enzymes, individual amino acids
cannot be assigned such clear roles. It is the tertiary structure as a whole which determines
the mode in which an enzyme will function; there is no way one can say, "This




Self-Rep and Self-Rep                                                                      515

amino acid's presence means that such-and-such an operation will get performed". In other
words, in real genetics, an individual amino acid's contribution to the enzyme's overall
function is not "context-free". However, this fact should not be construed in any way as
ammunition for an anti reductionist argument to the effect that "the whole [enzyme] cannot
be explained as the sum of its parts". That would he wholly unjustified. What is justified is
rejection of the simpler claim that "each amino acid contributes to the sum in a manner which
is independent of the other amino acids present". In other words, the function of a protein
cannot be considered to be built up from context-free functions of its parts; rather, one must
consider how the parts interact. It is still possible in principle to write a computer program
which takes as input the primary structure of a protein,

FIGURE 95. The structure of myoglobin, deduced from high-resolution X-ray data. The
large-scale "twisted pipe" appearance is the tertiary structure; the finer helix inside-the
"alpha helix"-is the secondary structure. [From A. Lehninger, Biochemistry]




Self-Rep and Self-Rep                                                                     516

and firstly determines its tertiary structure, and secondly determines the function of the
enzyme. This would be a completely reductionistic explanation of the workings of proteins,
but the determination of the "sum" of the parts would require a highly complex algorithm.
The elucidation of the function of an enzyme, given its primary, or even its tertiary, structure,
is another great problem of contemporary molecular biology.
         Perhaps, in the last analysis, the function of the whole enzyme can be considered to
be built up from functions of parts in a context-free manner, but where the parts are now
considered to be individual particles, such as electrons and protons, rather than "chunks",
such as amino acids. This exemplifies the "Reductionist's Dilemma": In order to explain
everything in terms of context free sums, one has to go down to the level of physics; but then
the number of particles is so huge as to make it only a theoretical "in-principle" kind of thing.
So, one has to settle for a context-dependent sum, which has two disadvantages. The first is
that the parts are much larger units, whose behavior is describable only on a high level, and
therefore indeterminately. The second is that the word "sum" carries the connotation that
each part can be assigned a simple function and that the function of the whole is just a
context-free sum of those individual functions. This just cannot be done when one tries to
explain a whole enzyme's function, given its amino acids as parts. But for better or for worse,
this is a general phenomenon which arises in the explanations of complex systems. In order
to acquire an intuitive and manageable understanding of how parts interact-in short, in order
to proceed-one often has to sacrifice the exactness yielded by a microscopic, context-free
picture, simply because of its unmanageability. But one does not sacrifice at that time the
faith that such an explanation exists in principle.


                               Transfer RNA and Ribosomes

Returning, then, to ribosomes and RNA and proteins, we have stated that a protein is
manufactured by a ribosome according to the blueprint carried from the DNA's "royal
chambers" by its messenger, RNA. This seems to imply that the ribosome can translate from
the language of codons into the language of amino acids, which amounts to saying that the
ribosome "knows" the Genetic Code. However, that amount of information is simply not
present in a ribosome. So how does it do it? Where is the Genetic Code stored? The curious
fact is that the Genetic Code is stored-where else?-in the DNA itself. This certainly calls for
some explanation.
         Let us back off from a total explanation for a moment, and give a partial explanation.
There are, floating about in the cytoplasm at any given moment, large numbers of four-leaf-
clover-shaped molecules; loosely fastened (i.e., hydrogen-bonded) to one leaf is an amino
acid, and on the opposite leaf there is a triplet of nucleotides called an anticodon. For our
purposes, the other two leaves are irrelevant. Here is how these "clovers" are used by the
ribosomes in their production of proteins. When a new




Self-Rep and Self-Rep                                                                        517

FIGURE 96. A section of mRNA passing through a ribosome. Floating nearby are tRNA
molecules, carrying amino acids which are stripped off by the ribosome and appended to the
growing protein. The Genetic Code is contained in the tRNA molecules, collectively. Note
how the base-pairing (A-U, C-G) is represented by interlocking letter-forms in the diagram.
[Drawing by Scott E. Kim]

codon of mRNA clicks into position in the ribosome's "playing head", the ribosome reaches
out into the cytoplasm and latches onto a clover whose anticodon is complementary to the
mRNA codon. Then it pulls the clover into such a position that it can rip off the clover's
amino acid, and stick it covalently onto the growing protein. (Incidentally, the bond between
an amino acid and its neighbor in a protein is a very strong covalent bond, called a "peptide
bond". For this reason, proteins are sometimes called "polypeptides".) Of course it is no
accident that the "clovers" carry the proper amino acids, for they have all been manufactured
according to precise instructions emanating from the "throne room".




Self-Rep and Self-Rep                                                                    518

The real name for such a clover is transfer RNA. A molecule of tRNA is quite small-about
the size of a very small protein-and consists of a chain of about eighty nucleotides. Like
mRNA, tRNA molecules are made by transcription off of the grand cellular template, DNA.
However, tRNA's are tiny by comparison with the huge mRNA molecules, which may
contain thousands of nucleotides in long, long chains. Also, tRNA's resemble proteins (and
are unlike strands of mRNA) in this respect: they have fixed, well-defined tertiary structures-
determined by their primary structure. A tRNA molecule's tertiary structure allows precisely
one amino acid to bind to its amino-acid site: to be sure, it is that one dictated according to
the Genetic Code by the anticodon on the opposite arm. A vivid image of the function of
tRNA molecules is as flashcards floating in a cloud around a simultaneous interpreter, who
snaps one out of the air-invariably the right one!-whenever he needs to translate a word. In
this case, the interpreter is the ribosome, the words are codons, and their translations are
amino acids.
         In order for the inner message of DNA to get decoded by the ribosomes, the tRNA
flashcards must be floating about in the cytoplasm. In some sense, the tRNAs contain the
essence of the outer message of the DNA, since they are the keys to the process of
translation. But they themselves came from the DNA. Thus, the outer message is trying to be
part of the inner message, in a way reminiscent of the message-in-a-bottle which tells what
language it is written in. Naturally, no such attempt can be totally successful: there is no way
for the DNA to hoist itself by its own bootstraps. Some amount of knowledge of the Genetic
Code must already be present in the cell beforehand, to allow the manufacture of those
enzymes which transcribe tRNA's themselves off of the master copy of DNA. And this
knowledge resides in previously manufactured tRNA molecules. This attempt to obviate the
need for any outer message at all is like the Escher dragon, who tries as hard as he can,
within the context of the two-dimensional world to which he is constrained, to be
threedimensional. He seems to go a long way-but of course he never makes it, despite the
fine imitation he gives of three-dimensionality.

                            Punctuation and the Reading Frame

How does a ribosome know when a protein is done? Just as in Typogenetics, there is a signal
inside the mRNA which indicates the termination or initiation of a protein. In fact, three
special codons-UAA, CAG, UGA act as punctuation marks instead of coding for amino
acids. Whenever such a triplet clicks its way into the "reading head" of a ribosome, the
ribosome releases the protein under construction and begins a new one.
        Recently, the entire genome of the tiniest known virus, φ)X174, has been laid bare.
One most unexpected discovery was made en route: some of its nine genes overlap-that is,
two distinct proteins are coded for by the same stretch of DNA! There is even one gene
contained entirely inside another!




Self-Rep and Self-Rep                                                                       519

This is accomplished by having the reading frames of the two genes shifted relative to each
other, by exactly one unit. The density of information packing in such a scheme is incredible.
This is, of course, the inspiration behind the strange "5/17 haiku" in Achilles' fortune cookie,
in the Canon byIntervallic Augmentation,

                                               Recap

In brief, then, this picture emerges: from its central throne, DNA sends off long strands of
messenger RNA to the ribosomes in the cytoplasm; and the ribosomes, making use of the
"flashcards" of tRNA hovering about them, efficiently construct proteins, amino acid by
amino acid, according to the blueprint contained in the mRNA. Only the primary structure of
the proteins is dictated by the DNA; but this is enough, for as they emerge from the
ribosomes, the proteins "magically" fold up into complex conformations which then have the
ability to act as powerful chemical machines.

                  Levels of Structure and Meaning in Proteins and Music

We have been using this image of ribosome as tape recorder, mRNA as tape, and protein as
music. It may seem arbitrary, and yet there are some beautiful parallels. Music is not a mere
linear sequence of notes. Our minds perceive pieces of music on a level far higher than that.
We chunk notes into phrases, phrases into melodies, melodies into movements, and
movements into full pieces. Similarly, proteins only make sense when they act as chunked
units. Although a primary structure carries all the information for the tertiary structure to be
created, it still "feels" like less, for its potential is only realized when the tertiary structure is
actually physically created.
  Incidentally, we have been referring only to primary and tertiary structures, and you may
well wonder whatever happened to the secondary structure. Indeed, it exists, as does a
quaternary structure, as well. The folding-up of a protein occurs at more than one level.
Specifically, at some points along the chain of amino acids, there may be a tendency to form
a kind of helix, called the alpha helix (not to be confused with the DNA double helix). This
helical twisting of a protein is on a lower level than its tertiary structure. This level of
structure is visible in Figure 95. Quaternary structure can be directly compared with the
building of 'a musical piece out of independent movements, for it involves the assembly of'
several distinct polypeptides, already in their full-blown tertiary beauty, into a larger
structure. The binding of these independent chains is usually accomplished by hydrogen
bonds, rather than covalent bonds; this is of course just as with pieces of music composed of
several movements, which are far less tightly bound to each other than they are internally, but
which nevertheless form a tight "organic" whole.
  The four levels of primary, secondary, tertiary, and quaternary structure can also be
compared to the four levels of the MU-picture (Fig. 60) in




Self-Rep and Self-Rep                                                                            520

Self-Rep and Self-Rep   521

the Prelude, Ant Fugue. The global structure-consisting of the letters 'M' and 'U'-is its
quaternary structure; then each of those two parts has a tertiary structure, consisting of
"HOLISM" or "REDUCTIONISM"; and then the opposite word exists on the secondary
level, and at bottom, the primary structure is once again the word "MU", over and over again.

                         Polyribosomes and Two-Tiered Canons

Now we come to another lovely parallel between tape recorders translating tape into music
and ribosomes translating mRNA into proteins. Imagine a collection of many tape recorders,
arranged in a row, evenly spaced. We might call this array a "polyrecorder". Now imagine a
single tape passing serially through the playing heads of all the component recorders. If the
tape contains a single long melody, then the output will be a many-voiced canon, of course,
with the delay determined by the time it takes the tape to get from one tape recorder to the
next. In cells, such "molecular canons" do indeed exist, where many ribosomes, spaced out in
long lines-forming what is called a polyribosome-all "play" the same strand of mRNA,
producing identical proteins, staggered in time (see Fig. 97).
        Not only this, but nature goes one better. Recall that mRNA is made by transcription
off of DNA; the enzymes which are responsible for this process are called RNA polymerases
("-ase" is a general suffix for enzymes). It happens often that a series of RNA polymerases
will be at work in parallel on a single strand of DNA, with the result that many separate (but
identical) strands of mRNA are being produced, each delayed with respect to the other by the
time required for the DNA to slide from one RNA polymerase to the next. At the same time,
there can be several different ribosomes working on each of the parallel emerging mRNA's.
Thus one arrives at a double-decker, or two-tiered, "molecular canon" (Fig. 98). The
corresponding image in music is a rather fanciful but amusing scenario: several

FIGURE 98. Here, an even more complex scheme. Not just one but several strands of
mRNA, all emerging by transcription from a single strand of DNA, are acted upon by
polyribosomes. The result is a two-tiered molecular canon. [From Hanawalt and Haynes, The
Chemical Basis of Life, p. 271]




Self-Rep and Self-Rep                                                                     522

different copyists are all at work simultaneously, each one of them copying the same original
manuscript from a clef which flutists cannot read into a clef which they can read. As each
copyist finishes a page of the original manuscript, he passes it on to the next copyist, and
starts transcribing a new page himself. Meanwhile, from each score emerging from the pens
of the copyists, a set of flutists are reading and tooting the melody, each flutist delayed with
respect to the others who are reading from the same sheet. This rather wild image gives,
perhaps, an idea of some of the complexity of the processes which are going on in each and
every cell of your body during every second of every day ...

                     Which Came First-The Ribosome or the Protein?

We have been talking about these wonderful beasts called ribosomes; but what are they
themselves composed of? How are they made? Ribosomes are composed of two types of
things: (1) various kinds of proteins, and (2) another kind of RNA, called ribosomal RNA
(rRNA). Thus, in order for a ribosome to be made, certain kinds of proteins must be present,
and rRNA must be present. Of course, for proteins to be present, ribosomes must be there to
make them. So how do you get around the vicious circle? Which comes first-the ribosome or
the protein? Which makes which? Of course there is no answer because one always traces
things back to previous members of the same class just as with the chicken-and-the-egg
question-until everything vanishes over the horizon of time. In any case, ribosomes are made
of two pieces, a large and a small one, each of which contains some rRNA and some proteins.
Ribosomes are about the size of large proteins; they are much much smaller than the strands
of mRNA which they take as input, and along which they move.

                                      Protein Function

We have spoken somewhat of the structure of proteins-specifically enzymes; but we have not
really mentioned the kinds of tasks which they perform in the cell, nor how they do them. All
enzymes are catalysts, which means that in a certain sense, they do no more than selectively
accelerate various chemical processes in the cell, rather than make things happen which
without them never could happen. An enzyme realizes certain pathways out of the myriad
myriad potentialities. Therefore, in choosing which enzymes shall be present, you choose
what shall happen and what shall not happen-despite the fact that, theoretically speaking,
there is a nonzero probability for any cellular process to happen spontaneously, without the
aid of catalysts.
   Now how do enzymes act upon the molecules of the cell? As has been mentioned,
enzymes are folded-up polypeptide chains. In every enzyme, there is a cleft or pocket or
some other clearly-defined surface feature where the enzyme hinds to some other kind of
molecule. This location is




Self-Rep and Self-Rep                                                                       523

called its active site, and any molecule which gets bound there is called a substrate. Enzymes
may have more than one active site, and more than one substrate. Just as in Typogenetics,
enzymes are indeed very choosy about what they will operate upon. The active site usually is
quite specific, and allows just one kind of molecule to bind to it, although there are
sometimes "decoys"-other molecules which can fit in the active site and clog it up, fooling
the enzyme and in fact rendering it inactive.
         Once an enzyme and its substrate are bound together, there is some disequilibrium of
electric charge, and consequently charge-in the form of electrons and protons-flows around
the bound molecules and readjusts itself. By the time equilibrium has been reached, some
rather profound chemical changes may have occurred to the substrate. Some examples are
these: there may have been a "welding", in which some standard small molecule got tacked
onto a nucleotide, amino acid, or other common cellular molecule; a DNA strand may have
been "nicked" at a particular location; some piece of a molecule may have gotten lopped off;
and so forth. In fact, bio-enzymes do operations on molecules which are quite similar to the
typographical operations which Typo-enzymes perform. However, most enzymes perform
essentially only a single task, rather than a sequence of tasks. There is one other striking
difference between Typoenzymes and bio-enzymes, which is this: whereas Typo-enzymes
operate only on strands, bio-enzymes can act on DNA, RNA, other proteins, ribosomes, cell
membranes-in short, on anything and everything in the cell. In other words, enzymes are the
universal mechanisms for getting things done in the cell. There are enzymes which stick
things together and take them apart and modify them and activate them and deactivate them
and copy them and repair them and destroy them .. .
         Some of the most complex processes in the cell involve "cascades" in which a single
molecule of some type triggers the production of a certain kind of enzyme; the manufacturing
process begins and the enzymes which come off the "assembly line" open up a new chemical
pathway which allows a second kind of enzyme to be produced. This kind of thing can go on
for three or four levels, each newly produced type of enzyme triggering the production of
another type. In the end a "shower" of copies of the final type of enzyme is produced, and all
of the copies go off and do their specialized thing, which may be to chop up some "foreign"
DNA, or to help make some amino acid for which the cell is very "thirsty", or whatever.


                      Need for a Sufficiently Strong Support System

Let us describe nature's solution to the puzzle posed for Typogenetics: "What kind of strand
of DNA can direct its own replication?" Certainly not every strand of DNA is inherently a
self-rep. The key point is this: any strand which wishes to direct its own copying must
contain directions for assembling precisely those enzymes which can carry out the task. Now
it is futile to hope that a strand of DNA in isolation could be a self-rep; for in




Self-Rep and Self-Rep                                                                     524

order for those potential proteins to be pulled out of the DNA, there must not only be
ribosomes, but also RNA polymerase, which makes the mRNA that gets transported to the
ribosomes. And so we have to begin by assuming a kind of "minimal support system" just
sufficiently strong that it allows transcription and translation to be carried out. This minimal
support system will thus consist in (1) some proteins, such as RNA polymerase, which allow
mRNA to be made from DNA, and (2) some ribosomes.

                                 How DNA Self-Replicates

It is not by any means coincidental that the phrases "sufficiently strong support system" and
"sufficiently powerful formal system" sound alike. One is the precondition for a self-rep to
arise, the other for a self-ref to arise. In fact there is in essence only one phenomenon going
on in two very different guises, and we shall explicitly map this out shortly. But before we do
so, let us finish the description of how a strand of DNA can be a self-rep.
         The DNA must contain the codes for a set of proteins which will copy it. Now there
is a very efficient and elegant way to copy a double-stranded piece of DNA, whose two
strands are complementary. This involves two steps:

      (1) unravel the two strands from each other;
      (2) mate" a new strand to each of the two new single strands.

This process will create two new double strands of DNA, each identical to the original one.
Now if our solution is to be based on this idea, it must involve a set of proteins, coded for in
the DNA itself, which will carry out these two steps.
        It is believed that in cells, these two steps are performed together in a coordinated
way, and that they require three principal enzymes: DNA endonuclease, DNA polymerase,
and DNA ligase. The first is an "unzipping enzyme": it peels the two original strands apart
for a short distance, and then stops. Then the other two enzymes come into the picture. The
DNA polymerase is basically a copy-and-move enzyme: it chugs down the short single
strands of DNA, copying them complementarily in a fashion reminiscent of the Copy mode
in Typogenetics. In order to copy, it draws on raw materials-specifically nucleotides-which
are floating about in the cytoplasm. Because the action proceeds in fits and starts, with some
unzipping and some copying each time, some short gaps are created, and the DNA ligase is
what plugs them up. The process is repeated over and over again. This precision three-
enzyme machine proceeds in careful fashion all the way down the length of the DNA
molecule, until the whole thing has been peeled apart and simultaneously replicated, so that
there are now two copies of it.




Self-Rep and Self-Rep                                                                       525

                  Comparison of DNA's Self-Rep Method with Quining

Note that in the enzymatic action on the DNA strands, the fact that information is stored in
the DNA is just plain irrelevant; the enzymes are merely carrying out their symbol-shunting
functions, just like rules of inference in the MMIU-system. It is of no interest to the three
enzymes that at some point they are actually copying the very genes which coded for them.
The DNA, to them, is just a template without meaning or interest.
         It is quite interesting to compare this with the Quine sentence's method of describing
how to construct a copy of itself. There, too, one has a sort of "double strand"-two copies of
the same information, where one copy acts as instructions, the other as template. In DNA, the
process is vaguely parallel, since the three enzymes (DNA endonuclease, DNA polymerase,
DNA ligase) are coded for in just one of the two strands, which therefore acts as program,
while the other strand is merely a template. The parallel is not perfect, for when the copying
is carried out, both strands are used as template, not just one. Nevertheless, the analogy is
highly suggestive. There is a biochemical analogue to the use-mention dichotomy: when
DNA is treated as a mere sequence of chemicals to be copied, it is like mention of
typographical symbols; when DNA is dictating what operations shall he carried out, it is like
use of typographical symbols.


                                 Levels of Meaning of DNA

There are several levels of meaning which can be read from a strand of DNA, depending on
how big the chunks are which you look at, and how powerful a decoder you use. On the
lowest level, each DNA strand codes for an equivalent RNA strand-the process of decoding
being transcription. If one chunks the DNA into triplets, then by using a "genetic decoder",
one can read the DNA as a sequence of amino acids. This is translation (on top of
transcription). On the next natural level of the hierarchy, DNA is readable as a code for a set
of proteins. The physical pulling-out of proteins from genes is called gene expression.
Currently, this is the highest level at which we understand what DNA means.
        However, there are certain to be higher levels of DNA meaning which are harder to
discern. For instance, there is every reason to believe that the DNA of, say, a human being
codes for such features as nose shape, music talent, quickness of reflexes, and so on. Could
one, in principle, learn to read off such pieces of information directly from a strand of DNA,
without going through the actual physical process of epigenesis-the physical pulling-out of
phenotype from genotype Presumably, yes, since-in theory-one could have an incredibly
powerful computer program simulating the entire process, including every cell, every protein,
every tiny feature involved in the replication of DNA, of cells, to the bitter end. The output of
such a pseudo-epigenesis program would be a high-level description of the phenotype.




Self-Rep and Self-Rep                                                                        526

  There is another (extremely faint) possibility: that we could learn to read the phenotype off
of the genotype without doing an isomorphic simulation of the physical process of
epigenesis, but by finding some simpler sort of decoding mechanism. This could be called
"shortcut pseudoepigenesis". Whether shortcut or not, pseudo-epigenesis is, of course, totally
beyond reach at the present time-with one notable exception: in the species Felis catus, deep
probing has revealed that it is indeed possible to read the phenotype directly off of the
genotype. The reader will perhaps better appreciate this remarkable fact after directly
examining the following typical section of the DNA of Felis catus:

                   CATCATCATCATCATCATCATCATCATCAT ...

Below is shown a summary of the levels of DNA-readability, together with the names of the
different levels of decoding. DNA can be read as a sequence of:

bases (nucleotides)                                                    transcription
amino acids                                                            translation
proteins (primary structure)                                        . gene expression
proteins (tertiary structure)
protein clusters                                       higher levels of gene expression
…
….      unknown levels of DNA meaning
….
(N-1) ????
(N) physical, mental, and
psychological traits                                           pseudo-epigenesis


                                    The Central Dogmap
.
With this background, now we are in a position to draw an elaborate comparison between F.
Crick's "Central Dogma of Molecular Biology" (.DOGMA I) upon which all cellular
processes are based; and what I, with poetic license, call the "Central Dogma of
Mathematical Logic" (.DOGMA II), upon which G6del's Theorem is based. The mapping
from one onto the other is laid out in Figure 99 and the following chart, which together
constitute the Central Dogmap.

FIGURE 99. The Central Dogmap. An analogy is established between two fundamental
Tangled Hierarchies: that of molecular biology and that of mathematical logic.




Self-Rep and Self-Rep                                                                      527

Self-Rep and Self-Rep   528

        Note the base-pairing of A and T (Arithmetization and Translation), as well as of G
and C (Godel and Crick). Mathematical logic gets the purine side, and molecular biology gets
the pyrimidine side.
        To complete the esthetic side of this mapping, I chose to model my Godel-numbering
scheme on the Genetic Code absolutely faithfully. In fact, under the following
correspondence, the table of the Genetic Code becomes the table of the Godel Code:




Each amino acid-of which there are twenty-corresponds to exactly one symbol of TNT-of
which there are twenty. Thus, at last, my motive for concocting "austere TNT" comes out-so
that there would be exactly twenty symbols! The Godel Code is shown in Figure 100.
Compare it with the Genetic Code (Fig. 94).
         There is something almost mystical in seeing the deep sharing of such an abstract
structure by these two esoteric, yet fundamental, advances in knowledge achieved in our
century. This Central Dogmap is by no means a rigorous proof of identity of the two theories;
but it clearly shows a profound kinship, which is worth deeper exploration.

                          Strange Loops in the Central Dogmap

One of the more interesting similarities between the two sides of the map is the way in which
"loops" of arbitrary complexity arise on the top level of both: on the left, proteins which act
on proteins which act on proteins and so on, ad infinitum; and on the right, statements about
statements about statements of meta-TNT and so on, ad infinitum. These are like
heterarchies, which we discussed in Chapter V, where a sufficiently complex substratum
allows high-level Strange Loops to occur and to cycle around, totally sealed off from lower
levels. We will explore this idea in greater detail in Chapter XX.
        Incidentally, you may be wondering about this question: "What, according to the
Central Dogmap, is Godel's Incompleteness Theorem itself mapped onto?" This is a good
question to think about before reading ahead.

                    The Central Dogmap and the Contracrostipuntus

  It turns out that the central dogmap is quite similar to the mapping that was laid out in
Chapter IV between the Contracrostipunctus and Godel’s Theorem. One can therefore drew
parallels between all three systems.




Self-Rep and Self-Rep                                                                      529

FIGURE 100. The Godel Code. Under this Godel-numbering scheme, each T V7- s'mbol gets
one or more codons. The small ovals show how this table subsumes the earlier
Godelnumhering table of Chapter IX.

(1) formal systems and strings
(2) cells and strands of DNA
(3) record players and records


In the following chart, the mapping between systems 2 and 3 is explained carefully




Self-Rep and Self-Rep                                                                530

The analogue of Godel's Theorem is seen to be a peculiar fact, probably little useful to
molecular biologists (to whom it is likely quite obvious):

    It is always possible to design a strand of DNA which, if injected into a cell, would,
    upon being transcribed, cause such proteins to be manufactured as would destroy the cell
    (or the DNA), and thus result in the non-reproduction of that DNA

This conjures tip a somewhat droll scenario, at least if taken in light of evolution: an invading
species of virus enters a cell by some surreptitious



Self-Rep and Self-Rep                                                                        531

FIGURE 101. The T4 bacterial virus is an assembly of protein components (a). The "head" is
a protein membrane, shaped like a kind of prolate irosahedron with thirty facets and filled
with DNA. It is attached by a neck to a tail consisting of a hollow core surrounded by a
contractile sheathh and based on a spiked end plate to which six fibers are attached. The
spikes and fibers affix the virus to a bacterial cell wall (h). The sheath contracts, driving the
core through the wall, and viral DNA enters the cell. [From Hanawalt and Haynes, The
Chemical Basis of Life, p. 230.1


means, and then carefully ensures the manufacture of proteins which will have the effect of
destroying the virus itself! It is a sort of suicide-or Epimenides sentence, if you will-on the
molecular level. Obviously it would not prove advantageous from the point of view of
survival of the species. However, it demonstrates the spirit, if not the letter, of the
mechanisms of protection and subversion which cells and their invaders have developed.

                                         E. Coli vs. T4

Let us consider the biologists' favorite cell, that of the bacterium Escherichia coli (no relation
to M. C. Escher), and one of their favorite invaders of that cell: the sinister and eerie T4
phage, pictures of which you can see in Figure 101. (Incidentally, the words "phage" and
"virus" are synonymous and mean "attacker of bacterial cells".) The weird tidbit looks like a
little like a cross between a LEM (Lunar Excursion Module) and a mosquito-and it is much
more sinister than the latter. It has a "head" wherein is stored all its "knowledge"-namely its
DNA; and it has six "legs" wherewith to fasten itself to the cell it has chosen to invade; and it
has a "stinging tube" (more properly called its "tail") like a mosquito. The major difference is
that unlike a mosquito, which uses its stinger for sucking blood, the T4 phage uses its stinger
for injecting its hereditary substance into the cell against the will of its victim. Thus the
phage commits "rape" on a tiny scale.




Self-Rep and Self-Rep                                                                         532

FIGURE 102. Viral infection begins when viral DNA enters a bacterium. Bacterial DNA is
disrupted and viral DNA replicated. Synthesis of viral structural proteins and their assembly
into virus continues until the cell bursts, releasing particles. [From Hanawalt and Haynes,
The Chemical Basis of Life, p. 230.]

                                 A Molecular Trojan Horse

What actually happens when the viral DNA enters a cell? The virus "hopes", to speak
anthropomorphically, that its DNA will get exactly the same treatment as the DNA of the
host cell. This would mean getting transcribed and translated, thus allowing it to direct the
synthesis of its own special proteins, alien to the host cell, which will then begin to do their
thing. This amounts to secretly transporting alien proteins "in code" (viz., the Genetic Code)
into the cell, and then "decoding" (viz., producing) them. In a way this resembles the story of
the Trojan horse, according to which hundreds of soldiers were sneaked into Troy inside a
harmless seeming giant wooden horse; but once inside the city, they broke loose and captured
it. The alien proteins, once they have been "decoded" (synthesized) from their carrier DNA,
now jump into action. The sequence of actions directed by the T4 phage has been carefully
studied, and is more or less as follows (see also Figs. 102 and 103):

Time elapsed Action taking place

0 min.         Injection of viral DNA.

1 min.         Breakdown of host DNA. Cessation of production of native proteins and
               initiation of production of alien (T4) proteins. Among the earliest produced
               proteins are those which direct the replication of the alien (T4) DNA.

5 min.         Replication of viral DNA begins.

8 min.         Initiation of production of structural proteins which will form the "bodies" of
               new phages.




Self-Rep and Self-Rep                                                                       533

Self-Rep and Self-Rep   534

13 min.         First complete replica of T4 invader is produced.
25 min.         Lysozyme (a protein) attacks host cell wall. breaking open the bacterium, and
                the "bicentuplets" emerge.

Thus, when a T4 phage invades an E. coli cell, after the brief span of about twenty-four or
twenty-five minutes, the cell has been completely subverted, and breaks open. Out pop about
two hundred exact copies of the original virus-"bicentuplets"-ready to go attack more
bacterial cells, the original cell having been largely consumed in the process.
        Although from a bacterium's point of view this kind of thing is a deadly serious
menace, from our large-scale vantage point it can be looked upon as an amusing game
between two players: the invader, or "T" player (named after the T-even class of phages,
including the T2, T4, and others), and the "C" player (standing for "Cell"). The objective of
the T player is to invade and take over the cell of the C player from within, for the purpose of
reproducing itself. The objective of the C player is to protect itself and destroy the invader.
When described this way, the molecular TC-game can be seen to be quite parallel to the
macroscopic TC-game described in the preceding Dialogue. (The reader can doubtless figure
out which player-T or C-corresponds to the Tortoise, and which to the Crab.)

                               Recognition, Disguises, Labeling
.
This "game" emphasizes the fact that recognition is one of the central themes of cellular and
subcellular biology. How do molecules (or higher-level structures) recognize each other? It is
essential for the functioning of enzymes that they should be able to latch onto special
"binding sites" on their substrates; it is essential that a bacterium should be able to distinguish
its own DNA from that of phages; it is essential that two cells should be able to recognize
each other and interact in a controlled way. Such recognition problems may remind you of
the original, key problem about formal systems: How can you tell if a string has, or does not
have, some property such as theoremhood? Is there a decision procedure? This kind of
question is not restricted to mathematical logic: it permeates computer science and, as we are
seeing, molecular biology.
        The labeling technique described in the Dialogue is in fact one of E. colt's tricks for
outwitting its phage invaders. The idea is that strands of DNA can be chemically labeled by
tacking on a small molecule-methyl-to various nucleotides. Now this labeling operation does
not change the usual biological properties of the DNA; in other words, methylated (labeled)
DNA can be transcribed just as well as unmethylated (unlabeled) DNA, and so it can direct
the synthesis of proteins. But if the host cell has some special




Self-Rep and Self-Rep                                                                          535

mechanisms for examining whether DNA is labeled or not, then the label may make all the
difference in the world. In particular, the host cell may have an enzyme system which looks
for unlabeled DNA, and destroys any that it finds by unmercifully chopping it to pieces. In
that case, woe to all unlabeled invaders!
         The methyl labels on the nucleotides have been compared to serifs on letters. Thus,
using this metaphor, we could say that the E. colt cell is looking for DNA written in its
"home script", with its own particular typeface-and will chop up any strand of DNA written
in an "alien" typeface. One counterstrategy, of course, is for phages to learn to label
themselves, and thereby become able to fool the cells which they are invading into
reproducing them.
         This TC-battle can continue to arbitrary levels of complexity, but we shall not pursue
it further. The essential fact is that it is a battle between a host which is trying to reject all
invading DNA, and a phage which is trying to infiltrate its DNA into some host which will
transcribe it into mRNA (after which its reproduction is guaranteed). Any phage DNA which
succeeds in getting itself reproduced this way can be thought of as having this high-level
interpretation: "I Can Be Reproduced in'Cells of Type X". This is to be distinguished from
the evolutionarily pointless kind of phage mentioned earlier, which codes for proteins that
destroy it, and whose high-level interpretation is the self-defeating sentence: "I Cannot Be
Reproduced in Cells of Type X".

                                Henkin Sentences and Viruses

Now both of these contrasting types of self-reference in molecular biology have their
counterparts in mathematical logic. We have already discussed the analogue of the self-
defeating phages-namely, strings of the G6del type, which assert their own unproducibility
within specific formal sstems. But one can also make a counterpart sentence to a real phage:
the' phage asserts its own producibility in a specific cell, and the sentence asserts its own
producibility in a specific formal system. Sentences of this type are called Henkin sentences,
after the mathematical logician Leon Henkin. They can be constructed exactly along the lines
of Godel sentences, the only difference being the omission of a negation. One begins with an
"uncle", of course:

3a:3a':<TNT-PROOF-PAIR{a,a'}-°ARITHMOQUINE{a",a'}>

and then proceeds by the standard trick. Say the Godel number of the above "uncle" is h.
Now by arithmoquining this very uncle, you get a Henkin sentence:

3a:3a':<TNT-PROOF-PAIR{a,a'} ^ ARITHMOQUINE{SSS ... SSSO/a",a'}> M




Self-Rep and Self-Rep                                                                         536

 (By the way, can you spot how this sentence differs from -G?) The reason I show it
explicitly is to point out that a Henkin sentence does not give a full recipe for its own
derivation; it just asserts that there exists one. You might well wonder whether its claim is
justified. Do Henkin sentences indeed possess derivations? Are they, as they claim,
theorems? It is useful to recall that one need not believe a politician who says. "I am honest"-
he may be honest, and yet he may not be. Are Henkin sentences any more trustworthy than
politicians? Or do Henkin sentences, like politicians, lie in cast-iron sinks?
         It turns out that these Henkin sentences are invariably truth tellers. Why this is so is
not obvious; but we will accept this curious fact without proof.

                           Implicit vs. Explicit Henkin Sentences

I mentioned that a Henkin sentence tells nothing about its own derivation; it just asserts that
one exists. Now it is possible to invent a variation on the theme of Henkin sentences-namely
sentences which explicitly describe their own derivations. Such a sentence's high-level
interpretation would not be "Some Sequence of Strings Exists Which is a Derivation of Me",
but rather, "The Herein-described Sequence of Strings           Is a Derivation of Me". Let us
call the first type of sentence an implicit Henkin sentence. The new sentences will be called
explicit Henkin sentences, since they explicitly describe their own derivations. Note that,
unlike their implicit brethren, explicit Henkin sentences need not be theorems. In fact, it is
quite easy to write a string which asserts that its own derivation consists of the single string
0=0-a false statement, since 0=0 is not a derivation of anything. However, it is also possible
to write an explicit Henkin sentence which is a theorem-that is, a sentence which in fact gives
a recipe for its own derivation.

                            Henkin Sentences and Self-Assembly

The reason I bring up this distinction between explicit and implicit Henkin sentences is that it
corresponds very nicely to a significant distinction between types of virus. There are certain
viruses, such as the so-called "tobacco mosaic virus", which are called self-assembling
viruses; and then there are others, such as our favorite T-evens, which are non-self-
assembling. Now what is this distinction? It is a direct analogue to the distinction between
implicit and explicit Henkin sentences.
        The DNA of a self-assembling virus codes only for the parts of a new virus, but not
for any enzymes. Once the parts are produced, the sneaky virus relies upon them to link up to
each other without help from any enzymes. Such a process depends on chemical affinities
which the parts have for each other, when swimming in the rich chemical brew of a cell. Not
only viruses, but also some organelles-such as ribosomes-assemble




Self-Rep and Self-Rep                                                                        537

themselves. Somtiems enzymes may be needed – but in such cases, they are recruited from
the host cell, and enslaved. This is what is meant by self-assembly.
         By contrast, the DNA of more complex viruses, such as the T-evens, codes not only
for the parts, but in addition for various enzymes which play special roles in the assembly of
the parts into wholes. Since the assembly process is not spontaneous but requires "machines",
such viruses are not considered to be self-assembling. The essence of the distinction, then,
between self-assembling units and non-self-assembling units is that the former get away with
self-reproduction without telling the cell anything about their construction, while the latter
need to give instructions as to how to assemble themselves.
         Now the parallel to Henkin sentences, implicit and explicit, ought to be quite clear.
Implicit Henkin sentences are self-proving but do not tell anything at all about their proofs-
they are analogous to self-assembling viruses; explicit Henkin sentences direct the
construction of their own proofs-they are analogous to more complex viruses which direct
their host cells in putting copies of themselves together.
         The concept of self-assembling biological structures as complex as viruses raises the
possibility of complex self-assembling machines as well. Imagine a set of parts which, when
placed in the proper supporting environment, spontaneously group themselves in such a way
as to form a complex machine. It seems unlikely, yet this is quite an accurate way to describe
the process of the tobacco mosaic virus' method of selfreproduction via self-assembly. The
information for the total conformation of the organism (or machine) is spread about in its
parts; it is not concentrated in some single place.
         Now this concept can lead in some strange directions, as was shown in the Edifying
Thoughts of a Tobacco Smoker. There, we saw how the Crab used the idea that information
for self-assembly can be distributed around, instead of being concentrated in a single place.
His hope was that this would prevent his new phonographs from succumbing to the Tortoise's
phonograph-crashing method. Unfortunately, just as with the most sophisticated axiom
schemata, once the system is all built and packaged into a box, its well-definedness renders it
vulnerable to a sufficiently clever "Godelizer"; and that was the sad tale related by the Crab.
Despite its apparent absurdity, the fantastic scenario of that Dialogue is not so far from
reality, in the strange, surreal world of the cell.


                               Two Outstanding Problems:
                            Differentiation and Morphogenesis

Now self-assembly may be the trick whereby certain subunits of cells are constructed, and
certain viruses-but what of the most complex macroscopic structures, such as the body of an
elephant or a spider, or the shape of a Venus's-Hyt-ap? How are homing instincts built into
the brain of" a




Self-Rep and Self-Rep                                                                      538

bird, or hunting instincts into the brain of a dog\% In short, how is it that merely by dictating
which proteins are to be produced in cells, DNA exercises such spectacularly precise control
over the exact structure and function of macroscopic living objects? There are two major
distinct problems here. One is that of cellular differentiation: how do different cells, sharing
exactly the same DNA, perform different roles-such as a kidney cell, a bone marrow cell, and
a brain cell? The other is that of morphogenesis ("birth of form"): how does intercellular
communication on a local level give rise to large-scale, global structures and organizations-
such as the various organs of the body, the shape of the face, the suborgans of the brain, and
so on? Although both cellular differentiation and rnorphogenesis are poorly understood at
present. the trick appears to reside in exquisitely fine-tuned feedback and "feedforward"
mechanisms within cells and between cells, which tell a cell when to "turn on" and when to
"turn off" production of various proteins.

                                Feedback and Feedforward

Feedback takes place when there is too much or too little of some desired substance in the
cell: then the cell must somehow regulate the production line which is assembling that
substance. Feedforward also involves the regulation of" an assembly line, but not according
to the amount of end product present: rather, according to the amount of some precursor of
the end product of that assembly line. There are two major devices for achieving negative
feedforward or feedback. One way is to prevent the relevant enzymes from being able to
perform-that is, to "clog up" their active sites. This is called inhibition. The other way is to
prevent the relevant enzymes from ever being manufactured! This is called repression.
Conceptually, inhibition is simple: you just block up the active site of the first enzyme in the
assembly line, and the whole process of synthesis gets stopped dead.

                                  Repressors and Inducers

Repression is trickier. How does a cell stop a gene from being expressed? The answer is, it
prevents it from ever getting transcribed. This means that it has to prevent RNA polymerase
from doing its job. This can be accomplished by placing a huge obstacle in its path, along the
DNA. precisely in front of that gene which the cell wants not to get transcribed. Such
obstacles do exist, and are called repressors. They are themselves proteins, and they bind to
special obstacle-holding sites on the DNA, called (I am not sure why) operators. An operator
therefore is a site of control for the gene (or genes) which immediately follow it: those genes
are called its operon. Because a series of enzymes often act in concert in carrying out a long
chemical transformation, they are often coded for in sequence; and this is why operons often
contain several genes, rather than just one. The effect of the successful repression of an
operon is that a whole series of genes is




Self-Rep and Self-Rep                                                                       539

prevented from being transcribed, which means that a whole set of related enzymes remains
unsynthesized.
         What about positive feedback and feedforward? Here again, there are two options: (1)
unclog the clogged enzymes, or (2) stop the repression of the relevant operon. (Notice how
nature seems to love double-negations! Probably there is some very deep reason for this.)
The mechanism by which repression is repressed involves a class of molecules called
inducers. The role of an inducer is simple: it combines with a repressor protein before the
latter has had a chance to bind to an operator on a DNA molecule; the resulting "repressor-
inducer complex" is incapable of binding to an operator, and this leaves the door open for the
associated operon to be transcribed into mRNA and subsequently translated into protein.
Often the end product or some precursor of the end product can act as an
inducer.

                         Feedback and Strange Loops Compared

Incidentally, this is a good time to distinguish between simple kinds of feedback, as in the
processes of inhibition and repression, and the looping-hack between different informational
levels, shown in the Central Dogrnap. Both are "feedback" in some sense; but the latter is
much deeper than the former. When an amino acid, such as tryptophan or isoleucine, acts as
feedback (in the form of an inducer) by binding to its repressor so that more of it gets made,
it is not telling how to construct itself; it is just telling enzymes to make more of it. This
could be compared to a radio's volume, which, when fed through a listener's ears, may cause
itself to be turned down or up. This is another thing entirely from the case in which the
broadcast itself tells you explicitly to turn your radio on or off, or to tune to another
wavelength-or even how to build another radio! The latter is much more like the looping-
back between informational levels, for here, information inside the radio signal gets
"decoded" and translated into mental structures. The radio signal is composed of symbolic
constituents whose symbolic meaning matters-a case of use, rather than mention. On the
other hand, when the sound is just too loud, the symbols are not conveying meaning: they are
merely being perceived as loud sounds, and might as well be devoid of meaning-a case of
mention, rather than use. This case more resembles the feedback loops by which proteins
regulate their own rates of synthesis.
         It has been theorized that the difference between two neighboring cells which share
the exact same genotype and yet have different functions is that different segments of their
genome have been repressed, and therefore they have different working sets of proteins. A
hypopothesis like this could account for the phenomenal differences between cells in
different organs of the body of a human being.




Self-Rep and Self-Rep                                                                     540

                          Two Simple Examples of Differentiation

The process by which one initial cell replicates over and over, giving rise to a myriad of
differentiated cells with specialized functions, can be likened to the spread of a chain letter
from person to person, in which each new participant is asked to propagate the message
faithfully, but also to add some extra personal touch. Eventually, there will be letters which
are tremendously different from each other.
         Another illustration of the ideas of differentiation is provided by this extremely
simple computer analogue of a differentiating self-rep. Consider a very short program which
is controlled by an up-down switch, and which has an internal parameter N-a natural number.
This program can run in two modes-the up-mode, and the down-mode. When it runs in the
upmode, it self-replicates into an adjacent part of the computer's memoryexcept it makes the
internal parameter N of its "daughter" one greater than in itself. When it runs in the down-
mode, it does not self-rep, but instead calculates the number

                                         (-1)'/(2N + 1)

and adds it to a running total.
         Well, suppose that at the beginning, there is one copy of the program in memory, N =
0, and the mode is up. Then the program will copy itself next door in memory, with N = 1.
Repeating the process, the new program will self-rep next door to itself, with a copy having
N = 2. And over and over again ... What happens is that a very large program is growing
inside memory. When memory is full, the process quits. Now all of memory can be looked
upon as being filled with one big program, composed of many similar, but differentiated,
modules-or "cells". Now suppose we switch the mode to down, and run this big program.
What happens? The first "cell" runs, and calculates 1/1. The second "cell" runs, calculating -
1/3, and adding it to the previous result. The third "cell" runs, calculating + 1/5 and adding it
on. .. The end result is that the whole "organism"-the big program-calculates the sum

                         l -1/3 +1/5 -117 +1/9 -1/11 +1/13 -1/15 + .. .

to a large number of terms (as many terms as "cells" can fit inside memory). And since this
series converges (albeit slowlv) to 7r/4, we have a "phenotype" whose function is to calculate
the value of a famous mathematical constant.

                                   Level Mixing in the Cell

I hope that the descriptions of processes such as labeling, self-assembly, differentiation,
morphogenesis, as well as transcription and translation, have helped to convey some notion
of the immensely complex system which is a cell-an information-processing system with
some strikingly




Self-Rep and Self-Rep                                                                        541

novel features. We have seen, in the Central Dogmap, that although we can try to draw a
clear line between program and data, the distinction is somewhat arbitrary. Carrying this line
of thought further, we find that not only are program and data intricately woven together, but
also the interpreter of programs, the physical processor, and even the language are included
in this intimate fusion. Therefore, although it is possible (to some extent) to draw boundaries
and separate out the levels, it is just as important-and ust as fascinating-to recognize the
level-crossings and mixings. Illustrative of this is the amazing fact that in biological systems,
all the various features necessary for self-rep (viz., language, program, data, interpreter, and
processor) cooperate to such a degree that all of them are replicated simultaneously-which
shows how much deeper is biological self-rep'ing than anything yet devised along those lines
by humans. For instance, the self-rep program exhibited at the beginning of this Chapter
takes for granted the pre-existence of three external aspects: a language, an interpreter, and a
processor, and does not replicate those.
         Let us try to summarize various ways in which the subunits of a cell can be classified
in computer science terms. First, let us take DNA. Since DNA contains all the information
for construction of proteins., which are the active agents of the cell, DNA can be viewed as a
program written in a higher-level language, which is subsequently translated (or interpreted)
into the "machine language" of the cell (proteins). On the other hand, DNA is itself a passive
molecule which undergoes manipulation at the hands of various kinds of enzymes; in this
sense, a DNA molecule is exactly like a long piece of data, as well. Thirdly, DNA contains
the templates off of which the tRNA "flashcards" are rubbed, which means that DNA also
contains the definition of its own higher-level language.
         Let us move on to proteins. Proteins are active molecules, and carry out all the
functions of the cell; therefore it is quite appropriate to think of them as programs in the
"machine language" of the cell (the cell itself being the processor). On the other hand, since
proteins are hardware and most programs are software, perhaps it is better to think of the
proteins as processors. Thirdly, proteins are often acted upon by other proteins, which means
that proteins are often data. Finally, one can view proteins as interpreters; this involves
viewing DNA as a collection of' high-level language programs, in which case enzymes are
merely carrying out the programs written in the DNA code, which is to say, the proteins are
acting as interpreters.
         Then there are ribosomes and tRNA molecules. They mediate the translation from
DNA to proteins, which can be compared to the translation of a program from a high-level
language to a machine language; in other words, the ribosomes are functioning as interpreters
and the tRNA molecules provide the definition of the higher-level language. But an
alternative view of translation has it that the ribosomes are processors, while the tRNA's are
interpreters.
         We have barely scratched the surface in this analysis of interrelations between all
these biomolecules. What we have seen is that nature feels quite




Self-Rep and Self-Rep                                                                        542

comfortable in mixing levels which we tend to see as quite distinct. Actually, in computer
science there is already a visible tendency to nix all these seemingly distinct aspects of an
information-processing system. This is particularly so in Artificial Intelligence research,
which is usually at the forefront of computer language design.

                                    The Origin of Life

A natural and fundamental question to ask, on learning of these incredibly intricately
interlocking pieces of software and hardware is: "How did they ever get started in the first
place?" It is truly a baffling thing. One has to imagine some sort of a bootstrap process
occurring, somewhat like that which is used in the development of new computer languages-
but a bootstrap from simple molecules to entire cells is almost beyond one's power to
imagine. There are various theories on the origin of life. They all run aground on this most
central of all central questions: "How did the Genetic Code, along with the mechanisms for
its translation (ribosomes and tRNA molecules), originate" For the moment, we will have to
content ourselves with a sense of wonder and awe, rather than with an answer. And perhaps
experiencing that sense of wonder and awe is more satisfying than having an answer-at least
for a while.




Self-Rep and Self-Rep                                                                    543

                         The Magnificrab, Indeed,
      It is spring, and the Tortoise and Achilles are taking a Sunday promenade in
      the woods together. They have decided to climb a hill at the top of which, it is
      said, there is a wonderful teahouse, with all sorts of delicious pastries.

Achilles: Man oh man! If a crab
Tortoise: If a crab??
Achilles: I was about to say, if a crab ever were intelligent, then surely it would be our
  mutual friend the Crab. Why, he must be at least two times as smart as any crab alive.
  Or maybe even three times as smart as any crab alive. Or perhaps
Tortoise: My soul! How you magnify the Crab!
Achilles: Well, I just happen to be an admirer of his ...
Tortoise: No need to apologize. I admire him, too. Speaking of Crab admirers, did I tell
  you about the curious fan letter which the Crab received not too long ago?
Achilles: I don't believe so. Who sent it?
Tortoise: It bore a postmark from India, and was from someone neither of us had ever
  heard of before-a Mr. Najunamar, I believe.
Achilles: I wonder why someone who never knew Mr. Crab would send him a letter-or
  for that matter, how they would get his address. Tortoise: Apparently whoever it was
  was under the illusion that the Crab is a mathematician. It contained numerous results,
  all of which were But, ho! Speak of the devil! Here comes Mr. Crab now, down the
  hill. Crab: Good-bye! It was nice to talk with you again. Well, I guess I had best be
  off. But I'm utterly stuffed-couldn't eat one more bite if I had to! I've just been up there
  myself-recommend it highly. Have you ever been to the teahouse at the crest of the
  hill? How are you, Achilles? Oh, there's Achilles. Hello, hello. Well, well, if it isn't
  Mr. T!
Tortoise: Hello, Mr. C. Are you headed up to the hilltop teahouse? Crab: Why, yes
  indeed, I am; how did you guess it? I'm quite looking forward to some of their special
  napoleons-scrumptious little morsels. I'm so hungry I could eat a frog. Oh, there's
  Achilles. How are you, Achilles?
Achilles: Could be worse, I suppose.
Crab: Wonderful! Well, don't let me interrupt your discussion. I'll just tag along.
Tortoise: Curiously enough, I was just about to describe your mysterious letter from that
  Indian fellow a few weeks back-but now that you're here. I'll let Achilles get the story
  from the Crab’s mouth.




The Magnificrab, Indeed,                                                                  544

            FIGURE 104. Castrovalva, by'M. C. Escher (lithograph, 1930).




The Magnificrab, Indeed,                                                   545

Crab: Well, it was this way. This fellow Najunamar had apparently never had any formal
  training in mathematics, but had instead worked out some of his own methods for
  deriving new truths of mathematics. Some of his discoveries defeated me completely;
  I had never seen anything in the least like them before. For instance, he exhibited a
  map of India that he had managed to color using no fewer than 1729 distinct colors.
Achilles: 1729! Did you say 1729? Crab: Yes-why do you ask?
Achilles: Well, 1729 is a very interesting number, you know. Crab: Indeed. I wasn't
  aware of it.
Achilles: In particular, it so happens that 1729 is the number of the taxicab which I took
  to Mr. Tortoise's this morning!
Crab: How fascinating! Could you possibly tell me the number of the trolley car which
  you'll take to Mr. Tortoise's tomorrow morning?
Achilles (after a moment's thought): It's not obvious to me; however, I should think it
  would be very large.
Tortoise: Achilles has a wonderful intuition for these things.
Crab: Yes. Well, as I was saying, Najunamar in his letter also proved that every even
  prime is the sum of two odd numbers, and that there are no solutions in positive
  integers to the equation

                                  a n + bn = c n   for n = 0.

Achilles: What? All these old classics of mathematics resolved in one fell swoop? He
  must be a genius of the first rank! Tortoise: But Achilles-aren't you even in the
  slightest skeptical?
Achilles: What? Oh, yes-skeptical. Well, of course I am. You don't think I believe that
  Mr. Crab got such a letter, do you? I don't fall for just anything, you know. So it must
  have been 5'ou, Mr. T, who received the letter!
Tortoise: Oh, no, Achilles, the part about Mr. C receiving the letter is quite true. What I
  meant was, aren't you skeptical about the content of the letter-its extravagant claims?
Achilles: Why should I be? Hmm ... Well, of course I am. I'm a very skeptical person, as
  both of you should well know by now. It's very hard to convince me of anything, no
  matter how true or false it is.
Tortoise: Very well put, Achilles. You certainly have a first-class awareness of your own
  mental workings.
Achilles: Did it ever occur to you, my friends, that these claims of Najunamar might be
  incorrect?
Crab: Frankly, Achilles, being rather conservative and orthodox myself, I was a bit
  concerned about that very point on first receiving the letter. In fact, I suspected at first
  that here was an out-and-out fraud. But on second thought, it occurred to me that not
  many types of people could manufacture such strange-sounding and complex results
  purely from their imagination. In fact, what it boiled down to was this question:




The Magnificrab, Indeed,                                                                  546

   "Which is the more likely: a charlatan of such extraordinary ingenuity, or a
   mathematician of great genius?" And before long, I realized that the probabilities clearly
   favored the former.
Achilles: Didn't you directly checkout any of his amazing claims, however?
Crab: Why should I? The probability argument was the most convincing thing I had ever
   thought of; no mathematical proof would have equaled it. But Mr. T here insisted on
   rigor. I finally gave in to his insistence, and checked all of Najunamar's results. To my
   great surprise, each one of them was right. How he discovered them, I'll never know,
   however. He must have some amazing and inscrutable Oriental type of insight which we
   here in the Occident can have no inkling of. At present, that's the only theory which
   makes an sense to me.
Tortoise: Mr. Crab has always been a little more susceptible to mystical or fanciful
   explanations than I am. I have full confidence that whatever Najunamar did in his way
   has a complete parallel inside orthodox mathematics. There is no way of doing
   mathematics which is fundamentally different from what we now know, in my opinion.
Achilles: That is an interesting opinion. I suppose it has something to do with the Church-
   Turing Thesis and related topics.
Crab: Oh, well, let us leave these technical matters aside on such a fine day, and enjoy the
   quiet of the forest, the chirping of the birds, and the play of sunlight on the new leaves
   and buds. Ho!
Tortoise: I second the motion. After all, all generations of Tortoises have reveled in such
   delights of nature.
Crab: As have all generations of Crabs.
Achilles: You don't happen to have brought your flute along, by any chance, Mr. C?
Crab: Why, certainly! I take it with me everywhere. Would you like to hear a tune or two?
Achilles: It would be delightful, in this pastoral setting. Do you play from memory?
Crab: Sad to say, that is beyond my capability. I have to read my music
from a sheet. But that is no problem. I have several very pleasant pieces here in this case.

  (He opens up a thin case and draws out a few pieces of paper. The topmost one has the
  following symbols on it:

                                          Va:-Sa=O

  He sticks the top sheet into a little holder attached to his flute, and plays. The tune is very
  short.)

Achilles: That was charming. (Peers over at the sheet on the flute, and a quizzical expression
  beclouds his face.) What is that statement of number theory doing, attached to your flute
  like that?

(The Crab looks at his flute, then his music, turns his head all around, and appears slightly
   confused.)




  The Magnificrab, Indeed,                                                                  547

Crab: I don't understand. What statement of number theory?
Achilles: "Zero is not the successor of any natural number." Right there, in the holder on
  your flute!
Crab: That's the third Piano Postulate. There are five of them, and I've arranged them all for
  flute. They're obvious, but catchy.
Achilles: What's not obvious to me is how a number-theoretical statement can be played as
  music.
Crab: But I insist, it's 'NOT a number-theoretical statement-it's a Piano Postulate! Would
  you like to hear another?
Achilles: I'd be enchanted.

(The Crab places another piece of paper on his flute, and this time Achilles watches more
   carefully.)

  Well, I watched your eyes, and they were looking at that FORMULA on the sheet. Are
  you sure that that is musical notation? I swear, it most amazingly resembles the notation
  which one might use in a formalized version of number theory.
Crab: How odd! But certainly that is music, not any kind of statement of mathematics, as
  far as I can tell! Of course, I am not a mathematician in any sense of the word. Would
  you like to hear any other tunes?
Achilles: By all means. Have you some others?
Crab: Scads.

   (He takes a new sheet, and attaches it to his flute. It contains the following symbols:

                            3a:3b:(SSa. SSb)=SSSSSSSSSSSSSO

   Achilles peers at it, while the Crab plays it.)

  Isn't it lovely?
Achilles: Yes, it certainly is a tuneful little piece. But I have to say, it's looking more and
  more like number theory to me.
Crab: Heavens! It is just my usual music notation, nothing more. I simply don't know how
  you read all these extramusical connotations into a straightforward representation for
  sounds.
Achilles: Would you be averse to playing a piece of my own composition?
Crab: Not in the least. Have you got it with you?
Achilles: Not yet, but I have a hunch I might be able to compose some tunes all by myself.
Tortoise: I must. tell you, Achilles, that Mr. C is a harsh judge of music composed by others,
  so do not be disappointed if, by some chance, he is not an enthusiast for your efforts.
Achilles: That is very kind of you to forewarn me. Still, I'm willing to give it a try .

   (He writes:

                 ((SSSO . SSSO) +(SSSSO. SSSSO))=(SSSSSO. SSSSSO)



   The Magnificrab, Indeed,                                                                  548

  The Crab takes it, looks it over for a moment, then sets it in his music holder, and pipes.)

Crab: Why, that's quite nice, Achilles. I enjoy strange rhythms.
Achilles: What's strange about the rhythms in that piece?
Crab: Oh, naturally, to you as the composer it must seem quite bland, but to my ears,
  shifting from a 3/3 rhythm to 4/4 and then to 5/5 is quite exotic. If you have any other
  songs, I'd be glad to play them. Achilles: Thank you very much. I've never composed
  anything before, and I must say composing is quite different from how I had imagined it
  to be. Let me try my hand at another one. (jots down a line.)

                         3a:3b:(SSa - SSb) =SSSSSSSSSSSSSSO

Crab: Hmmm ... Isn't that just a copy of my earlier piece?
Achilles: Oh, no! I've added one more S. Where you had thirteen in a row, I have fourteen.
Crab: Oh, yes. Of course. (He plays it, and looks very stern.)
Achilles: I do hope you didn't dislike my piece!
Crab: I am afraid, Achilles, that you completely failed to grasp the subtleties of my piece,
  upon which yours is modeled. But how could I expect you to understand it on first
  hearing? One does not always understand what is at the root of beauty. It is so easy to
  mistake the superficial aspects of a piece for its beauty, and to imitate them, when the
  beauty itself is locked deep inside the music, in a way which seems always to elude
  analysis.
Achilles: I am afraid that you have lost me a little in your erudite commentary. I understand
  that my piece does not measure up to your high standards, but I do not know exactly
  where I went astray. Could you perhaps tell me some specific way in which you find fault
  with my composition?
Crab: One possible way to save your composition, Achilles, would be to insert another three
  S's-five would do as well-into that long group of S's near the end. That would create a
  subtle and unusual effect.
Achilles: I see.
Crab: But there are other ways you might choose to change your piece. Personally, I would
  find it most appealing to put another tilde in the front. Then there would be a nice balance
  between the beginning and the end. Having two tildes in a row never fails to give a gay
  little twist to .a piece, you know.
Achilles: How about if I take both of your suggestions, and make the following piece?

                        -3a:3b:(SSa.SSb)=SSSSSSSSSSSSSSSSSO

Crab (a painful grimace crossing his face): Now, Achilles, it is important to learn the
  following lesson: never try to put too much into any single piece. There is always a point
  beyond which it cannot be improved,




  The Magnificrab, Indeed,                                                                 549

  and further attempts to improve it will in fact destroy it. Such is the case in this example.
  Your idea of incorporating both of my suggestions together does not yield the desired
  extra amount of beauty, but on the contrary creates an imbalance which quite takes away
  all the charm.
Achilles: How is it that two very similar pieces, such as yours with thirteen 5's, and mine
  with fourteen S's, seem to you to be so different in their musical worth? Other than in that
  minor respect, the two are identical.
Crab: Gracious! There is a world of difference between your piece and mine. Perhaps this is
  a place where words fail to convey what the spirit can feel. Indeed, I would venture to say
  that there exists no set of rules which delineate what it is that makes a piece beautiful, nor
  could there ever exist such a set of rules. The sense of Beauty is the exclusive domain of
  Conscious Minds, minds which through the experience of living have gained a depth that
  transcends explanation by any mere set of rules.
Achilles: I will always remember this vivid clarification of the nature of Beauty. I suppose
  that something similar applies to the concept of Truth, as well?
Crab: Without doubt. Truth and Beauty are as interrelated as-as
Achilles: As interrelated as, say, mathematics and music?
Crab: Oh! You took the words right out of my mouth! How did you know that that is what I
  was thinking?
Tortoise: Achilles is very clever, Mr. C. Never underestimate the potency of his insight.
Achilles: Would you say that there could conceivably be any relationship between the truth
  or falsity of a particular statement of mathematics, and the beauty, or lack of beauty, of
  an associated piece of music? Or is that just a far-fetched fancy of mine, with no basis in
  reality?
Crab: If you are asking me, that is carrying things much too far. When I spoke of the
  interrelatedness of music and mathematics, I was speaking very figuratively, you know.
  As for a direct connection between specific pieces of music and specific statements of
  mathematics, however, I harbor extremely grave doubts about its possibility. I would
  humbly counsel you not to give too' much time to such idle speculations.
Achilles: You are no doubt right. It would be most unprofitable. Perhaps I ought to
  concentrate on sharpening my musical sensitivity by composing some new pieces. Would
  you be willing to serve as my mentor,
Mr. C?
Crab: I would be very happy to aid you in your steps towards musical understanding.

       (So Achilles takes pen in hand, and, with what appears to be a great deal of
          concentration, writes:

                      AOOaV'\/--nn:b+cS(33=OAD((-d)<v(VS-+(>v




   The Magnificrab, Indeed,                                                                 550

   The Crab looks very startled.)

Y ou really want me to play that-that-that whatever-it-is?
Achilles: Oh, please do!

(So the Crab plays it, with evident difficulty.)

Tortoise: Bravo! Bravo! Is John Cage your favorite composer, Achilles? Achilles: Actually,
  he's my favorite anti-composer. Anyway, I'm glad you liked MY music.
Crab: The two of you may find it amusing to listen to such totally meaningless cacophony,
  but I assure you it is not at all pleasant for a sensitive composer to be subjected to such
  excruciating, empty dissonances and meaningless rhythms. Achilles, I thought you had a
  good feeling for music. Could it be that your previous pieces had merit merely by
  coincidence?
Achilles: Oh, please forgive me, Mr. Crab. I was trying to explore the limits of your musical
  notation. I wanted to learn directly what kinds of sound result when I write certain types
  of note sequences, and also how you evaluate pieces written in various styles.
Crab: Harrumph! I am not just an automatic music-machine, you know. Nor am I a garbage
  disposal for musical trash.
Achilles: I am very sorry. But I feel that I have learned a great deal by writing that small
  piece, and I am convinced that I can now write much better music than I ever could have
  if I hadn't tried that idea. And if you'll just play one more piece of mine, I have high
  hopes that you will feel better about my musical sensitivities.
Crab: Well, all right. Write it down and I'll give it a chance.

   (Achilles writes:

                       Ya:Vb:<(a -a) =(SSO -(b               > b))Da=0>

   and the Crab plays.)

  You were right, Achilles. You seem to have completely regained your musical acuity.
  This is a little gem! How did you come to compose it? I have never heard anything like it.
  It obeys all the rules of harmony, and yet has a certain-what shall I say?-irrational appeal
  to it. I can't put my finger on it, but I like it for that very reason.
Achilles: I kind of thought you might like it.
Tortoise: Have you got a name for it, Achilles? Perhaps you might call it "The Song of
  Pythagoras". You remember that Pythagoras and his followers were among the first to
  study musical sound.
Achilles: Yes, that's true. That would be a fine title.
Crab: Wasn't Pythagoras also the first to discover that the ratio of two squares can never be
  equal to 2? Tortoise: I believe you're right. It was considered a truly sinister discovery at
  the time, for never before had anyone realized that there are




   The Magnificrab, Indeed,                                                                551

  numbers-such as the square root of 2-which are not ratios of integers. And thus the
  discovery was deeply disturbing to the Pythagoreans, who felt that it revealed an
  unsuspected and grotesque defect in the abstract world of numbers. But I don't know
  what this has to do with the price of tea in China.
Achilles: Speaking of tea, isn't that the teahouse just up there ahead of us?
Tortoise: Yes, that's it, all right. We ought to be there in a couple of minutes.
Achilles: Hmm ... That's just enough time for me to whistle for you the tune which the taxi
  driver this morning had on his radio. It went like this.
Crab: Hold on for a moment; I'll get some paper from my case, and jot
  down your tune. (Scrounges around inside his case, and finds a blank sheet.)
  Go ahead; I'm ready.

   (Achilles whistles a rather long tune, and the Crab scrambles to keep up with him.)

  Could you whistle the last few bars again?
Achilles: Why, certainly.

   (After a couple of such repeats, the session is complete, and the Crab proudly displays
   his transcription:

   <((SSSSSO.SSSSSO)+(SSSSSO.SSSSSO))=((SSSSSSSO.SSSSSSSO)+(SO.SO))n---
            3b:<3c:(Sc+b)=((SSSSSSSO . SSSSSSSO)+(SO • SO))n3d:3d':3e:3e':
         <---<d=evd=e' >n<b=((Sd • Sd)+(Sd' • Sd'))nb=((Se • Se)+(Se' • Se'))>> >>

   The Crab then plays it himself.)

Tortoise: It's peculiar music, isn't it? It sounds a wee bit like music from India, to me.
Crab. Oh, I think it's too simple to be from India. But of course I know precious little about
  such things.
Tortoise: Well, here we are at the teahouse. Shall we sit outside here, on the verandah?
Crab: If you don't mind, I'd prefer to go inside. I've gotten perhaps enough sun for the day.

   (They go inside the teahouse and are seated at a nice wooden table, and order cakes and
   tea. Soon a cart of scrumptious-looking pastries is wheeled up, and each of them chooses
   his favorite.)

Achilles: You know, Mr. C, I would love to know what you think of another piece which I
  have just composed in my head. Crab: Can you show it to me? Here, write it down on this
  napkin.

(Achilles writes:

              da:3b:3c:<- 3d:3e:<(SSd•SSe)=bv(SSd•SSe)=c>n(a+a)=(b+c)>

The Crab and Tortoise study it with interest.)



   The Magnificrab, Indeed,                                                               552

Tortoise: Is it another beautiful piece, Mr. C, in your opinion?
Crab: Well, uh ... (Shifts in his chair, and looks somewhat uncomfortable.)
Achilles: What's the matter? Is it harder to decide whether this piece is beautiful than it is for
  other pieces?
Crab: Ahm ... No, it's not that-not at all. It's just that, well ... I really have to HEAR a piece
  before I can tell how much I like it.
Achilles: So go ahead and play it! I'm dying to know whether you find it beautiful or not.
Crab: Of course, I'd be extremely glad to play it for you. The only thing is
Achilles: Can't you play it for me? What's the matter? Why are you balking?
Tortoise: Don't you realize, Achilles, that for Mr. Crab to fulfill your request would be most
  impolite and disturbing to the clientele and employees of this fine establishment?
Crab (suddenly looking relieved): That's right. We have no right to impose our music on
  others.
Achilles (dejectedly): Oh, PHOOEY! And I so much wanted to know what he thinks of this
  piece!
Crab: Whew! That was a close call!
Achilles: What was that remark?
Crab: Oh-nothing. It's just that that waiter over there, he got knocked into by another waiter,
  and almost dropped a whole pot of tea into a lady's lap. A narrow escape, I must say.
  What do you say, Mr. Tortoise? Tortoise: Very good teas, I'd say. Wouldn't you agree,
  Achilles? Achilles: Oh, yes. Prime teas, in fact.
Crab: Definitely. Well, I don't know about you two, but I should perhaps be going, for I've a
  long steep trail back to my house, on the other side of this hill.
Achilles: You mean this is a big bluff?
Crab: You said it, Achilles.
Achilles: I see. Well, I'll have to remember that.
Crab: It has been such a jolly afternoon, Achilles, and I sincerely hope we will exchange
  more musical compositions another day.
Achilles: I'm looking forward to that very much, Mr. C. Well, good-bye. Tortoise: Good-
  bye, Mr. C.
(And the Crab heads off down his side of the hill.)
Achilles: Now there goes a brilliant fellow ... In my estimation, he's at least four times as
  smart as any crab alive. Or he might even be five
Tortoise: As you said in the beginning, and probably shall be saying forevermore, words
  without end.




   The Magnificrab, Indeed,                                                                   553


                Church, Turing, Tarski, and Others
                          Formal and Informal Systems
WE HAVE COME to the point where we can develop one of the main theses of this
book: that every aspect of thinking can be viewed as a high-level description of a system
which, on a low level, is governed by simple, even formal, rules. The "system", of course,
is a brain-unless one is' speaking of thought processes flowing in another medium, such
as a computer's circuits. The image is that of a formal system underlying an "informal
system"-a system which can, for instance, make puns, discover number patterns, forget
names, make awful blunders in chess, and so forth. This is what one sees from the
outside: its informal, overt, software level. By contrast, it has a formal, hidden, hardware
level (or "substrate") which is a formidably complex mechanism that makes transitions
from state to state according to definite rules physically embodied in it, and according to
the input of signals which impinge on it.
        A vision of the brain such as this has many philosophical and other consequences,
needless to say. I shall try to spell some of them out in this Chapter. Among other things,
this vision seems to imply that, at bottom, the brain is some sort of a "mathematical"
object. Actually, that is at best a very awkward way to look at the brain. The reason is
that, even if a brain is, in a technical and abstract sense, some sort of formal system, it
remains true that mathematicians only work with simple and elegant systems, systems in
which everything is extremely clearly defined-and the brain is a far cry from that, with its
ten billion or more semi-independent neurons, quasi-randomly connected up to each
other. So mathematicians would never study a real brain's networks. And if you define
"mathematics" as what mathematicians enjoy doing, then the properties of brains are not
mathematical.
        The only way to understand such a complex system as a brain is by chunking it on
higher and higher levels, and thereby losing some precision at each step. What emerges at
the top level is the "informal system" which obeys so many rules of such complexity that
we do not yet have the vocabulary to think about it. And that is what Artificial
Intelligence research is hoping to find. It has quite a different 'flavor from mathematics
research. Nevertheless, there is a loose connection to mathematics: Al people often come
from a strong mathematics background, and




Church, Turing, Tarski, and Others                                                      554

mathematicians sometimes are intrigued by the workings of their own brains. The
following passage, quoted from Stanislaw Ulam's autobiographical Adventures of a
Mathematician, illustrates this point:

   It seems to me that more could be done to elicit ... the nature of associations, with
   computers providing the means for experimentation. Such a study would have to
   involve a gradation of notions, of symbols, of classes of symbols, of classes of
   classes, and so on, in the same way that the complexity of mathematical or
   physical structures is investigated.
        There must be a trick to the train of thought, a recursive formula. A group of
   neurons starts working automatically, sometimes without external impulse. It is a
   kind of iterative process with a growing pattern. It wanders about in the brain, and
   the way it happens must depend on the memory of similar patterns.'

                       Intuition and the Magnificent Crab
Artificial Intelligence is often referred to as "Al". Often, when I try to explain what is
meant by the term, I say that the letters "AI" could just as well stand for "Artificial
Intuition", or even "Artificial Imagery". The aim of Al is to get at what is happening
when one's mind silently and invisibly chooses, from a myriad alternatives, which one
makes most sense in a very complex situation. In many real-life situations, deductive
reasoning is inappropriate, not because it would give wrong answers, but because there
are too many correct but irrelevant statements which can be made; there are just too many
things to take into account simultaneously for reasoning alone to be sufficient. Consider
this mini-dialogue:

   "The other day I read in the paper that the--
   "Oh-you were reading? It follows that you have eyes. Or at least one eye. Or
   rather, that you had at least one eye then."

A sense of judgment-"What is important here, and what is not?"-is called for. Tied up
with this is a sense of simplicity, a sense of beauty. Where do these intuitions come from?
How can they emerge from an underlying formal system?
        In the Magnificrab, some unusual powers of the Crab's mind are revealed. His
own version of his powers is merely that he listens to music and distinguishes the
beautiful from the non-beautiful. (Apparently for him there is a sharp dividing line.) Now
Achilles finds another way to describe the Crab's abilities: the Crab divides statements of
number theory into the categories true and false. But the Crab maintains that, if he
chances to do so, it is only by the purest accident, for he is, by his own admission,
incompetent in mathematics. What makes the Crab's performance all the more mystifying
to Achilles, however, is that it seems to be in direct violation of a celebrated result of
metamathematics with which Achilles is familiar:

CHURCH's THEOREM: There is no infallible method for telling theorems of TNT from
nontheorems.


Church, Turing, Tarski, and Others                                                     555

It was proven in 1936 by the American logician Alonzo Church. Closely related is what I
call the

TARSKI-CHURCH-TURING THEOREM: There is no infallible method for telling true
from false statements of number theory.

                             The Church-Turing Thesis
To understand Church's Theorem and the Tarski-Church-Turing Theorem better, we
should first describe one of the ideas on which they are based; and that is the Church-
Turing Thesis (often called "Church's Thesis"). For the Church-Turing Thesis is certainly
one of the most important concepts in the philosophy of mathematics, brains, and
thinking.
        Actually, like tea, the Church-Turing Thesis can be given in a variety of different
strengths. So I will present it in various versions, and we will consider what they imply.
The first version sounds very innocent-in fact almost pointless:

       CHURCH-TURING THESIS, TAUTOLOGICAL VERSION: Mathematics
       problems can be solved only by doing mathematics.

Of course, its meaning resides in the meaning of its constituent terms. By "mathematics
problem" I mean the problem of deciding whether some number possesses or does not
possess a given arithmetical property. It turns out that by means of Godel-numbering and
related coding tricks, almost any problem in any branch of mathematics can be put into
this form, so that "mathematics problem" retains its ordinary meaning. What about "doing
mathematics"? When one tries to ascertain whether a number has a property, there seem
to be only a small number of operations which one uses in combination over and over
again-addition, multiplication, checking for equality or inequality. That is, loops
composed of such operations seem to be the only tool we have that allows us to probe the
world of numbers. Note the word "seem". This is the critical word which the Church-
Turing Thesis is about. We can give a revision:

   CHURCH-TURING THESIS, STANDARD VERSION: Suppose there is a
    method which a sentient being follows in order to sort numbers into two classes.
    Suppose further that this method always yields an answer within a finite amount
    of time, and that it always gives the same answer for a given number. Then:
    Some terminating FlooP program (i.e., some general recursive function) exists
    which gives exactly the same answers as the sentient being's method does.

The central hypothesis, to make it very clear, is that any mental process which divides
numbers into two sorts can be described in the form of a FlooP program. The intuitive
belief is that there are no other tools than those in FlooP, and that there are no ways to use
those tools other than by




Church, Turing, Tarski, and Others                                                        556

unlimited iterations (which FlooP allows). The Church-Turing Thesis is not a provable
fact in the sense of a Theorem of mathematics-it is a hypothesis about the processes
which human brains use.

                           The Public-Processes Version
Now some people might feel that this version asserts too much. These people might put
their objections as follows: "Someone such as the Crab might exist-someone with an
almost mystical insight into mathematics, but who is just as much in the dark about his
own peculiar abilities as anyone else-and perhaps that person's mental mechanisms carry
out operations which have no counterpart in FlooP." The idea is that perhaps we have a
subconscious potential for doing things which transcend the conscious processes-things
which are somehow inexpressible in terms of the elementary FlooP operations. For these
objectors, we shall give a weaker version of the Thesis, one which distinguishes between
public and private mental processes:

   CHURCH-TURING THESIS, PUBLIC-PROCESSES VERSION: Suppose there
    is a method which a sentient being follows in order to sort numbers into two
    classes. Suppose further that this method always yields an answer within a finite
    amount of time, and that it always gives the same answer for a given number.
    Proviso. Suppose also that this method can be communicated reliably from one
    sentient being to another by means of language. Then: Some terminating FlooP
    program (i.e., general recursive function) exists which gives exactly the same
    answers as the sentient beings' method does.

This says that public methods are subject to "FlooPification", but asserts nothing about
private methods. It does not say that they are un-FlooP-able, but it at least leaves the door
open.

                                Srinivasa Ramanujan
As evidence against any stronger version of the Church-Turing Thesis, let us consider the
case of the famous Indian mathematician of the first quarter of the twentieth century,
Srinivasa Ramanujan (1887-1920). Ramanujan (Fig. 105) came from Tamilnadu, the
southernmost part of India, and studied mathematics a little in high school. One day,
someone who recognized Ramanujan's talent for math presented him with a copy of a
slightly out-of-date textbook on analysis, which Ramanujan devoured (figuratively
speaking). He then began making his own forays into the world of analysis, and by the
time he was twenty-three, he had made a number of discoveries which he considered
worthwhile. He did not know to whom to turn, but somehow was told about a professor
of mathematics in faraway England, named G. H. Hardy. Ramanujan compiled his best




Church, Turing, Tarski, and Others                                                       557

                                                 FIGURE 105. Srinivasa Ramanujan and
                                                 one of his strange Indian melodies.




results together in a packet of papers, and sent them all to the, unforewarned Hardy with a
covering letter which friends helped him express
in English. Below are some excerpts taken from Hardy's description of his reaction upon
receiving the bundle:

   ... It soon became obvious that Ramanujan must possess much more general
   theorems and was keeping a great deal up his sleeve.... [Some formulae]defeated
   me completely; I had never seen anything in the least like them before. A single
   look at them is enough to show that they could only be written down by a
   mathematician of the highest class. They must be true because, if they were not
   true, no one would have had the imagination to invent them. Finally ... the writer
   must be completely honest, because great mathematicians are commoner than
   thieves or humbugs of such incredible skill2

What resulted from this correspondence was that Ramanujan came to England in 1913,
sponsored by Hardy; and then followed an intense collaboration which terminated in
Ramanujan's early demise at age thin thirty-three from tuberculosis.
        Ramanujan had several extraordinary characteristics which set him apart from the
majority of mathematicians. One was his lack of rigor. Very often he would simply state
a result which, he would insist, had just come to




Church, Turing, Tarski, and Others                                                     558

him from a vague intuitive source, far out of the realm of conscious probing. In fact, he
often said that the goddess Namagiri inspired him in his dreams. This happened time and
again, and what made it all the more mystifying-perhaps even imbuing it with a certain
mystical quality-was the fact that many of his "intuition-theorems" were wrong. Now
there is a curious paradoxical effect where sometimes an event which you think could not
help but make credulous people become a little more skeptical, actually has the reverse
effect, hitting the credulous ones in some vulnerable spot of their minds, tantalizing them
with the hint of some baffling irrational side of human nature. Such was the case with
Ramanujan's blunders: many educated people with a yearning to believe in something of
the sort considered Ramanujan's intuitive powers to be evidence of a mystical insight into
Truth, and the fact of his fallibility seemed, if anything, to strengthen, rather than
weaken, such beliefs.
        Of course it didn't hurt that he was from one of the most backward parts of India,
where fakirism and other eerie Indian rites had been practiced for millennia, and were
still practiced with a frequency probably exceeding that of the teaching of higher
mathematics. And his occasional wrong flashes of insight, instead of suggesting to people
that he was merely human, paradoxically inspired the idea that Ramanujan's wrongness
always had some sort of "deeper rightness" to it-an "Oriental" rightness, perhaps touching
upon truths inaccessible to Western minds. What a delicious, almost irresistible thought!
Even Hardy-who would have been the first to deny that Ramanujan had any mystical
powers-once wrote about one of Ramanujan's failures, "And yet I am not sure that, in
some ways, his failure was not more wonderful than any of his triumphs."
        The other outstanding feature of Ramanujan's mathematical personality was his
"friendship with the integers", as his colleague Littlewood put it. This is a characteristic
that a fair number of mathematicians share to some degree or other, but which
Ramanujan possessed to an extreme. There are a couple of anecdotes which illustrate this
special power. The first one is related by Hardy:

   I remember once going to see him when he was lying ill at Putney. I had ridden in
   taxi-cab No. 1729, and remarked that the number seemed to me rather a dull one,
   and that I hoped it was not an unfavorable omen. "No," he replied, "it is a very
   interesting number; it is the smallest number expressible as a sum of two cubes in
   two different ways." .1 asked him, naturally, whether he knew the answer to the
   corresponding problem for fourth powers; and he replied, after a moment's
   thought, that he could see no obvious example, and thought that the first such
   number must be very large.'

It turns out that the answer for fourth powers is:

                          635318657 = 134 + 1334 = 1584 + 594

The reader may find it interesting to tackle the analogous problem for squares, which is
much easier.
       It is actually quite interesting to ponder why it is that Hardy im-




Church, Turing, Tarski, and Others                                                      560

mediately jumped to fourth powers. After all, there are several other reasonably natural
generalizations of the equation

                                         u3 + v3 = x3 + y3

along different dimensions. For instance, there is the question about representing a
number in three distinct ways as a sum of two cubes:

                                  r3 + 3 = u3 + v3 = x3 + y3

Or, one can use three different cubes:
                                 u3
                                      + v3 + w3 = x3 + y3 + z3,

Or one can even make a Grand Generalization in all dimensions at once:

                          r4 + s4 + t4 = u4 + v4 + w4 = x4 + v4 + z4

There is a sense, however, in which Hardy's generalization is "the most mathematician-
like". Could this sense of mathematical esthetics ever be programmed?
        The other anecdote is taken from a biography of Ramanujan by his
countryman S. R. Ranganathan, where it is called "Ramanujan's Flash". It is related by a
Indian friend of Ramanujan's from his Cambridge days, Dr. P. C. Mahalanobis.

   On another occasion, I went to his room to have lunch with him. The First World
   War had started some time earlier. I had in my hand a copy of the monthly "Strand
   Magazine" which at that time used to publish a number of puzzles to be solved by
   readers. Ramanujan was stirring something in a pan over the fire for our lunch. I
   was sitting near the table, turning over the pages of the Magazine. I got interested in
   a problem involving a relation between two numbers. I have forgotten the details;
   but I remember the type of the problem. Two British officers had been billeted in
   Paris in two different houses in a long street; the door numbers of these houses
   were related in a special way; the problem was to find out the two numbers. It was
   not at all difficult. I got the solution in a few minutes by trial and error.

   MAHALANOBIS (in a joking way): Now here is a problem for you.
   RAMANUJAN: What problem, tell me.'(He went on stirring the pan.) I read out the
   question from the "Strand Magazine".
   RAMANUJAN: Please take down the solution. (He dictated a continued fraction.)
        The first term was the solution which I had obtained. Each successive term
   represented successive solutions for the same type of relation between two numbers,
   as the number of houses in the street would increase indefinitely. I was amazed.
   MAHALANOBIS: Did you get the solution in a flash?
   RAMANUJAN: Immediately I heard the problem, it was clear that the Solution was
   obviously a continued fraction; I then thought, "Which continued fraction?" and the
   answer came to my mind. It was just as simple as this.'

Hardy, as Ramanujan's closest co-worker, was often asked after

Church, Turing, Tarski, and Others                                                       561

Ramanujan's death if there had been any occult or otherwise exotically flavored elements
to Ramanujan's thinking style. Here is one comment which he gave:

  I have often been asked whether Ramanujan had any special secret; whether his
  methods differed in kind from those of other mathematicians; whether there was
  anything really abnormal in his mode of thought. I cannot answer these questions
  with any confidence or conviction; but I do not believe it. My belief is that all
  mathematicians think, at bottom, in the same kind of way, and that Ramanujan was
  no exception .5

Here Hardy states in essence his own version of the Church-Turing Thesis. I paraphrase:
CHURCH-TURING THESIS, HARDY'S VERSION: At bottom, all mathematicians are
isomorphic.

This does not equate the mathematical potential of mathematicians with that of general
recursive functions; for that, however, all you need is to show that some mathematician's
mental capacity is no more general than recursive functions. Then, if you believe Hardy's
Version, you know it for all mathematicians.
Then Hardy compares Ramanujan with calculating prodigies:

  His memory, and his powers of calculation, were very unusual, but they could not
  reasonably be called "abnormal". If he had to multiply two large numbers, he
  multiplied them in the ordinary way; he could do it with unusual rapidity and
  accuracy, but not more rapidly and accurately than any mathematician who is
  naturally quick and has the habit of computations

Hardy describes what he perceived as Ramanujan's outstanding intellectual attributes:

  With his memory, his patience, and his power of calculation, he combined a power
  of generalisation, a feeling for form, and a capacity for rapid modification of his
  hypotheses, that were often really startling, and made him, in his own field, without a
  rival in his day.'

The part of this passage which I have italicized seems to me to be an excellent
characterization of some of the subtlest features of intelligence in general. Finally, Hardy
concludes somewhat nostalgically:

  [His work has not the simplicity and inevitableness of the very greatest work; it
  would be greater if it were less strange. One gift it has which no one can deny-
  profound and invincible originality. He would probably have been a greater
  mathematician if he had been caught and tamed a little in his youth; he would have
  discovered more that was new, and that, no doubt, of greater importance. On the
  other hand he would have been less of a Ramanujan, and more of a European
  professor and the loss might have been greater than the gain.8

The esteem in which Hardy held Ramanujan is revealed by the romantic way in which he
speaks of him.


Church, Turing, Tarski, and Others                                                      562

                                    "Idiots 'Savants"
There is another class of people whose mathematical abilities seem to defy rational
explanation-the so-called "idiots savants", who can perform complex calculations at
lightning speeds in their heads (or wherever they do it). Johann Martin Zacharias Dase,
who lived from 1824 to 1861 and was employed by various European governments to
perform computations, is an outstanding example. He not only could multiply two
numbers each of 100 digits in his head; he also had an uncanny sense of quantity. That is,
he could just "tell", without counting, how many sheep were in a field, or words in a
sentence, and so forth, up to about 30-this in contrast to most of us, who have such a
sense up to about 6, with reliability. Incidentally, Dase was not an idiot.
        I shall not describe the many fascinating documented cases of "lightning
calculators", for that is not my purpose here. But I do feel it is important to dispel the idea
that they do it by some mysterious, unanalyzable method. Although it is often the case
that such wizards' calculational abilities far exceed their abilities to explain their results,
every once in a while, a person with other intellectual gifts comes along who also has this
spectacular ability with numbers. From such people's introspection, as well as from
extensive research by psychologists, it has been ascertained that nothing occult takes
place during the performances of lightning calculators, but simply that their minds race
through intermediate steps with the kind of self-confidence that a natural athlete has in
executing a complicated motion quickly and gracefully. They do not reach their answers
by some sort of instantaneous flash of enlightenment (though subjectively it may feel that
way to some of them), but-like the rest of us-by sequential calculation, which is to say, by
FlooP-ing (or BlooP-ing) along. -
        Incidentally, one of the most obvious clues that no "hot line to God" is involved is
the mere fact that when the numbers involved get bigger, the answers are slower in
coming. Presumably, if God or an "oracle" were supplying the answers, he wouldn't have
to slow up when the numbers got bigger. One could probably make a nice plot showing
how the time taken by a lightning calculator varies with the sizes of the numbers
involved, and the operations involved, and from it deduce some features of the algorithms
employed.

           The Isomorphism Version of the Church-Turing Thesis
This finally brings us to a strengthened standard version of the Church-Turing Thesis:

   CHURCH-TURING THESIS, ISOMORPHISM VERSION: Suppose there is a
    method which a sentient being follows in order to sort numbers into two classes.
    Suppose further that this method always yields an answer within a finite amount
    of time, and that it always gives the same answer for a given number. Then:
    Some terminating FlooP program (i.e.,




Church, Turing, Tarski, and Others                                                         563

     general recursive function) exists which gives exactly the same answers as the
     sentient being's method does. Moreover: The mental process and the FlooP
     program are isomorphic in the sense that on some level there is a correspondence
     between the steps being carried out in both computer and brain.

Notice that not only has the conclusion been strengthened, but also the proviso of
communicability of the faint-hearted Public-Processes Version has been dropped. This
bold version is the one which we now shall discuss.
        In brief, this version asserts that when one computes something, one's mental
activity can be mirrored isomorphically in some FlooP program. And let it be very clear
that this does not mean that the brain is actually running a FlooP program, written in the
FlooP language complete with BEGIN's, END's, ABORT's, and the rest-not at all. It is
just that the steps are taken in the same order as they could be in a FlooP program, and
the logical structure of the calculation can be mirrored in a FlooP program.
        Now in order to make sense of this idea, we shall have to make some level
distinctions in both computer and brain, for otherwise it could be misinterpreted as utter
nonsense. Presumably the steps of the calculation going on inside a person's head are on
the highest level, and are supported by lower levels, and eventually by hardware. So if we
speak of an isomorphism, it means we've tacitly made the assumption that the highest
level can be isolated, allowing us to discuss what goes on there independently of other
levels, and then to map that top level into FlooP. To be more precise, the assumption is
that there exist software entities which play the roles of various mathematical constructs,
and which are activated in ways which can be mirrored exactly inside FlooP (see Fig.
106). What enables these software entities to come into existence is the entire
infrastructure discussed in Chapters XI and XI I, as well as in the Prelude, Ant Fugue.
There is no assertion of isomorphic activity on the lower levels of brain and computer
(e.g., neurons and bits).
        The spirit of the Isomorphism Version, if not the letter, is gotten across by saying
that what an idiot savant does in calculating, say, the logarithm of 7r, is isomorphic to
what a pocket calculator does in calculating it-where the isomorphism holds on the
arithmetic-step level, not on the lower levels of, in the one case, neurons, and in the other,
integrated circuits. (Of course different routes can be followed in calculating anything-but
presumably the pocket calculator, if not the human, could be instructed to calculate the
answer in any specific manner.)

FIGURE 106. The behavior of natural numbers can be mirrored in a human brain or in
the programs of a computer. These two different representations can then be mapped
onto each other on an appropriately abstract level.




Church, Turing, Tarski, and Others                                                        564

             Representation of Knowledge about the Real World
Now this seems quite plausible when the domain referred to is number theory, for there
the total universe in which things happen is very small and clean. Its boundaries and
residents and rules are well-defined, as in a hard-edged maze. Such a world is far less
complicated than the open-ended and ill-defined world which we inhabit. A number
theory problem, once stated, is complete in and of itself. A real-world problem, on the
other hand, never is sealed off from any part of the world with absolute certainty. For
instance, the task of replacing a burnt-out light bulb may turn out to require moving a
garbage bag; this may unexpectedly cause the spilling of a box of pills, which then forces
the floor to be swept so that the pet dog won't eat any of the spilled pills, etc., etc. The
pills and the garbage and the dog and the light bulb are all quite distantly related parts of
the world-yet an intimate connection is created by some everyday happenings. And there
is no telling what else could be brought in by some other small variations on the
expected. By contrast, if you are given a number theory problem, you never wind up
having to consider extraneous things such as pills or dogs or bags of garbage or brooms
in order to solve your problem. (Of course, your intuitive knowledge of such objects may
serve you in good stead as you go about unconsciously trying to manufacture mental
images to help you in visualizing the problem in geometrical terms-but that is another
matter.)
        Because of the complexity of the world, it is hard to imagine a little pocket
calculator that can answer questions put to it when you press a few buttons bearing labels
such as "dog", "garbage", "light bulb", and so forth. In fact, so far it has proven to be
extremely complicated to have a full-size high-speed computer answer questions about
what appear to us to be rather simple subdomains of the real world. It seems that a large
amount of knowledge has to be taken into account in a highly integrated way for
"understanding" to take place. We can liken real-world thought processes to a tree whose
visible part stands sturdily above ground but depends vitally on its invisible roots which
extend way below ground, giving it stability and nourishment. In this case the roots
symbolize complex processes which take place below the conscious level of the mind-
processes whose effects permeate the way we think but of which we are unaware. These
are the "triggering patterns of symbols" which were discussed in Chapters XI and XII.
        Real-world thinking is quite different from what happens when we do a
multiplication of two numbers, where everything is "above ground", so to speak, open to
inspection. In arithmetic, the top level can be "skimmed off " and implemented equally
well in many different sorts of hardware: mechanical adding machines, pocket
calculators, large computers, people's brains, and so forth. This is what the Church-
Turing Thesis is all about. But when it comes to real-world understanding, it seems that
there is no simple way to skim off the top level, and program it. alone. The triggering
patterns of symbols are just too complex. There must he several levels through which
thoughts may "percolate" and "bubble".




Church, Turing, Tarski, and Others                                                       565

        In particular-and this comes back to a major theme of Chapters XI ' and XII-the
representation of the real world in the brain, although rooted in isomorphism to some
extent, involves some elements which have no counterparts at all in the outer world. That
is, there is much more to it than simple mental structures representing "dog", "broom",
etc. All of these symbols exist, to be sure-but their internal structures are extremely
complex and to a large degree are unavailable for conscious inspection. Moreover, one
would hunt in vain to map each aspect of a symbol's internal structure onto some specific
feature of the real world.

                     Processes That Are Not So Skimmable
For this reason, the brain begins to look like a very peculiar formal system, for on its
bottom level-the neural level-where the "rules" operate and change the state, there may be
no interpretation of the primitive elements (neural firings, or perhaps even lower-level
events). Yet on the top level, there emerges a meaningful interpretation-a mapping from
the large "clouds" of neural activity which we have been calling "symbols", onto the real
world. There is some resemblance to the Gödel construction, in that a high-level
isomorphism allows a high level of meaning to be read into strings; but in the Gödel
construction, the higher-level meaning "rides" on the lower level-that is, it is derived
from the lower level, once the notion of Gödel-numbering has been introduced. But in the
brain, the events on the neural level are not subject to real-world interpretation; they are
simply not imitating anything. They are there purely as the substrate to support the higher
level, much as transistors in a pocket calculator are there purely to support its number-
mirroring activity. And the implication is that there is no way to skim off just the highest
level and make an isomorphic copy in a program; if one is to mirror the brain processes
which allow real-world understanding, then one must mirror some of the lower-level
things which are taking place: the "languages of the brain". This doesn't necessarily mean
that one must go all the way down to the level of the hardware, though that may turn out
to be the case.
         In the course of developing a program with the aim of achieving an "intelligent"
(viz., human-like) internal representation of what is "out there", at some point one will
probably be forced into using structures and processes which do not admit of any
straightforward interpretations-that is, which cannot be directly mapped onto elements of
reality. These lower layers of the program will be able to be understood only by virtue of
their catalytic relation to layers above them, rather than because of some direct
connection they have to the outer world. (A concrete image of this idea was suggested by
the Anteater in the Ant Fugue: the "indescribably boring nightmare" of trying to
understand a book on the letter level.)
         Personally, I would guess that such multilevel architecture of concept-handling
systems becomes necessary just when processes involving images and analogies become
significant elements of the program-in




Church, Turing, Tarski, and Others                                                      566

contrast to processes which are supposed to carry out strictly deductive reasoning.
Processes which carry out deductive reasoning can be programmed in essentially one
single level, and are therefore skimmable, by definition. According to my hypothesis,
then, imagery and analogical thought processes intrinsically require several layers of
substrate and are therefore intrinsically non-skimmable. I believe furthermore that it is
precisely at this same point that creativity starts to emerge-which would imply that
creativity intrinsically depends upon certain kinds of "uninterpretable" lower-level
events. The layers of underpinning of analogical thinking are, of course, of extreme
interest, and. some speculations on their nature will be offered in the next two Chapters.

                         Articles of Reductionistic Faith
One way to think about the relation between higher and lower levels in the brain is this.
One could assemble a neural net which, on a local (neuron-to-neuron) level, performed in
a manner indistinguishable from a neural net in a brain, but which had no higher-level
meaning at all. The fact that the lower level is composed of interacting neurons does not
necessarily force any higher level of meaning to appear-no more than the fact that
alphabet soup contains letters forces meaningful sentences to be found, swimming about
in the bowl. High-level meaning is an optional feature of a neural network-one which
may emerge as a consequence of evolutionary environmental pressures.
       Figure 107 is a diagram illustrating the fact that emergence of a higher level of
meaning is optional. The upwards-pointing arrow indicates that a substrate can occur
without a higher level of meaning, but not vice versa: the higher level must be derived
from properties of a lower one.

FIGURE 107. Floating on neural activity, the symbol level of the brain mirrors the
world. But neural activity per se, which can be simulated on a computer, does not create
thought; that calls for higher levels of organization.




Church, Turing, Tarski, and Others                                                    567

The diagram includes an indication of a computer simulation of a neural network. This is
in principle feasible, no matter how complicated the network, provided that the behavior
of individual neurons can be described in terms of computations which a computer can
carry out. This is a subtle postulate which few people even think of questioning.
Nevertheless it is a piece of "reductionistic faith"; it could be considered a "microscopic
version" of the Church-Turing Thesis. Below we state it explicitly:

   CHURCH-TURING THESIS, MICROSCOPIC VERSION: The behavior of the
     components of a living being can be simulated on a computer. That is, the
     behavior of any component (typically assumed to be a cell) can be calculated by
     a FlooP program (i.e., general recursive function) to any desired degree of
     accuracy, given a sufficiently precise description of the component's internal
     state and local environment.

This version of the Church-Turing Thesis says that brain processes do not possess any
more mystique-even though they possess more levels of organization-than, say, stomach
processes. It would be unthinkable in this day and age to suggest that people digest their
food, not by ordinary chemical processes, but by a sort of mysterious and magic
"assimilation". This version of the CT-Thesis simply extends this kind of commonsense
reasoning to brain processes. In short, it amounts to faith that the brain operates in a way
which is, in principle, understandable. It is a piece of reductionist faith.
       A corollary to the Microscopic CT-Thesis is this rather terse new macroscopic
version:

CHURCH-TURING THESIS, REDUCTIONIST'S VERSION: All brain processes are
  derived from a computable substrate.

This statement is about the strongest theoretical underpinning one could give in support
of the eventual possibility of realizing Artificial Intelligence.
        Of course, Artificial Intelligence research is not aimed at simulating neural
networks, for it is based on another kind of faith: that probably there are significant
features of intelligence which can be floated on top of entirely different sorts of substrates
than those of organic brains. Figure 108 shows the presumed relations among Artificial
Intelligence, natural intelligence, and the real world.

                 Parallel Progress in Al and Brain Simulation?
The idea that, if Al is to be achieved, the actual hardware of the brain might one day have
to be simulated or duplicated, is, for the present at least, quite an abhorrent thought to
many Al workers. Still one wonders, "How finely will we need to copy the brain to
achieve Al?" The real answer is probably that it all depends on how many of the features
of human consciousness you want to simulate.




Church, Turing, Tarski, and Others                                                        568

FIGURE 108. Crucial to the endeavor of Artificial Intelligence research is the notion that
the symbolic levels of the mind can be "skimmed off " of their neural substrate and
implemented in other media, such as the electronic substrate of computers. To what depth
the copying of brain must go is at present completely unclear.

         Is an ability to play checkers well a sufficient indicator of intelligence? If so, then
Al already exists, since checker-playing programs are of world class. Or is intelligence an
ability to integrate functions symbolically, as in a freshman calculus class? If so, then AI
already exists, since symbolic integration routines outdo the best people in most cases. Or
is intelligence the ability to play chess well? If so, then AI is well on its way, since chess-
playing programs can defeat most good amateurs; and the level of artificial chess will
probably continue to improve slowly.
         Historically, people have been naive about what qualities, if mechanized, would
undeniably constitute intelligence. Sometimes it seems as though each new step towards
Al, rather than producing something which everyone agrees is real intelligence, merely
reveals what real intelligence is not. If intelligence involves learning, creativity,
emotional responses, a sense of beauty, a sense of self, then there is a long road ahead,
and it may be that these will only be realized when we have totally duplicated a living
brain.

                           Beauty, the Crab, and the Soul
Now what, if anything, does all this have to say about the Crab's virtuoso performance in
front of Achilles? There are two issues clouded together here. They are:




Church, Turing, Tarski, and Others                                                         569

   (1) Could any brain process, under any circumstances, distinguish completely
     reliably between true and false statements of TNT without being in violation of
     the Church-Turing Thesis-or is such an act in principle impossible?
   (2) Is perception of beauty a brain process?

First of all, in response to (1), if violations of the Church-Turing Thesis are allowed, then
there seems to be no fundamental obstacle to the strange events in the Dialogue. So what
we are interested in is whether a believer in the Church-Turing Thesis would have to
disbelieve in the Crab's ability. Well, it all depends on which version of the CT-Thesis
you believe. For example, if you only subscribe to the Public-Processes Version, then you
could reconcile the Crab's behavior with it very easily by positing that the Crab's ability is
not communicable. Contrariwise, if you believe the Reductionist's Version, you will have
a very hard time believing in the Crab's ostensible ability (because of Church's Theorem-
soon to be demonstrated). Believing in intermediate versions allows you a certain amount
of wishy-washiness on the issue. Of course, switching your stand according to
convenience allows you to waffle even more.
        It seems appropriate to present a new version of the CT-Thesis, one which is
tacitly held by vast numbers of people, and which has been publicly put forth by several
authors, in various manners. Some of the more famous ones are: philosophers Hubert
Dreyfus, S. Jaki, Mortimer Taube, and J. R. Lucas; the biologist and philosopher Michael
Polanyi (a holist par excellence); the distinguished Australian neurophysiologist John
Eccles. I am sure there are many other authors who have expressed similar ideas, and
countless readers who are sympathetic. I have attempted below to summarize their joint
position. I have probably not done full justice to it, but I have tried to convey the flavor as
accurately as I can:

   CHURCH-TURING THESIS, SOULISTS' VERSION: Some kinds of things which
     a brain can do can be vaguely approximated on a computer but not most, and
     certainly not the interesting ones. But anyway, even if they all could, that would
     still leave the soul to explain, and there is no way that computers have any
     bearing on that.

This version relates to the tale of the Magnificrab in two ways. In the first place, its
adherents would probably consider the tale to be silly and implausible, but-not forbidden
in principle. In the second place, they would probably claim that appreciation of qualities
such as beauty is one of those properties associated with the elusive soul, and is therefore
inherently possible only for humans, not for mere machines.
        We will come back to this second point in a moment; but first, while we are on
the subject of "soulists", we ought to exhibit this latest version in an even more extreme
form, since that is the form to which large numbers of well-educated people subscribe
these days:

CHURCH-TURING THESIS, THEODORE ROSZAK VERSION: Computers are
    ridiculous. So is science in general.



Church, Turing, Tarski, and Others                                                         570

This view is prevalent among certain people who see in anything smacking of numbers or
exactitude a threat to human values. It is too bad that they do not appreciate the depth and
complexity and beauty involved in exploring abstract structures such as the human mind,
where, indeed, one comes in intimate contact with the ultimate questions of what to be
human is.
        Getting back to beauty, we were about to consider whether the appreciation of
beauty is a brain process, and if so, whether it is imitable by a computer. Those who
believe that it is not accounted for by the brain are very unlikely to believe that a
computer could possess it. Those who believe it is a brain process again divide up
according to which version of the CT-Thesis they believe. A total reductionist would
believe that any brain process can in principle be transformed into a computer program;
others, however, might feel that beauty is too ill-defined a notion for a computer program
ever to assimilate. Perhaps they feel that the appreciation of beauty requires an element of
irrationality, and therefore is incompatible with the very fiber of computers.

          Irrational and Rational Can Coexist on Different Levels
         However, this notion that "irrationality is incompatible with computers" rests on a
         severe confusion of levels. The mistaken notion stems from the idea that since
         computers are faultlessly functioning machines, they are therefore bound to be
         "logical" on all levels. Yet it is perfectly obvious that a computer can be
         instructed to print out a sequence of illogical statements-or, for variety's sake, a
         batch of statements having random truth values. Yet in following such
         instructions, a computer would not be making any mistakes! On the contrary, it
         would only be a mistake if the computer printed out something other than the
         statements it had been instructed to print. This illustrates how faultless
         functioning on one level may underlie symbol manipulation on a higher level-and
         the goals of the higher level may be completely unrelated to the propagation of
         Truth.
Another way to gain perspective on this is to remember that a brain, too, is a collection of
faultlessly functioning elements-neurons. Whenever a neuron's threshold is surpassed by
the sum of the incoming signals, BANG!-it fires. It never happens that a neuron forgets
its arithmetical knowledge-carelessly adding its inputs and getting a wrong answer. Even
when a neuron dies, it continues to function correctly, in the sense that its components
continue to obey the laws of mathematics and physics. Yet as we all know, neurons are
perfectly capable of supporting high-level behavior that is wrong, on its own level, in the
most amazing ways. Figure 109 is meant to illustrate such a clash of levels: an incorrect
belief held in the software of a mind, supported by the hardware of a faultlessly
functioning brain.
         The point-a point which has been made several times earlier in various contexts-is
simply that meaning can exist on two or more different levels of a symbol-handling
system, and along with meaning, rightness and wrongness can exist on all those levels.
The presence of meaning on a given




Church, Turing, Tarski, and Others                                                       571

FIGURE 109. The brain is rational; the mind may not be. [Drawing by the author.]




Church, Turing, Tarski, and Others                                                 572

level is determined by whether or- not reality is mirrored in an isomorphic (or looser)
fashion on that level. So the fact that neurons always perform correct additions (in fact,
much more complex calculations) has no bearing whatsoever on the correctness of the
top-level conclusions supported by their machinery. Whether one's top level is engaged in
proving koans of Boolean Buddhism or in meditating on theorems of Zer1 Algebra, one's
neurons are functioning rationally. By the same token, the high-level symbolic processes
which in a brain create the experience of appreciating beauty are perfectly rational on the
bottom level, where the faultless functioning is taking place; any irrationality, if there is
such, is on the higher level, and is an epiphenomenon-a consequence-of the events on the
lower level.
        To make the same point in a different way, let us say you are having a hard time
making up your mind whether to order a cheeseburger or a pineappleburger. Does this
imply that your neurons are also balking, having difficulty deciding whether or not to
fire? Of course not. Your hamburger-confusion is a high-level state which fully depends
on the efficient firing of thousands of neurons in very organized ways. This is a little
ironic, yet it is perfectly obvious when you think about it. Nevertheless, it is probably fair
to say that nearly all confusions about minds and computers have their origin in just such
elementary level-confusions.
        There is no reason to believe that a. computer's faultlessly functioning hardware
could not support high-level symbolic behavior which would represent such complex
states as confusion, forgetting, or appreciation of beauty. It would require that there exist
massive subsystems interacting with each other according to a complex "logic". The
overt behavior could appear either rational or irrational; but underneath it would be the
performance of reliable, logical hardware.

                                 More Against Lucas
Incidentally, this kind of level distinction provides us with some new fuel in arguing
against Lucas. The Lucas argument is based on the idea that Gödel’s Theorem is
applicable, by definition, to machines. In fact, Lucas makes a most emphatic
pronunciation:

   Gödel’s theorem must apply to cybernetical machines, because it is of the essence
   of being a machine, that it should be a concrete instantiation of a formal system.°

This is, as we have seen, true on the hardware level-but since there may be higher levels,
it is not the last word on the subject. Now Lucas gives the impression that in the mind-
imitating machines he discusses, there is only one level on which manipulation of
symbols takes place. For instance, the Rule of Detachment (called "Modus Ponens" in his
article) would be wired into the hardware and would be an unchangeable feature of such
a machine. He goes further and intimates that if Modus Ponens were not an




Church, Turing, Tarski, and Others                                                        573

immutable pillar of the machine's system, but could be overridden on occasion, then:

   The system will have ceased to be a formal logical system. and the machine will
   barely qualify for the title of a model for the mind.10

Now many programs which are being developed in At research have very little in
common with programs for generating truths of number theory--programs with inflexible
rules of inference and fixed sets of axioms. Yet they are certainly intended as "models for
the mind". On their top level the "informal" level-there may be manipulation of images,
formulation of analogies, forgetting of ideas, confusing of concepts, blurring of
distinctions, and so forth. But this does not contradict the fact that they rely on the correct
functioning of their underlying hardware as much as brains rely on the correct
functioning of their neurons. So At programs are still "concrete instantiations of formal
systems"-but they are not machines to which Lucas' transmogrification of Gödel’s proof
can be applied. Lucas' argument applies merely to their bottom level, on which their
intelligence-however great or small it may be-does not lie.
        There is one other way in which Lucas betrays his oversimplified vision of how
mental processes would have to be represented inside computer programs. In discussing
the matter of consistency, he writes

   If we really were inconsistent machines, we should remain content with our
   inconsistencies, and would happily affirm both halves of a contradiction. Moreover,
   we would be prepared to say absolutely anything-which we are not. It is easily
   shown that in an inconsistent formal system everything is provable."

This last sentence shows that Lucas assumes that the Propositional Calculus must of
necessity be built into any formal system which carries out reasoning. In particular, he is
thinking of the theorem <<PA-P>DQ> of the Propositional Calculus; evidently he has
the erroneous belief that it is an inevitable feature of mechanized reasoning. However, it
is perfectly plausible that logical thought processes, such as propositional reasoning, will
emerge as consequences of the general intelligence of an At program, rather than being
preprogrammed. This is what happens in humans! And there is no particular reason to
assume that the strict Propositional Calculus, with its rigid rules and the rather silly
definition of consistency that they entail, would emerge from such a program.

                                An Underpinning of Al
We can summarize this excursion into level distinctions and come away with one final,
strongest version of the Church-Turing Thesis:

CHURCH-TURING THESIS, At VERSION: Mental processes of any sort can be
simulated by a computer program whose underlying language is of




Church, Turing, Tarski, and Others                                                         574

power equal to that of FlooP-that is, in which all partial recursive functions can be
programmed.

        It should also be pointed out that in practice, many AI researchers rely on another
article of faith which is closely related to the CT-Thesis, and which I call the AI Thesis. It
runs something like this:

   AI THESIS:       As the intelligence of machines evolves, its underlying
      mechanisms will gradually converge to the mechanisms underlying human
      intelligence.

In other words, all intelligences are just variations on a single theme; to create true
intelligence, At workers will just have to keep pushing to ever lower levels, closer and
closer to brain mechanisms, if they wish their machines to attain the capabilities which
we have.

                                  Church's Theorem
Now let us come back to the Crab and to the question of whether his decision procedure
for theoremhood (which is presented in the guise of a filter for musical beauty) is
compatible with reality. Actually, from the events which occur in the Dialogue, we have
no way of deducing whether the Crab's gift is an ability to tell theorems from
nontheorems, or alternatively, an ability to tell true statements from false ones. Of course
in many cases this amounts to the same thing but Gödel’s Theorem shows that it doesn't
always. But no matter: both of these alternatives are impossible, if you believe the At
Version of the Church-Turing Thesis. The proposition that it is impossible to have a
decision procedure for theoremhood in any formal system with the power of TNT is
known as Church's Theorem. The proposition that it is impossible to have a decision
procedure for number theoretical truth-if such truth exists, which one can well doubt after
meeting up with all the bifurcations of TNT-follows quickly from Tarski's Theorem
(published in 1933, although the ideas were known to Tarski considerably earlier).
        The proofs of these two highly important results of metamathematics are very
similar. Both of them follow quite quickly from self-referential constructions. Let us first
consider the question of a decision procedure for TNT-theoremhood. If there were a
uniform way by which people could decide which of the classes "theorem" and
"nontheorem" any given formula X fell into, then, by the CT-Thesis (Standard Version),
there would exist a terminating FlooP program (a general recursive function) which could
make the same decision, when given as input the Gödel number of formula X. The
crucial step is to recall that any property that can be tested for by a terminating FlooP
program is represented in TNT. This means that the property of TNT-theoremhood would
be represented (as distinguished from merely expressed) inside TNT. But as we shall see
in a moment, this,




Church, Turing, Tarski, and Others                                                        575

would put us in hot water, for if theoremhood is a representable attribute, then Gödel’s
formula G becomes as vicious as the Epimenides paradox.
        It all hinges on what G says: "G is not a theorem of TNT". Assume that G were a
theorem. Then, since theoremhood is supposedly represented, the TNT-formula which
asserts "G is a theorem" would be a theorem of TNT. But this formula is -G, the negation
of G, so that TNT is inconsistent. On the other hand, assume G were not a theorem. Then
once again by the supposed representability of theoremhood, the formula which asserts
"G is not a theorem" would be a theorem of TNT. But this formula is G, and once again
we get into paradox. Unlike the situation before, there is no resolution of the paradox.
The problem is created by the assumption that theorernhood is represented by some
formula of TNT, and therefore we must backtrack and erase that assumption. This forces
us also to conclude that no FlooP program can tell the Gödel numbers of theorems from
those of nontheorems. Finally, if we accept the Al Version of the CT-Thesis, then we
must backtrack further, and conclude that no method whatsoever could exist by which
humans could reliably tell theorems from nontheorems-and this includes determinations
based on beauty. Those who subscribe only to the Public-Processes Version might still
think the Crab's performance is possible; but of all the versions, that one is perhaps the
hardest one to find any justification for.

                                 Tarski's Theorem
Now let us proceed to Tarski's result. Tarski asked whether there could be a way of
expressing in TNT the concept of number-theoretical truth. That theoremhood is
expressible (though not representable) we have seen; Tarski was interested in the
analogous question regarding the notion of truth. More specifically, he wished to
determine whether there is any TNT-formula with a single free variable a which can be
translated thus:

               "The formula whose Gödel number is a expresses a truth."

Let us suppose, with Tarski, that there is one-which we'll abbreviate as TRUE{a}. Now
what we'll do is use the diagonalization method to produce a sentence which asserts about
itself that it is untrue. We copy the Godel method exactly, beginning with an "uncle":

                       3a:<-TRUE{a}nARITHMOQUINE{a",a}>

Let us say the Gödel number of the uncle is t. We arithmoquine this very uncle, and
produce the Tarski formula T:

               3a:<--TRUE{a}AARITHMOQUINE{SSS ... SSSO/a",a}>
                                             t S's




Church, Turing, Tarski, and Others                                                    576

When interpreted, it says:

                             "The arithmoquinification of t is the
                             Gödel number of a false statement."

But since the arithmoquinification of t is T's own Gödel number, Tarski's formula T
reproduces the Epimenides paradox to a tee inside TNT, saying of itself, "I am a falsity".
Of course, this leads to the conclusion that it must be simultaneously true and false (or
simultaneously neither). There arises now an interesting matter: What is so bad about
reproducing the Epimenides paradox? Is it of any consequence? After all, we already
have it in English, and the English language has not gone up in smoke.

                      The Impossibility of the Magnificrab
The answer lies in remembering that there are two levels of meaning involved here. One
level is the level we have just been using; the other is as a statement of number theory. If
the Tarski formula T actually existed, then it would be a statement about natural numbers
that is both true and false at once! There is the rub. While we can always just sweep the
English-language Epimenides paradox under the rug, saying that its subect matter (its
own truth) is abstract, this is' not so when it becomes a concrete statement about
numbers! If we believe this is a ridiculous state of affairs, then we have to undo our
assumption that the formula TRUE{a} exists. Thus, there is no way of expressing the
notion of truth inside TNT. Notice that this makes truth a far more elusive property than
theoremhood, for the latter is expressible. The same backtracking reasons as before
(involving the Church-Turing Thesis, Al Version) lead us to the conclusion that

      The Crab's mind cannot be a truth-recognizer any more than it is a TNT-
      theorem-recognizer.

The former would violate the Tarski-Church-Turing Theorem ("There is no decision
procedure for arithmetical truth"), while the latter would violate Church's Theorem.

                                  Two Types of Form
It is extremely interesting, then, to think about the meaning of the word "form" as it
applies to constructions of arbitrarily complex shapes. For instance, what is it that we
respond to when we look at a painting and feel its beauty? Is it the "form" of the lines and
dots on our retina? Evidently it must be, for that is how it gets passed along to the
analyzing mechanisms in our heads-but the complexity of the processing makes us feel
that we are not merely looking at a two-dimensional surface; we are responding to




Church, Turing, Tarski, and Others                                                      577

some sort of inner meaning inside the picture, a multidimensional aspect trapped
somehow inside those two dimensions. It is the word "meaning" which is important here.
Our minds contain interpreters which accept
two-dimensional patterns and then "pull" from them high-dimensional notions which are
so complex that we cannot consciously describe them. The same can be said about how
we respond to music, incidentally.
        It feels subjectively that the pulling-out mechanism of inner meaning is not at all
akin to a decision procedure which checks for the presence or absence of some particular
quality such as well-formedness in a string. Probably this is because inner meaning is
something which reveals more of itself over a period of time. One can never be sure, as
one can about well-formedness, that one has finished with the issue.
        This suggests a distinction that could be drawn between two senses of "form" in
patterns which we analyze. First, there are qualities such as well-formedness, which can
be detected by predictably terminating tests, as in BlooP programs. These I propose to
call syntactic qualities of form. One intuitively feels about the syntactic aspects of form
that they lie close to the surface, and therefore they do not provoke the creation of
multidimensional cognitive structures.
        By contrast, the semantic aspects of form are those which cannot be tested for in
predictable lengths of time: they require open-ended tests. Such an aspect is theoremhood
of TNT-strings, as we have seen. You cannot just apply some standard test to a string and
find out if it is a theorem. Somehow, the fact that its meaning is involved is crucially
related to the difficulty of telling whether or not a string is a TNT-theorem. The act of
pulling out a string's meaning involves, in essence, establishing all the implications of its
connections to all other strings, and this leads, to be sure, down an open-ended trail. So
"semantic" properties are connected to open-ended searches because, in an important
sense, an object's meaning is not localized within the object itself. This is not to say that
no understanding of any object's meaning is possible until the end of time, for as time
passes, more and more of the meaning unfolds. However, there are always aspects of its
meaning which will remain hidden arbitrarily long.

             Meaning Derives from Connections to Cognitive Structures

       Let us switch from strings to pieces of music, just for variety. You may still
       substitute the term "string" for every reference to a piece of music, if you prefer.
       The discussion is meant to be general, but its flavor is better gotten across, I feel,
       by referring to music. There is a strange duality about the meaning of a piece of
       music: on the one hand, it seems to be spread around, by virtue of its relation to
       many other things in the world-and yet, on the other hand, the meaning of a piece
       of music is obviously derived from the music itself, so it must be localized
       somewhere inside the music.
       The resolution of this dilemma comes from thinking about the interpreter-the
       mechanism which does the pulling-out of meaning. (By "inter




Church, Turing, Tarski, and Others                                                       578

preter in this context, I mean not -the performer of the piece, but the mental mechanism
in the listener which derives meaning when the piece is played.) The interpreter may
discover many important aspects of a piece's meaning while hearing it for the first time;
this seems to confirm the notion that the meaning is housed in the piece itself, and is
simply being read off. But that is only part of the story. The music interpreter works by
setting up a multidimensional cognitive structure-a mental representation of *\_he piece-
which it tries to integrate with pre-existent information by finding links to other
multidimensional mental structures which encode previous experiences. As this process
takes place, the full meaning gradually unfolds. In fact, years may pass before someone
comes to feel that he has penetrated to the core meaning of a piece. This seems to support
the opposite view: that musical meaning is spread around, the interpreter's role being to
assemble it gradually.
         The truth undoubtedly lies somewhere in between: meanings-both musical and
linguistic-are to some extent localizable, to some extent spread around. In the
terminology of Chapter VI, we can say that musical pieces and pieces of text are partly
triggers, and partly carriers of explicit meaning. A vivid illustration of this dualism of
meaning is provided by the example of a tablet with an ancient inscription: the meaning
is partially stored in the libraries and the brains of scholars around the world, and yet it is
also obviously implicit in the tablet itself.
         Thus, another way of characterizing the difference between "syntactic" and
"semantic" properties (in the just-proposed sense) is that the syntactic ones reside
unambiguously inside the object under consideration, whereas semantic properties
depend on its relations with a potentially infinite class of other objects, and therefore are
not completely localizable. There is nothing cryptic or hidden, in principle, in syntactic
properties, whereas hiddenness is of the essence in semantic properties. That is the reason
for my suggested distinction between "syntactic" and "semantic" aspects of visual form.

                              Beauty, Truth, and Form
What about beauty? It is certainly not a syntactic property, according to the ideas above.
Is it even a semantic property? Is beauty a property which, for instance, a particular
painting has? Let us immediately restrict our consideration to a single viewer. Everyone
has had the experience of finding something beautiful at one time, dull another time-and
probably intermediate at other times. So is beauty an attribute which varies in time? One
could turn things around and say that it is the beholder who has varied in time. Given a
particular beholder of a particular painting at a particular time, is it reasonable to assert
that beauty is a quality that is definitely present or absent? Or is there still something ill-
defined and intangible about it?
        Different levels of interpreter probably could be invoked in every




Church, Turing, Tarski, and Others                                                         579

person, depending on the circumstances. These various interpreters pull out different
meanings, establish different connections, and generally evaluate all deep aspects
differently. So it seems that this notion of beauty is extremely hard to pin down. It is for
this reason that I chose to link beauty, in the Magnificrab, with truth, which we have seen
is also one of the most intangible notions in all of metamathematics.

               The Neural Substrate of the Epimenides Paradox
I would like to conclude this Chapter with some ideas about that central problem of truth,
the Epimenides paradox. I think the Tarski reproduction of the Epimenides paradox
inside TNT points the way to a deeper understanding of the nature of the Epimenides
paradox in English. What Tarski found was that his version of the paradox has two
distinct levels to it. On one level, it is a sentence about itself which would be true if it
were false, and false if it were true. On the other level-which I like to call the arithmetical
substrate-it is a sentence about integers which is true if and only if false.
        Now for some reason this latter bothers people a lot more than the former. Some
people simply shrug off the former as "meaningless", because of its self-referentiality.
But you can't shrug off paradoxical statements about integers. Statements about integers
simply cannot be both true and false.
        Now my feeling is that the Tarski transformation of the Epimenides paradox
teaches us to look for a substrate in the English-language version. In the arithmetical
version, the upper level of meaning is supported by the lower arithmetical level. Perhaps
analogously, the self-referential sentence which we perceive ("This sentence is false") is
only the top level of a dual-level entity. What would be the lower level, then? Well, what
is the mechanism that language rides on? The brain. Therefore one ought to look for a
neural substrate to the Epimenides paradox-a lower level of physical events which clash
with each other. That is, two events which by their nature cannot occur simultaneously. If
this physical substrate exists, then the reason we cannot make heads or tails of the
Epimenides sentence is that our brains are trying to do an impossible task.
        Now what would be the nature of the conflicting physical events? Presumably
when you hear the Epimenides sentence, your brain sets up some "coding" of the
sentence-an internal configuration of interacting symbols. Then it tries to classify the
sentence as "true" or "false". This classifying act must involve an attempt to force several
symbols to interact in a particular way. (Presumably this happens when any sentence is
processed.) Now if it happens that the act of classification would physically disrupt the
coding of the sentence-something which would ordinarily never happen-then one is in
trouble, for it is tantamount to trying to force a record player to play its self-breaking
record. We have described the conflict in physical terms, but not in neural terms. If this
analysis is right so




Church, Turing, Tarski, and Others                                                        580

far, then presumably the rest of the discussion could be carried on when we know •
something about the constitution of the "symbols" in the brain out of neurons and their
firings, as well as about the way that sentences become converted into "codings".
         This sketch of the neural substrate of the Epimenides paradox suggests (to me, at
least) that the resolution of the English version of the Epimenides paradox might be
similar to that for the Tarski version. The resolution involves abandoning the notion that
a brain could ever provide a fully accurate representation for the notion of truth. The
novelty of this resolution lies in its suggestion that a total modeling of truth is impossible
for quite physical reasons: namely, such a modeling would require physically
incompatible events to occur in a brain.




Church, Turing, Tarski, and Others                                                        581

                                    SHRDLU,
                             Toy of Man's Designing
         One day, Eta Oin wanders into the MIT Artificial Intelligence Laboratory,
         where she runs into the brilliant young computer program, SHRDLU. It
         happens that SHRDLU is just dying for someone to try out the newly
         developed human being "Dr. Tony Earrwig". SHRDLU explains that this
         Earrwig is fairly intelligent in the limited domain of analyzing a conversation
         about a "toy world" containing blocks of various shapes, sizes, and colors-
         blocks which sit on a table, and which can be picked up and moved around.
         Eta Oin, fascinated, begins enthusiastically typing to SHRDLU. Earrwig
         dangles over her shoulder and provides a sort of running commentary, as
         promised.

1. Eta Oin: pick up a big red block. [Refer to Fig. 110.
   SHRDLU: OK.
        Dr. Tony Earrwig: SHRDLU answers "OK" when it carries out a command. In order
          to pick up the red block, it had to clear it off by finding a space for the green one
          and moving the green one away.
2. Eta Oin: grasp the pyramid.
   SHRDLU: I DON'T UNDERSTAND WHICH PYRAMID YOU MEAN.
   Dr. Tony Earrwig: There are three pyramids in the scene, and SHRDLU knows that
   phrases beginning with "the" are intended to refer to a specific object the speaker has in
   mind.

FIGURE 110. "Pick up a big red block." [From Terry Winograd, Understanding Natural
  Language, (New York: Academic Press, 1972), p. 8.1




   SHRDLU, Toy of Man's Designing                                                          582

FIGURE 111. "Find a block which is bigger than the one you are holding and put it into the
  box." Adapted from Terry Winograd, Understanding Natural Language, p.9.]

3. Eta Oin: find a block which is taller than the one you are holding and put it into the box.
   [Refer to Fig. 111.]
   Dr. Tony Earrwig: SHRDLU has to figure out what is meant by "one" and "it", as well as
   using deduction to decide which block it is holding.
   SHRDLU: BY "IT", I ASSUME YOU MEAN THE BLOCK WHICH IS TALLER
   THAN THE ONE I AM HOLDING.
   Dr. Tony Earrwig: There were two possible referents for "it": the one selected, and "the
   one you are holding". SHRDLU has heuristics for choosing between them.
   SHRDLU: OK.
   Dr. Tony Earrwig: SHRDLU's hand has to set down the red block pick up the taller blue
   one.
Eta Oin: what does the box contain?
   SHRDLU: THE BLUE PYRAMID AND THE BLUE BLOCK.
Eta Oin: what is the pyramid supported by?
   Dr. Tony Earrwig: Earlier SHRDLU responded to "the pyramid" by saying "I DON'T
   UNDERSTAND WHICH PYRAMID YOU MEAN." Now, it assumes that you mean the
   one just mentioned.
   SHRDLU: THE BOX.
Eta Oin: how many blocks are not in the box?




  SHRDLU, Toy of Man's Designing                                                          583

Dr. Tony Earrwig: Negatives such as "not" and "no" are handled whether they appear in
   clauses or noun groups (such as "no boxes").
        SHRDLU: FOUR OF THEM.
7. Eta Oin: is at least one of them narrower than the one which i told you to pick up?
        Dr. Tony Earrwig: The phrase "at least one of them" must be analyzed to refer
        back to the previous question, and "the one which ..." must also refer to blocks.
        Noun groups can have complicated numbers such as "at least one", "more than
        three", "exactly seven", etc.. The memory of past events is used to deduce which
        block is "the one which I told you to pick up". Notice the complicated syntactic
        structure in which "the one" is actually the object of a clause embedded two levels
        deeper. SHRDLU's parser incorporates a grammar that handles this kind of
        complexity in a general way.
        SHRDLU: YES, THE RED CUBE.
8. Eta Oin: is it supported?
        SHRDLU: YES, BY THE TABLE.
        Dr. Tony Earrwig: SHRDLU has heuristics about what information is relevant to
        a person who asks a question, even if the person does not request that information
        specifically.
9. Eta Oin: can the table pick up blocks?
        SHRDLU: NO,
        Dr. Tony Earrwig: SHRDLU's semantic rules include the knowledge that "picking
        up" must be done by an animate object, and that the table is not animate.
10. Eta Oin: can a pyramid be supported by a block?
        SHRDLU: YES.
        Dr. Tony Earrwig:• SHRDLU's deductive system finds an actual example, so it
        knows this is possible.
11. Eta Oin: can a pyramid support a pyramid?
        SHRDLU: I DON'T KNOW.
        Dr. Tony Earrwig: The question makes sense, but there is no example in the
        scene. SHRDLU has no way of looking into the details of its programs, even
        though these ultimately define its capabilities.
12. Eta Oin: stack up two pyramids.
        SHRDLU: I CANT.
        Dr. Tony Earrwig: SHRDLU tries and fails.
13. Eta Oin: the blue pyramid is mine.
        SHRDLU: I UNDERSTAND.
        Dr. Tony Earrwig: All of the information about colors, shapes, sizes and locations
        of the objects is known to SHRDLU initially. Possession is left as an unknown
        property of objects about which it can be told. It stores this information in its data
        base.
14. Eta Gin: i own all blocks which are not red, but i don't own anything which supports a
        pyramid.
SHRDLU: I UNDERSTAND.




SHRDLU, Toy of Man's Designing                                                            584

FIGURE 112. "Will you please stack up both of the red blocks and either a green cube or
a pyramid?" [Adapted from Terry Winograd, Understanding Natural Language, p. 12.

        Dr. Tony Earrwig: The result of understanding these sentences is the creation of
        new procedures within SHRDLU for deducing whether you own something.
15. Eta Oin: do i own the box?
        SHRDLU: NO.
        Dr. Tony Earrwig: The box supports whatever it contains, including a pyramid.
16. Eta Oin: do i own anything in the box?
        SHRDLU: YES, TWO THINGS: THE BLUE BLOCK AND PYRAMID.
17. Eta Oin: will you please stack up both of the red blocks and either a green cube or a
        pyramid? [Refer to Fig. 112.
        SHRDLU: OK.
        Dr. Tony Earrwig: Logical connectives, such as "and", "or", "either", etc. are
           handled in both the grammar and semantics. A command is executed even if it
           looks like a question. The. command is carried out by the following series of
           steps: SHRDLU puts a green cube on the large red block (it chooses the green
           cube with nothing on it), then removes the small pyramid from the little red
           cube, so that it can use that cube to complete the stack as specified.




SHRDLU, Toy of Man's Designing                                                       585

18. Eta Oin: which cube is sitting on the table?
        SHRDLU: THE LARGE GREEN ONE WHICH SUPPORTS THE RED
        PYRAMID.
        Dr. Tony Earrwig: Ordinarily SHRDLU would be satisfied with calling this "A
        BIG GREEN CUBE", but since you asked "which cube" it has to be more spec /lc
        in naming it. It has procedures for using substitute nouns (such as "one") and
        pronouns in answers to make the discourse more natural.
19. Eta Oin: is there a large block behind a pyramid?
        SHRDLU: YES, THREE OF THEM: A LARGE RED ONE, A LARGE GREEN
        CUBE AND THE BLUE ONE.
20. Eta Oin: put a small one onto the green cube which supports a pyramid.
        Dr. Tony Earrwig: "A small one" must contrast with your earlier phrase "a large
        block", so this means "a small block". SHRDLU: OK,
21. Eta Oin: put the littlest pyramid on top of it.
        Dr. Tony Earrwig: Words like "littlest" are not in SHRDLU's dictionary but are
        interpreted from the root farms like "little". SHRDLU compares the sizes of all
        the pyramids to carry out the action.
        SHRDLU• OK.
22. Eta Oin: how many things are on top of green cubes?
        SHRDLU: I'M NOT SURE WHAT YOU MEAN BY "ON TOP OF" IN THE
        PHRASE "ON TOP OF GREEN CUBES"
        DO YOU MEAN
        1-DIRECTLY ON THE SURFACE
        2-ANYWHERE ON TOP OF IT?
        Dr. Tony Earrwig: In the previous sentence, the phrase "on top of " was
        interpreted without trouble because the context made it clear what was meant. In
           this one, the heuristics for disambiguation do not come up with a clearcut
           definition, so SHRDLU had to ask you for more information. The paraphrases
           were included with the dictionary definitions of "on". A request like this can
           also occur when the meaning of a pronoun such as "it" or "they" is unclear.
23. Eta Din: had you touched any pyramid before you put the green one on the little
        cube?
        SHRDLU: YES, THE GREEN ONE.
        Dr. Tony Earrwig: SHRDLU deduces from its memory what event is being
            referred to, and can relate other events to it in time. The analysis includes
            complex tenses of verb groups, and there are special facilities for keeping
            track of the various forms of irregular verbs like "have".
24. Eta Oin: when did you pick it up?
        SHRDLU: WHILE I WAS STACKING UP THE RED CUBE, A LARGE RED
        BLOCK AND A LARGE GREEN ONE.
        Dr. Tony Earrwig: A time is described by naming the main goal which was being
        worked on.
25. Eta Oin: why?




SHRDLU, Toy of Man's Designing                                                       586

        Dr. Tony Earrwig: SHRDLU remembers immediately previous sentences to
        understand questions like this and the following ones.
        SHRDLU: TO GET RID OF IT.
        Dr. Tony Earrwig: By keeping track of selected parts of the original subgoal tree,
        SHRDLU has some understanding of its own motives.
26. Eta Oin: why did you do that?
        SHRDLU: TO CLEAN OFF THE RED CUBE.
        Dr. Tony Earrwig: "That" refers to the action mentioned in SHRDLU's last
        answer.
27. Eta Oin: why did you clear off that cube?
        Dr. Tony Earrwig: Here, "that" means the one just referred to. SHRDLU: TO
        PUT IT ON A LARGE GREEN CUBE.
28. Eta Oin: why did you do that?
        SHRDLU: TO STACK UP THE RED CUBE, A LARGE RED BLOCK, AND A
        LARGE GREEN CUBE.
29. Eta Oin: why did you do that?
        SHRDLU.' BECAUSE YOU ASKED ME TO.
        Dr. Tony Earrwig: SHRDLU remembers what it did, not how the request was
        worded.
30. Eta Oin: how did you do it?
        Dr. Tony Earrwig: "It" refers to the action mentioned in the last question,
        rather than the answer.
        SHRDLU: BY PUTTING A LARGE RED BLOCK ON THE TABLE; THEN
        LETTING GO OF IT; THEN PUTTING A LARGE GREEN CUBE ON IT;
        THEN LETTING GO OF THAT CUBE; THEN PUTTING THE RED CUBE
‘       ON THAT CUBE; THEN LETTING GO OF THAT CUBE.
        Dr. Tony Earrwig: This is a description of the top-level subgoals used to
        achieve the goal.
31. Eta Oin: how many objects did you touch while you were doing it?
        Dr. Tony Earrwig: Time references can refer to a time span, using words like
        "while".
        SHRDLU: FOUR OF THEM.
32. Eta Oin: what did the red cube support before you started to clean it off?
        SHRDLU: THE GREEN PYRAMID.
        Dr. Tony Earrwig: By remembering what has been moved, it is possible to
        deduce things about earlier states of the scene. You can use complex nestings
        of clauses like "before you started to clean it off ".
33. Eta Oin: there were five blocks to the left of the box then.
        Dr. Tony Earrwig: "Then" must refer to the time mentioned in the last sentence.
        SHRDLU: NO, ONLY FOUR OF THEM: THE RED CUBE, TWO LARGE
        GREEN CUBES, AND A LARGE RED BLOCK.
        Dr. Tony Earrwig: Earlier in the dialogue, information about "owning" was
        accepted at face value. But SHRDLU knows that it has complete infornia-




SHRDLU, Toy of Man's Designing                                                        587

        tion about the basic properties of objects, including location. It checks any
        statements about these to see whether it agrees or disagrees.
34. Eta Oin: put the blue pyramid on the block in the box.
        Dr. Tony Earrwig: This is syntactically ambiguous, but is understood without
        ambiguity since the parsing corresponding to "put the blue pyramid which is
        on the block into the box" does not make sense in the current scene, whereas
        "put the blue pyramid onto the block which is in the box" does. SHRDLU: OK.
35. Eta Oin: is there anything which is bigger than every pyramid but is not as wide as
        the thing that supports it?
        Dr. Tony Earrwig: SHRDLU understands quantifiers such as "every", "any",
        "some", etc., and structures such as "bigger than", "as wide as", etc.. Here "it"
        refers to a noun group within the same sentence but two levels up in the parsing
        tree.
        SHRDLU: YES, THE BLUE BLOCK.
36. Eta Oin: does a steeple
        SHRDLU: SORRY, I DON'T KNOW THE WORD "STEEPLE".
        Dr. Tony Earrwig: SHRDLU interrupts immediately when it sees a word it doesn't
        know.
37. Eta Oin: a "steeple" is a stack which contains two green cubes and a pyramid.
        Dr. Tony Earrwig: The quote marks indicate that you are defining a new word.
        The definition is added to SHRDLU's dictionary in a form compatible with all the
        words it already knows.
        SHRDLU: I UNDERSTAND.
38. Eta Oin: are there any steeples now?
        SHRDLU: NO.
        Dr. Tony Earrwig: The new definition is used as a description to recognize
        steeples.
39. Eta Oin: build one.
        SHRDLU: OK.
        Dr. Tony Earrwig: The description given in the definition is enough to deduce a
        plan for building new objects as well as recognizing them. The word "one" means
        "a steeple", and the plan is carried out. First, one of the large green cubes is
        cleared off by removing the red cube and two pyramids. Then the two green cubes
        and the tall red pyramid are stacked up to fit the description.
40. Eta Oin: call the biggest block "superblock".
        SHRDLU: OK.
        Dr. Tony Earrwig: You can also give objects names.
41. Eta Oin: have you picked up superblock since we began?
        SHRDLU: YES.
42. Eta Oin: why did you drop it?
        SHRDLU: BY "IT", I ASSUME YOU MEAN THE SUPERBLOCK. I CAN'T
        EXPLAIN A NON-EXISTENT EVENT.




SHRDLU, Toy of Man's Designing                                                       588

        Dr. Tony Earrwig: There is no record in SHRDLU's memory of any block being
        dropped. The answer routines have special checks for this type of "loaded
        question" which presupposes things which aren't true.
43. Eta Oin: is there anything to the right of the red pyramid?
        SHRDLU: YES, FOUR THINGS: THE BOX, THE RED CUBE, THE BLUE
        PYRAMID AND THE BLOCK NAMED SUPERBLOCK.
        Dr. Tony Earrwig: Once an object has a name, the name is used to help identify it
        in responses.
44. Eta Oin: thank you.
        SHRDLU: YOU'RE WELCOME!
        Dr. Tony Earrwig: SHRDLU's grammar is written as a procedure, which gives
        SHRDLU the capability to carry out special sorts of parsing for
        idioms.




SHRDLU, Toy of Man's Designing                                                       589


                            Artificial Intelligence:
                                 Retrospects
                                       Turing
                                            IN 1950, ALAN TURING wrote a most
                                            prophetic and provocative article on
                                            Artificial Intelligence. It was entitled
                                            "Computing          Machinery          and
                                            Intelligence" and appeared in the journal
                                            Mind.' I will say some things about that
                                            article, but I would like to precede them
                                            with some remarks about Turing the
                                            man.
                                            Alan Mathison Turing was born in
                                            London in 1912. He was a child full of
                                            curiosity and humor. Gifted in
                                            mathematics, he went to Cambridge
                                            where his interests in machinery and
                                            mathematical logic cross-fertilized and
                                            resulted in his famous paper on
                                            "computable numbers", in which he
                                            invented the theory of Turing machines
                                            and demonstrated the unsolvability of
                                            the halting problem; it was published in
                                            1937. In the 1940's, his interests turned
                                            from the theory of computing machines
                                            to the actual building of real computers.
                                            He was a major figure in the
                                            development of computers in Britain,
                                            and a staunch defender of Artificial In-

                                            FIGURE 113. Alan Turing, after a
                                            successful race (May, 1950). [From
                                            Sara Turing, Alan M. Turing
                                            (Cambridge, U. K.:W. Heffer \& Sons,
                                            1959).




Artificial Intelligence: Retrospects                                              590

telligence when it first came under attack. One of his best friends was David
Champernowne (who later worked on computer composition of music). Champernowne
and Turing were both avid chess players and invented "round-the-house" chess: after
your move, run around the house-if you get back before your opponent has moved, you're
entitled to another move. More seriously, Turing and Champernowne invented the first
chess playing program, called "Turochamp . Turing died young, at 41-apparently of an
accident with chemicals. Or some say suicide. His mother, Sara Turing, wrote his
biography. From the people she quotes, one gets the sense that Turing was highly
unconventional, even gauche in some ways, but so honest and decent that he was
vulnerable to the world. He loved games, chess, children, and bike riding; he was a strong
long-distance runner. As a student at Cambridge, he bought himself a second-hand violin
and taught himself to play. Though not very musical, he derived a great deal of
enjoyment from it. He was somewhat eccentric, given to great bursts of energy in the
oddest directions. One area he explored was the problem of morphogenesis in biology.
According to his mother, Turing "had a particular fondness for the Pickwick Papers", but
"poetry, with the exception of Shakespeare's, meant nothing to him." Alan Turing was
one of the true pioneers in the field of computer science.

                                       The Turing Test
Turing's article begins with the sentence: "I propose to consider the question `Can
machines think?"' Since, as he points out, these are loaded terms, it is obvious that we
should search for an operational way to approach the question. This, he suggests, is
contained in what he calls the "imitation game"; it is nowadays known as the Turing test.
Turing introduces it as follows:

   It is played with three people: a man (A), a woman (B), and an interrogator (C)
   who may be of either sex. The interrogator stays in a room apart from the other
   two. The object of the game for the interrogator is to determine which of the other
   two is the man and which is the woman. He knows them by labels X and Y, and at
   the end of the game he says either "X is A and Y is B" or "X is B and Y is A". The
   interrogator is allowed to put questions to A and B thus:

   C: Will X please tell me the length of his or her hair?

   Now suppose X is actually A, then A must answer. It is A's object in the game to
   try to cause C to make the wrong identification. His answer might therefore be

   "My hair is shingled, and the longest strands are about nine inches long."

   In order that tones of voice may not help the interrogator the answers should be
   written, or better still, typewritten. The ideal arrangement is to have a teleprinter
   communicating between the two rooms. Alternatively the questions and answers
   can be repeated by an intermediary. The object of the game for the third player (B)
   is to help the interrogator. The best strategy for her is probably to give truthful
   answers. She can add such things as "I am the woman, don't listen to him!" to her
   answers, but it will avail nothing as the man can make similar remarks.

Artificial Intelligence: Retrospects                                                   591

   We now ask the question, "What will happen when a machine takes the part of A
   in this game Will the interrogator decide wrongly as often when the game is
   played like this as he does when the game is played between a man and a woman?
   These questions replace our original, "Can machines think?"2

After having spelled out the nature of his test, Turing goes on to make some
commentaries on it, which, given the year he was writing in, are quite sophisticated. To
begin with, he gives a short hypothetical dialogue between interrogator and interrogatee:3

       Q. Please write me a sonnet on the subject of the Forth Bridge [a bridge over the
       Firth of Forth, in Scotland].
       A. Count me out on this one. I never could write poetry.
       Q. Add 34957 to 70764.
       A. (Pause about 30 seconds and then give as answer) 105621.
       Q. Do you play chess?
       A. Yes.
       Q. I have K at my K1, and no other pieces. You have only K at K6 and R at R1. It
       is your move. What do you play?
       A. (After a pause of 15 seconds) R-R8 mate.

Few readers notice that in the arithmetic problem, not only is there an inordinately long
delay, but moreover, the answer given is wrong! This would be easy to account for if the
respondent were a human: a mere calculational error. But if the respondent were a
machine, a variety of explanations are possible. Here are some:

       (1) a run-time error on the hardware level (i.e., an irreproducible fluke);
       (2) an unintentional hardware (or programming) (reproducibly) causes
       arithmetical mistakes;
       (3) a ploy deliberately inserted by the machine's programmer (or builder) to
       introduce occasional arithmetical mistakes, so as to trick interrogators;
       (4) an unanticipated epiphenomenon: the program has a hard time thinking
       abstractly, and simply made "an honest mistake", which it might not make the
       next time around;
       (5) a joke on the part of the machine itself, deliberately teasing its interrogator.

Reflection on what Turing might have meant by this subtle touch opens up just about all
the major philosophical issues connected with Artificial Intelligence.
       Turing goes on to point out that

   The new problem has the advantage of drawing a fairly sharp line between the
   physical and the intellectual capacities of a man. . . . We do not wish to
   penalize the machine for its inability to shine in beauty competitions, nor to
   penalize a man for losing in a race against an airplane.'

One of the pleasures of the article is to see how far Turing traced out each



Artificial Intelligence: Retrospects                                                   592

line of thought, usually turning up a seeming contradiction at some stage and, by refining
his concepts, resolving it at a deeper level of analysis. Because of this depth of
penetration into the issues. the article still shines after nearly thirty years of tremendous
progress in computer development and intensive work in Al. In the following short
excerpt you can see some of this rich back-and-forth working of ideas:

   The game may perhaps be criticized on the ground that the odds are weighted too
   heavily against the machine. If the man were to try to pretend to be the machine he
   would clearly make a very poor showing. He would be given away at once by
   slowness and inaccuracy in arithmetic. May not machines carry out something
   which ought to be described as thinking but which is very different from what a
   man does: This objection is a very strong one, but at least we can say that if,
   nevertheless, a machine can be constructed to play the imitation game satisfactorily,
   we need not be troubled by this objection.
           It might be urged that when playing the "imitation game" the best strategy
   for the machine may possibly be something other than imitation of the behaviour of
   a man. This may be, but I think it is unlikely that there is any greet effect of this
   kind. In any case there is ,no intention to investigate here the theory of the game,
   and it will be assumed that the best strategy is to try to provide answers that would
   naturally be given by a mans

Once the test has been proposed and discussed, Turing remarks:

   The original question "Can machines think 1 believe to be too meaningless to
   deserve discussion. Nevertheless, I believe that at the end of the century the use of
   words and general educated opinion will have altered so much that one will be able
   to speak of machines thinking without expecting to be contradicted.6

                           Turing Anticipates Objections
Aware of the storm of opposition that would undoubtedly greet this opinion, he then
proceeds to pick apart, concisely and with wry humor, a series of objections to the notion
that machines could think. Below I list the nine types of objections he counters, using his
own descriptions of them .7 Unfortunately there is not space to reproduce the humorous
and ingenious responses he formulated. You may enjoy pondering the objections
yourself, and figuring out your own responses.

   (1) The Theological Objection. Thinking is a function of man's immortal soul. God
     has given an immortal soul to every man and woman, but not to any other animal
     or to machines. Hence no animal or machine can think.
   (2) The "Heads in the Sand" Objection. The consequences of machines thinking
     would be too dreadful. Let us hope and believe that they cannot do so.
   (3) The Mathematical Objection. [This is essentially the Lucas argument.
   (4) The Argument from Consciousness. "Not until a machine can write a sonnet or
     compose a concerto because of thoughts and emotions felt, and not by the chance
     fall of symbols, could we agree that machine equals brainy that is, not only write
     it but know that it had written it. No mechanism


Artificial Intelligence: Retrospects                                                     593

     could feel (and not merely artificially signal, an easy contrivance) pleasure at its
     successes, grief when its valves fuse, be warmed by flattery, be made miserable
     by its mistakes, be charmed by sex, be angry or depressed when it cannot get
     what it wants." [A quote from a certain Professor Jefferson.]
Turing is quite concerned that he should answer this serious objection in full detail.
Accordingly, he devotes quite a bit of space to his answer, and in it he offers another
short hypothetical dialogue:'
   Interrogator: In the first line of your sonnet which reads "Shall I compare thee to a
     summer's day", would not "a spring day" do as well or better, Witness: It
     wouldn't scan.
   Interrogator: How about ' a winter's day''? That would scan all right. Witness: Yes,
     but nobody wants to be compared to a winter's day. Interrogator: Would you say
     Mr. Pickwick reminded you of Christmas? Witness: In a way.
   Interrogator: Yet Christmas is a winter's day, and I do not think Mr. Pickwick
     would mind the comparison.
   Witness: I don't think you're serious. By a winter's day one means a typical winter's
     day, rather than a special one like Christmas.
After this dialogue, Turing asks, "What would Professor Jefferson say if the sonnet-
writing machine was able to answer like this in the viva voce?"
        Further objections:
   (5) Arguments from various Disabilities. These arguments take the form, "I grant you that
      you can make machines do all the things that you have mentioned but you will never be
      able to make one to do X." Numerous features X are suggested in this connection. I offer
      a selection:
   Be kind, resourceful, beautiful, friendly, have initiative, have a sense of humor, tell right
      from wrong, make mistakes, fall in love, enjoy strawberries and cream, make someone
      fall in love with it, learn from experience, use words properly, be the subject of its own
      thought, have as much diversity of behaviour as a man, do something really new.
   (6) Lady Lovelace's Objection. Our most detailed information of Babbage's Analytical
      Engine comes from a memoir by Lady Lovelace. In it she states, "The Analytical Engine
      has no pretensions to originate anything. It can do whatever we know how to order it to
      perform" (her italics).
   (7) Argument from Continuity in the Nervous System. The nervous system is certainly not
      a discrete state machine. A small error in the information about the size of a nervous
      impulse impinging on a neuron may make a large difference to the size of the outgoing
      impulse. It may be argued that, this being so, one cannot expect to be able to mimic the
      behaviour of the nervous system with a discrete state system.
   (8) The Argument from Informality of Behaviour. It seems to run something like this. "If
      each man had a definite set of rules of conduct by which he regulated his life he would
      be no better than a machine. But there are no such rules, so men cannot be machines."
   (9) The Argument from Extra-Sensory Perception. Let us play the imitation game, using as
      witnesses a man who is good as a telepathic receiver, and a digital computer. The
      interrogator can ask such questions as "What suit




Artificial Intelligence: Retrospects                                                           594

     does the card in my right hand belong to?" The man by telepathy or clairvoyance gives
     the right answer 130 times out of 400 cards. The machine can only guess at random, and
     perhaps get 104 right, so the interrogator makes the right identification.

As you can see, many of the objections overlap and are related to each other. I have
attempted in this book to answer most of them in some way or other, usually far less
concisely than Turing did.
   Objection (9) I find remarkable. I have seen the Turing paper reprinted in a book-
   but with objection (9) omitted-which I find equally remarkable. My own point of
   view-contrary to Turing's-is that ESP does not exist. Turing was reluctant to accept
   the idea that ESP is real, but did so nonetheless, being compelled by his outstanding
   scientific integrity to accept the consequences of what he viewed as powerful
   statistical evidence in favor of ESP. I disagree, though I consider it an exceedingly
   complex and fascinating question.
                               "Parry Encounters the Doctor"

In the Dialogue preceding this Chapter, you have seen an authentic exchange between
a computer program and a human. Now there have been many programs with less
ability to "converse" than SHRDLU. Two rather famous ones are "Doctor", created by
Joseph Weizenbaum, and "Parry", created by Kenneth Colby. The former is supposed
to simulate a psychiatrist using "nondirective therapy", the latter to simulate the belief
structure of a paranoid. An amusing head-on collision between the two programs was
arranged over a computer network, and the following is a short excerpt from the full
transcript:9
       Parry: I don't understand your motives.
       Doctor: Why don't you understand my motives?
       Parry: Shouldn't I be?
       Doctor: Why do you say that just now?
       Parry: That's just the way it is.
       Doctor: Suppose that's just the way it were not.
       Parry: So?
       Doctor: I am not sure I understand you fully.
       Parry: You harp on the same questions too much.
       Doctor: What does that similarity suggest to you?
       Parry: I don't get you.
       Doctor: Do you wish that you did get I?
       Parry: It's about time you showed some interest in my feelings.
       Doctor: Oh? I showed some interest in your feelings?

These programs, of course, were made to talk to humans, so they are shown up rather
poorly here, pitted against each other. Mostly they rely on shrewd guesses as to the
nature of the input (which they analyze quite shallowly) and spit back canned answers
that have been carefully selected from a large repertoire. The answer may be only
partially canned: for example, a template with blanks that can be filled in. It is
assumed that their




Artificial Intelligence: Retrospects                                                      595

human partners will read much more into what they say than is actually underlying it.
And in fact, according to Weizenbaum, in his book Computer Power and Human
Reason, just that happens. He writes:

   ELIZA [the program from which Doctor was made created the most remarkable
   illusion of having understood in the minds of the many people who conversed with
   it.... They would often demand to be permitted to converse with the system in
   private, and would, after conversing with it for a time, insist, in spite of my
   explanations, that the machine really understood them.10

Given the above excerpt, you may find this incredible. Incredible, but true.
Weizenbaum has an explanation:

   Most men don't understand computers to even the slightest degree. So, unless they
   are capable of very great skepticism (the kind we bring to bear while
   watching a stage magician), they can explain the computer's intellectual feats only
   by bringing to hear the single analogy available to them, that is, their
   model of their own capacity to think. No wonder, then, that they overshoot the
   mark: it is truly impossible to imagine a human who could imitate ELIZA,
   for example, but for whom ELIZA's language abilities were his limit."

Which amounts to an admission that this kind of program is based on a shrewd
mixture of bravado and bluffing, taking advantage of people's gullibility.
   In light of this weird "ELIZA-effect", some people have suggested that the Turing
test needs revision, since people can apparently be fooled by simplistic gimmickry. It
has been suggested that the interrogator should be a Nobel Prize-winning scientist. It
might be more advisable to turn the Turing test on its head, and insist that the
interrogator should be another computer. Or perhaps there should be two
interrogators-a human and a computer-and one witness, and the two interrogators
should try to figure out whether the witness is a human or a computer.
In a more serious vein, I personally feel that the Turing test, as originally proposed, is
quite reasonable. As for the people who Weizenbaum claims were sucked in by
ELIZA, they were not urged to be skeptical, or to use all their wits in trying to
determine if the "person" typing to them were human or not. I think that Turing's
insight into this issue was sound, and that the Turing test, essentially unmodified, will
survive.

                                 A Brief History of AI

I would like in the next few pages to present the story, perhaps from an unorthodox point
of view, of some of the efforts at unraveling the algorithms behind intelligence: there
have been failures and setbacks and there will continue to be. Nonetheless, we are
learning a great deal, and it is an exciting period.
   Ever since Pascal and Leibniz, people have dreamt of machines that could perform
intellectual tasks. In the nineteenth century, Boole and De Morgan devised "laws of
thought"-essentially the Propositional

Artificial Intelligence: Retrospects                                                     596

Calculus-and thus took the first step towards At software; also Charles Babbage designed
the first "calculating engine"-the precursor to the hardware of computers and hence of AI.
One could define AI as coming into existence at the moment when mechanical devices
took over any tasks previously performable only by human minds. It is hard to look back
and imagine the feelings of those who first saw toothed wheels performing additions and
multiplications of large numbers. Perhaps they experienced a sense of awe at seeing
"thoughts" flow in their very physical hardware. In any case, we do know that nearly a
century later, when the first electronic computers were constructed, their inventors did
experience an awesome and mystical sense of being in the presence of another kind of
"thinking being". To what extent real thought was taking place was a source of much
puzzlement; and even now, several decades later, the question remains a great source of
stimulation and vitriolics.
     It is interesting that nowadays, practically no one feels that sense of awe any longer-
even when computers perform operations that are incredibly more sophisticated than
those which sent thrills down spines in the early days. The once-exciting phrase "Giant
Electronic Brain" remains only as a sort of "camp" cliché, a ridiculous vestige of the era
of Flash Gordon and Buck Rogers. It is a bit sad that we become blasé so quickly.
     There is a related "Theorem" about progress in Al: once some mental function is
programmed, people soon cease to consider it as an essential ingredient of "real
thinking". The ineluctable core of intelligence is always in that next thing which hasn't
yet been programmed. This "Theorem" was first proposed to me by Larry Tesler, so I call
it Tesler's Theorem. "Al is whatever hasn't been done vet."
     A selective overview of AI is furnished below. It shows several domains in which
workers have concentrated their efforts, each one seeming in its own way to require the
quintessence of intelligence. With some of the domains I have included a breakdown
according to methods employed, or more specific areas of concentration.

    mechanical translation
    direct (dictionary look-up with some word rearrangement)
    indirect (via some intermediary internal language)

    game playing

         chess
            with brute force look-ahead
         with heuristically pruned look-ahead
         with no look-ahead checkers
    go
    kalah
    bridge (bidding; playing)
    poker
    variations on tic-tac-toe
    etc.




Artificial Intelligence: Retrospects                                                    597

   proving theorems in various parts. of mathematics
       symbolic logic
              "resolution" theorem-proving
              elementary geometry

   symbolic manipulation of mathematical expressions
      symbolic integration
      algebraic simplification
      summation of infinite series

   vision
        printed matter:
                recognition of individual hand-printed characters drawn
                        from a small class (e.g., numerals)
                reading text in variable fonts reading passages in handwriting
                reading Chinese or Japanese printed characters
                reading Chinese or Japanese handwritten characters
        pictorial:
                locating prespecified objects in photographs
                decomposition of a scene into separate objects
                identification of separate objects in a scene
                recognition of objects portrayed in sketches by people
                recognition of human faces
   hearing
        understanding spoken words drawn from a limited vocabulary (e.g., names of
            the ten digits)
        understanding continuous speech in fixed domains finding boundaries between
            phonemes
        identifying phonemes
        finding boundaries between morphemes
        identifying morphemes
        putting together whole words and sentences

   understanding natural languages
     answering questions in specific domains
      parsing complex sentences
     making paraphrases of longer pieces of text
     using knowledge of the real world in order to understand passages
     resolving ambiguous references

   producing natural language
       abstract poetry (e.g., haiku)
       random sentences, paragraphs, or longer pieces of text producing output from
       internal representation of knowledge




Artificial Intelligence: Retrospects                                              598

   creating original thoughts or works of art
        poetry writing (haiku) story writing
        computer art
        musical composition
               atonal
               tonal

   analogical thinking
       geometrical shapes ("intelligence tests")
       constructing proofs in one domain of mathematics based on
               those in a related domain

   learning
        adjustment of parameters
        concept formation

                                 Mechanical Translation

   Many of the preceding topics will not be touched upon in my selective discussion
   below, but the list would not be accurate without them. The first few topics are
   listed in historical order. In each of them, early efforts fell short of expectations.
   For example, the pitfalls in mechanical translation came as a great surprise to many
   who had thought it was a nearly straightforward task, whose perfection, to be sure,
   would be arduous, but whose basic implementation should be easy. As it turns out,
   translation is far more complex than mere dictionary look-up and word rearranging.
   Nor is the difficulty caused by a lack of knowledge of idiomatic phrases. The fact is
   that translation involves having a mental model of the world being discussed, and
   manipulating symbols in that model. A program which makes no use of a model of
   the world as it reads the passage will soon get hopelessly bogged down in
   ambiguities and multiple meanings. Even people-who have a huge advantage over
   computers, for they come fully equipped with an understanding of the world-when
   given a piece of text and a dictionary of a language they do not know, find it next to
   impossible to translate the text into their own language. Thus-and it is not
   surprising in retrospect-the first problem of AI led immediately to the issues at the
   heart of AI.

                                       Computer Chess

   Computer chess, too, proved to be much more difficult than the early intuitive
   estimates had suggested. Here again it turns out that the way humans represent a
   chess situation in their minds is far more complex than just knowing which piece is
   on which square, coupled with knowledge of the rules of chess. It involves
   perceiving configurations of several related pieces, as well as knowledge of
   heuristics, or rules of thumb, which pertain to



Artificial Intelligence: Retrospects                                                    599

   such higher-level chunks. Even though heuristic rules are not rigorous in the way
   that the official rules are, they provide shortcut insights into what is going on on the
   board, which knowledge of the official rules does not. This much was recognized
   from the start; it was simply underestimated how large a role the intuitive, chunked
   understanding of the chess world plays in human chess skill. It was predicted that a
   program having some basic heuristics, coupled with the blinding speed and
   accuracy of a computer to look ahead in the game and analyze each possible move,
   would easily beat top-flight human players-a prediction which, even after twenty-
   five years of intense work by various people, still is far from being realized.
        People are nowadays tackling the chess problem from various angles. One of
   the most novel involves the hypothesis that looking ahead is a silly thing to do. One
   should instead merely look at what is on the board at present, and, using some
   heuristics, generate a plan, and then find a move which advances that particular
   plan. Of course, rules for the formulation of chess plans will necessarily involve
   heuristics which are, in some sense, "flattened" versions of looking ahead. That is,
   the equivalent of many games' experience of looking ahead is "squeezed" into
   another form which ostensibly doesn't involve looking ahead. In some sense this is
   a game of words. But if the "flattened" knowledge gives answers more efficiently
   than the actual look-ahead-even if it occasionally misleads- then something has
   been gained. Now this kind of distillation of knowledge into more highly usable
   forms is just what intelligence excels at-so look-ahead-less chess is probably a
   fruitful line of research to push. Particularly intriguing would be to devise a
   program which itself could convert knowledge gained from looking ahead into
   "flattened" rules-but that is an immense task.

                             Samuel's Checker Program
   As a matter of fact, such a method was developed by Arthur Samuel in his
   admirable checker-playing program. Samuel's trick was to use both dynamic (look-
   ahead) and static (no-look-ahead) ways of evaluating any given board position. The
   static method involved a simple mathematical function of several quantities
   characterizing any board position, and thus could be calculated practically
   instantaneously, whereas the dynamic evaluation method involved creating a "tree"
   of possible future moves, responses to them, responses to the responses, and so
   forth (as was shown in Fig. 38). In the static evaluation function there were some
   parameters which could vary; the effect of varying them was to provide a set of
   different possible versions of the static evaluation function. Samuel's strategy was
   to select, in an evolutionary way, better and better values of those parameters.
        Here's how this was done: each time the program evaluated a board position, it
   did so both statically and dynamically. The answer gotten by looking ahead-let us
   call it D-was used in determining the move to be made. The purpose of S, the static
   evaluation, was trickier: on each move, the variable parameters were readjusted
   slightly so that S approximated D




Artificial Intelligence: Retrospects                                                      600

   as accurately as possible. The effect was to partially encode in the values of the
   static evaluation’s parameters the knowledge gained by dynamically searching the
   tree. In short, the idea was to "flatten" the complex dynamic evaluation method into
   the much simpler and more efficient static evaluation function.
        There is a rather nice recursive effect here. The point is that the dynamic
   evaluation of any single board position involves looking ahead a finite number of
   moves-say seven. Now each of the scads of board positions which might turn up
   seven turns down the road has to be itself evaluated somehow as well. But when the
   program evaluates these positions, it certainly cannot look another seven moves
   ahead, lest it have to look fourteen positions ahead, then twenty-one, etc., etc.-an
   infinite regress. Instead, it relies on static evaluations of positions seven moves
   ahead. Therefore, in Samuel's scheme, an intricate sort of feedback takes place,
   wherein the program is constantly trying to "flatten" look-ahead evaluation into a
   simpler static recipe; and this recipe in turn plays a key role in the dynamic look-
   ahead evaluation. Thus the two are intimately linked together, and each benefits
   from improvements in the other in a recursive way.
        The level of play of the Samuel checkers program is extremely high: of the
   order of the top human players in the world. If this is so, why not apply the same
   techniques to chess? An international committee, convened in 1961 to study the
   feasibility of computer chess, including the Dutch International Grandmaster and
   mathematician Max Euwe, came to the bleak conclusion that the Samuel technique
   would be approximately one million times as difficult to implement in chess as in
   checkers, and that seems to close the book on that.
        The extraordinarily great skill of the checkers program cannot be taken as
   saying "intelligence has been achieved"; yet it should not be minimized, either. It is
   a combination of insights into what checkers is, how to think about checkers, and
   how to program. Some people might feel that all it shows is Samuel's own checkers
   ability. But this is not true, for at least two reasons. One is that skillful game players
   choose their moves according to mental processes which they do not fully
   understand-they use their intuitions. Now there is no known way that anyone can
   bring to light all of his own intuitions; the best one can do via introspection is to
   use "feeling" or "meta-intuition"-an intuition about one's intuitions-as a guide, and
   try to describe what one thinks one's intuitions are all about. But this will only give
   a rough approximation to the true complexity of intuitive methods. Hence it is
   virtually certain that Samuel has not mirrored his own personal methods of play in
   his program. The other reason that Samuel's program's play should not be confused
   with Samuel's own play is that Samuel does not play checkers as well as his
   program-it beats him. This is not a paradox at all-no more than is the fact that a
   computer which has been programmed to calculate 7T can outrace its programmer
   in spewing forth digits of π.




Artificial Intelligence: Retrospects                                                        601

                         When Is a Program Original?
     This issue of a program outdoing its programmer is connected with the question of
"originality" in AI. What if an AI program comes up with an idea, or a line of play in a
game, which its programmer has never entertained-who should get the credit? There are
various interesting instances of this having happened, some on a fairly trivial level, some
on a rather deep level. One of the more famous involved a program to find proofs of
theorems in elementary Euclidean geometry, written by E. Gelernter. One day the
program came up with a sparklingly ingenious proof of one of the basic theorems of
geometry-the so-called "pons asinorum", or "bridge of asses".
     This theorem states that the base angles of an isosceles triangle are equal. Its
standard proof requires constructing an altitude which divides the triangle into
symmetrical halves. The elegant method found by the program (see Fig. 114) used no
construction lines. Instead, it considered

                                             FIGURE 114. Pons Asinorum Proof (found
                                             by Pappus [-300 A.D. ] and Gelernter's
                                             program [--1960 A.D. ]). Problem: To show
                                             that the base angles of an isosceles triangle
                                             are equal. Solution: As the triangle is
                                             isosceles, AP and AP' are of equal length.
                                             Therefore triangles PAP' and PAP are
                                             congruent (side-side-side). This implies that
                                             corresponding angles are equal. In
                                             particular, the two base angles are equal.


     the triangle and its mirror image as two different triangles. Then, having proved
them congruent, it pointed out that the two base angles matched each other in this
congruence-QED.
     This gem of a proof delighted the program's creator and others; some saw evidence
of genius in its performance. Not to take anything away from this feat, it happens that in
A.D. 300 the geometer Pappus had actually found this proof, too. In any case, the
question remains: "Who gets the credit?" Is this intelligent behavior? Or was the proof
lying deeply hidden within the human (Gelernter), and did the computer merely bring it
to the surface? This last question comes close to hitting the mark. We can turn it around:
Was the proof lying deeply hidden in the program? Or was it close to the surface? That is,
how easy is it to see why the program did what it did? Can the discovery be attributed to
some simple mechanism, or simple combination of mechanisms, in the program? Or was
there a complex interaction which, if one heard it explained, would not diminish one's
awe at its having happened?
     It seems reasonable to say that if one can ascribe the performance to certain
operations which are easily traced in the program, then in some sense the program was
just revealing ideas which were in essence hiddenthough not too deeply-inside the
programmer's own mind. Conversely, if




Artificial Intelligence: Retrospects                                                     602

following the program does not serve to enlighten one as to why this particular discovery
popped out, then perhaps one should begin to separate the program's "mind" from that of
its programmer. The human gets credit for having invented the program, but not for
having had inside his own head the ideas produced by the program. In such cases, the
human can be referred to as the "meta-author"-the author of the author of the result, and
the program as the (just plain) author.
     In the particular case of Gelernter and his geometry machine, while Gelernter
probably would not have rediscovered Pappus'. proof, still the mechanisms which
generated that proof were sufficiently close to the surface of the program that one
hesitates to call the program a geometer in its own right. If it had kept on astonishing
people by coming up with ingenious new proofs over and over again, each of which
seemed to be based on a fresh spark of genius rather than on some standard method, then
one would have no qualms about calling the program a geometer-but this did not happen.

                          Who Composes Computer Music?

The distinction between author and meta-author is sharply pointed up in the case of
computer composition of music. There are various levels of autonomy which a program
may seem to have in the act of composition. One level is illustrated by a piece whose
"meta-author" was Max Mathews of Bell Laboratories. He fed in the scores of the two
marches "When Johnny Comes Marching Home" and "The British Grenadiers", and
instructed the computer to make a new score-one which starts out as "Johnny", but slowly
merges into "Grenadiers". Halfway through the piece, "Johnny" is totally gone, and one
hears "Grenadiers" by itself ... Then the process is reversed, and the piece finishes with
"Johnny", as it began. In Mathews' own words, this is

        ... a nauseating musical experience but one not without interest, particularly in
   the rhythmic conversions. "The Grenadiers" is written in 2/4 time in the key of F
   major. "Johnny" is written in 6/8 time in the key of E minor. The change from 2/4
   to 6/8 time can be clearly appreciated, yet would be quite difficult for a human
   musician to play. The modulation from the key of F major to E minor, which
   involves a change of two notes in the scale, is jarring, and a smaller transition
   would undoubtedly have been a better choice."

    The resulting piece has a somewhat droll quality to it, though in spots it is turgid and
confused.

       Is the computer composing? The question is best unasked, but it cannot be
   completely ignored. An answer is difficult to provide. The algorithms are
   deterministic, simple, and understandable. No complicated or hard-to understand
   computations are involved; no "learning" programs are used; no random processes
   occur; the machine functions in a perfectly mechanical and straightforward manner.
   However, the result is sequences of sound that are unplanned in fine detail by the
   composer, even though the over-all structure



Artificial Intelligence: Retrospects                                                    603

        of the section is completely and precisely specified. Thus the composer is often
   surprised, and pleasantly surprised, ar the details of the realization of his ideas. To
   this extent only is the computer composing. We call the process algorithmic
   composition, but we immediately re-emphasize that the algorithms are
   transparently simple."

     This is Mathews' answer to a question which he would rather "unask". Despite his
disclaimer, however, many people find it easier to say simply that the piece was
"composed by a computer". I believe this phrase misrepresents the situation totally. The
program contained no structures analogous to the brain's "symbols", and could not be said
in any sense to be "thinking" about what it was doing. To attribute the composition of
such a piece of music to the computer would be like attributing the authorship of this
book to the computerized automatically (often incorrectly) hyphenating phototypesetting
machine with which it was set.
     This brings up a question which is a slight digression from Al, but actually not a
huge one. It is this: When you see the word "I" or "me" in a text, what do you take it to be
referring to? For instance, think of the phrase "WASH ME" which appears occasionally
on the back of dirty trucks. Who is this "me"? Is this an outcry of some forlorn child who,
in desperation to have a bath, scribbled the words on the nearest surface? Or is the truck
requesting a wash? Or, perhaps, does the sentence itself wish to be given a shower? Or, is
it that the filthy English language is asking to be cleansed? One could go on and on in
this game. In this case, the phrase is a joke, and one is supposed to pretend, on some
level, that the truck itself wrote the phrase and is requesting a wash. On another level, one
clearly recognizes the writing as that of a child, and enjoys the humor of the misdirection.
Here, in fact, is a game based on reading the "me" at the wrong level.
     Precisely this kind of ambiguity has arisen in this book, first in the
Contracrostipunctus, and later in the discussions of Gödel’s string G (and its relatives).
The interpretation given for unplayable records was "I Cannot Be Played on Record
Player X", and that for unprovable statements was, "I Cannot Be Proven in Formal
System X". Let us take the latter sentence. On what other occasions, if any, have you
encountered a sentence containing the pronoun "I" where you automatically understood
that the reference was not to the speaker of the sentence, but rather to the sentence itself?
Very few, I would guess. The word "I", when it appears in a Shakespeare sonnet, is
referring not to a fourteen-line form of poetry printed on a page, but to a flesh-and-blood
creature behind the scenes, somewhere off stage.
     How far back do we ordinarily trace the "I" in a sentence? The answer, it seems to
me, is that we look for a sentient being to attach the authorship to. But what is a sentient
being? Something onto which we can map ourselves comfortably. In Weizenbaum's
"Doctor" program, is there a personality? If so, whose is it? A small debate over this very
question recently raged in the pages of Science magazine.
     This brings us back to the issue of the "who" who composes computer music. In
most circumstances, the driving force behind such pieces is a




Artificial Intelligence: Retrospects                                                     604

human intellect, and the computer has been employed, with more or less ingenuity, as a
tool for realizing an idea devised by the human. The program which carries this out is not
anything which we can identify with. It is a simple and single-minded piece of software
with no flexibility, no perspective on what it is doing, and no sense of self. If and when,
however, people develop programs which have those attributes, and pieces of music start
issuing forth from them, then I suggest that will be the appropriate, time to start splitting
up one's admiration: some to the programmer for creating such an amazing program, and
some to the program itself for its sense of music. And it seems to me that that will only
take place when the internal structure of such a program is based on something similar to
the "symbols" in our brains and their triggering patterns, which are responsible for the
complex notion of meaning. The fact of having this kind of internal structure would
endow the program with properties which would make us feel comfortable in identifying
with it, to some extent. But until then, I will not feel comfortable in saying "this piece
was composed by a computer".

                   Theorem Proving and Problem Reduction
Let us now return to the history of AI. One of the early things which people attempted to
program was the intellectual activity of theorem proving. Conceptually, this is no
different from programming a computer to look for a derivation of MU in the MIU-
system, except that the formal systems involved were often more complicated than the
MIU-system. They were versions of the Predicate Calculus, which is an extension of the
Propositional Calculus involving quantifiers. Most of the rules of the Predicate Calculus
are included in TNT, as a matter of fact. The trick in writing such a program is to instill a
sense of direction, so that the program does not wander all over the map, but works only
on "relevant" pathways-those which, by some reasonable criterion, seem to be leading
towards the desired string.
     In this book we have not dealt much with such issues. How indeed can you know
when you are proceeding towards a theorem, and how can you tell if what you are doing
is just empty fiddling? This was one thing which I hoped to illustrate with the MU-
puzzle. Of course, there can be no definitive answer: that is the content of the limitative
Theorems, since if you could always know which way to go, you could construct an
algorithm for proving any desired theorem, and that would violate Church's Theorem.
There is no such algorithm. (I will leave it to the reader to see exactly why this follows
from Church's Theorem.) However, this doesn't mean that it is impossible to develop any
intuition at all concerning what is and what is not a promising route; in fact, the best
programs have very sophisticated heuristics, which enable them to make deductions in
the Predicate Calculus at speeds which are comparable to those of capable humans.
     The trick in theorem proving is to use the fact that you have an overall goal-namely
the string you want to produce-in guiding you locally. One technique which was
developed for converting global goals




Artificial Intelligence: Retrospects                                                     605

into local strategies for derivations is called problem reduction. It is based on the idea that
whenever one has a long-range goal, there are usually subgoals whose attainment will aid
in the attainment of the main goal. Therefore if one breaks up a given problem into a
series of new subproblems, then breaks those in turn into subsubproblems, and so on, in a
recursive fashion, one eventually comes down to very modest goals which can
presumably be attained in a couple of steps. Or at least so it would seem ...
      Problem reduction got Zeno into hot water. Zeno's method, you recall, for getting
from A to B (think of B as the goal), is to "reduce" the problem into two subproblems:
first go halfway, then go the rest of the way. So now you have "pushed"-in the sense of
Chapter V-two subgoals onto your "goal stack". Each of these, in turn, will be replaced
by two subsubgoals and so on ad infinitum. You wind up with an infinite goal-stack,
instead of a single goal (Fig. 115). Popping an infinite number of goals off your stack will
prove to be tricky-which is just Zeno's point, of course.
      Another example of an infinite recursion in problem reduction occurred in the
Dialogue Little Harmonic Labyrinth, when Achilles wanted to have a Typeless Wish
granted. Its granting had to be deferred until permission was gotten from the Meta-Genie;
but in order to get permission to give permission, she had to summon the Meta-Meta-
Genie-and so on. Despite

              FIGURE 115. Zeno's endless goal tree, for getting from A to B.




Artificial Intelligence: Retrospects                                                       606

the infiniteness of the goal stack, Achilles got his wish. Problem reduction wins the day!
      Despite my mockery, problem reduction is a powerful technique for converting
global problems into local problems. It shines in certain situations, such as in the
endgame of chess, where the look-ahead technique often performs miserably, even when
it is carried to ridiculous lengths, such as fifteen or more plies. This is because the look-
ahead technique is not based on planning; it simply has no goals and explores a huge
number of pointless alternatives. Having a goal enables you to develop a strategy for the
achievement of that goal, and this is a completely different philosophy from looking
ahead mechanically. Of course, in the look-ahead technique, desirability or its absence is
measured by the evaluation function for positions, and that incorporates indirectly a
number of goals, principally that of not getting checkmated. But that is too indirect. Good
chess players who play against look-ahead chess programs usually come away with the
impression that their opponents are very weak in formulating plans or strategies.

                                Shandy and the Bone
     There is no guarantee that the method of problem reduction will work. There are
many situations where it flops. Consider this simple problem, for instance. You are a dog,
and a human friend has just thrown your favorite bone over a wire fence into another
yard. You can see your bone through the fence, just lying there in the grass-how luscious!
There is an open gate in the fence about fifty feet away from the bone. What do you do?
Some dogs will just run up to the\_ fence, stand next to it, and bark; others will dash up to
the open gate and double back to the lovely bone. Both dogs can be said to be exercising
the problem reduction technique; however, they represent the problem in their minds in
different ways, and this makes all the difference. The barking dog sees the subproblems
as (1) running to the fence, (2) getting through it, and (3) running to the bone-but that
second subproblem is a "toughie", whence the barking. The other dog sees the
subproblems as (1) getting to the gate; (2) going through the gate; (3) running to the
bone. Notice how everything depends on the way you represent the "problem space"-that
is, on what you perceive as reducing the problem (forward motion towards the overall
goal) and what you perceive as magnifying the problem (backward motion away from the
goal).

                           Changing the Problem Space
Some dogs first try running directly towards the bone, and when they encounter the
fence, something clicks inside their brain; soon they change course, and run over to the
gate. These dogs realize that what on first




Artificial Intelligence: Retrospects                                                     607

glance seemed as if it would increase the distance between the initial situation and the
desired situation-namely, running away from the bone but towards the open gate-actually
would decrease it. At first, they confuse physical distance with problem distance. Any
motion away from the bone seems, by definition, a Bad Thing. But then-somehow-they
realize that they can shift their perception of what will bring them "closer" to the bone. In
a properly chosen abstract space, moving towards the gate is a trajectory bringing the dog
closer to the bone! At every moment, the dog is getting "closer"-in the new sense-to the
bone. Thus, the usefulness of problem reduction depends on how you represent your
problem mentally. What in one space looks like a retreat can in another space look like a
revolutionary step forward.
        In ordinary life, we constantly face and solve variations on the dog and-bone
problem. For instance, if one afternoon I decide to drive one hundred miles south, but am
at my office and have ridden my bike to work, I have to make an extremely large number
of moves in what are ostensibly "wrong" directions before I am actually on my way in car
headed south. I have to leave my office, which means, say, heading east a few feet; then
follow the hall in the building which heads north, then west. Then I ride my bike home,
which involves excursions in all the directions of the compass; and I reach my home. A
succession of short moves there eventually gets me into my car, and I am off. Not that I
immediately drive due south. of course-I choose a route which may involve some
excursions north. west, or east, with the aim of getting to the freeway as quickly as
possible.
        All of this doesn't feel paradoxical in the slightest; it is done without even any
sense of amusement. The space in which physical backtracking is perceived as direct
motion towards the goal is built so deeply into my mind that I don't even see any irony
when I head north. The roads and hallways and so forth act as channels which I accept
without much fight, so that part of the act of choosing how to perceive the situation
involves just accepting what is imposed. But dogs in front of fences sometimes have a
hard time doing that, especially when that bone is sitting there so close, staring them in
the face, and looking so good. And when the problem space is just a shade more abstract
than physical space, people are often just as lacking in insight about what to do as the
barking dogs.
        In some sense all problems are abstract versions of the dog-and-bone problem.
Many problems are not in physical space but in some sort of conceptual space. When you
realize that direct motion towards the goal in that space runs you into some sort of
abstract "fence", you can do one of two things: (1) try moving away from the goal in
some sort of random way, hoping that you may come upon a hidden "gate" through
which you can pass and then reach your bone; or (2) try to find a new "space" in which
you can represent the problem, and in which there is no abstract fence separating you
from your goal-then you can proceed straight towards the goal in this new space. The first
method may seem like the lazy way to go, and the second method may seem like a
difficult and complicated way to go. And yet, solutions which involve restructuring the
problem space more often




Artificial Intelligence: Retrospects                                                     608

than not come as sudden flashes of insight rather than as products of a series of slow,
deliberate thought processes. Probably these intuitive flashes come from the extreme core
of intelligence-and, needless to say, their source is a closely protected secret of our
jealous brains.
        In any case, the trouble is not that problem reduction per se leads to failures; it is
quite a sound technique. The problem is a deeper one: how do you choose a good internal
representation for a problem? What kind of "space" do you see it in? What kinds of
action reduce the "distance" between you and your goal in the space you have chosen?
This can be expressed in mathematical language as the problem of hunting for an
approprate metric (distance function) between states. You want to find a metric in which
the distance between you and your goal is very small.
        Now since this matter of choosing an internal representation is itself a type of
problem-and a most tricky one, too-you might think of turning the technique of problem
reduction back on it! To do so, you would have to have a way of representing a huge
variety of abstract spaces, which is an exceedingly complex project. I am not aware of
anyone's having tried anything along these lines. It may be just a theoretically appealing,
amusing suggestion which is in fact wholly unrealistic. In any case, what Al sorely lacks
is programs which can "step back" and take a look at what is going on, and with this
perspective, reorient themselves to the task at hand. It is one thing to write a program
which excels at a single task which, when done by a human being, seems to require
intelligence-and it is another thing altogether to write an intelligent program! It is the
difference between the Sphex wasp (see Chapter XI), whose wired-in routine gives the
deceptive appearance of great intelligence, and a human being observing a Sphex wasp.

                           The I-Mode and the M-Mode Again

An intelligent program would presumably be one which is versatile enough to solve
problems of many different sorts. It would learn to do each different one and would
accumulate experience in doing so. It would be able to work within a set of rules and yet
also, at appropriate moments, to step back and make a judgment about whether working
within that set of rules is likely to be profitable in terms of some overall set of goals
which it has. It would be able to choose to stop working within a given framework, if
need be, and to create a new framework of rules within which to work for a while.
        Much of this discussion may remind you of aspects of the MU-puzzle. For
instance, moving away from the goal of a problem is reminiscent of moving away from
MU by making longer and longer strings which you hope may in some indirect way
enable you to make MU. If you are a naive "dog", you may feel you are moving away
from your "MU-bone" whenever your string increases beyond two characters; if you are a
more sophisticated dog, the use of such lengthening rules has an indirect justification,
something like heading for the gate to get your MU-bone.




Artificial Intelligence: Retrospects                                                      609

        Another connection between the previous discussion and the MU puzzle is the
two modes of operation which led to insight about the nature of the MU-puzzle: the
Mechanical mode, and the Intelligent mode. In the former, you are embedded within
some fixed framework; in the latter, you can always step back and gain an overview of
things. Having an overview is tantamount to choosing a representation within which to
work; and working within the rules of the system is tantamount to trying the technique of
problem reduction within that selected framework. Hardy's comment on Ramanujan's
style-particularly his willingness to modify his own hypotheses-illustrates this interplay
between the M-mode and the I-mode in creative thought.
        The Sphex wasp operates excellently in the M-mode, but it has absolutely no
ability to choose its framework or even to alter its M-mode in the slightest. It has no
ability to notice when the same thing occurs over and over and over again in its system,
for to notice such a thing would be to jump out of the system, even if only ever so
slightly. It simply does not notice the sameness of the repetitions. This idea (of not
noticing the identity of certain repetitive events) is interesting when we apply it to
ourselves. Are there highly repetitious situations which occur in our lives time and time
again, and which we handle in the identical stupid way each time, because we don't have
enough of an overview to perceive their sameness? This leads back to that recurrent
issue, "What is sameness?" It will soon come up as an Al theme, when we discuss pattern
recognition.

                            Applying Al to Mathematics
Mathematics is in some ways an extremely interesting domain to study from the Al point
of view. Every mathematician has the sense that there is a kind of metric between ideas in
mathematics-that all of mathematics is a network of results between which there are
enormously many connections. In that network, some ideas are very closely linked;
others require more elaborate pathways to be joined. Sometimes two theorems in
mathematics are close because one can be proven easily, given the other. Other times two
ideas are close because they are analogous, or even isomorphic. These are two different
senses of the word "close" in the domain of mathematics. There are probably a number of
others. Whether there is an objectivity or a universality to our sense of mathematical
closeness, or whether it is largely an accident of historical development is hard to say.
Some theorems of different branches of mathematics appear to us hard to link, and we
might say that they are unrelated-but something might turn up later which forces us to
change our minds. If we could instill our highly developed sense of mathematical
closeness-a "mathematician's mental metric", so to speak-into a program, we could
perhaps produce a primitive "artificial mathematician". But that depends on being able to
convey a sense of simplicity or "naturalness" as well, which is another major stumbling
block.
        These issues have been confronted in a number of AI projects. There




Artificial Intelligence: Retrospects                                                  610

is a collection of programs developed at MIT which go under the name «MACSYMA",
whose purpose it is to aid mathematicians in symbolic manipulation of complex
mathematical, expressions. This program has in it some sense of "where to go"-a sort of
"complexity gradient" which guides it from what we would generally consider complex
expressions to simpler ones. Part of MACSYMA's repertoire is a program called "SIN",
which does symbolic integration of functions; it is generally acknowledged to be superior
to humans in some categories. It relies upon a number of different skills, as intelligence
in general must: a vast body of knowledge, the technique of problem reduction, a large
number of heuristics, and also some special tricks.
          Another program, written by Douglas Lenat at Stanford, had as its aim to invent
concepts and discover facts in very elementary mathematics. Beginning with the notion
of sets, and a collection of notions of what is "interesting" which had been spoon-fed into
it, it "invented" the idea of counting, then the idea of addition, then multiplication, then-
among other things-the notion of prime numbers, and it went so far as to rediscover
Goldbach's conjecture! Of course these "discoveries" were all hundreds-even thousands-
of years old. Perhaps this may be explained in part by saying that the sense of
"interesting" was conveyed by Lenat in a large number of rules which may have been
influenced by his twentieth century training; nonetheless it is impressive. The program
seemed to run out of steam after this very respectable performance. An interesting thing
about it was that it was unable to develop or improve upon its own sense of what is
interesting. That seemed another level of difficulty up-or perhaps several levels up.

                The Crux of Al: Representation of Knowledge
Many of the examples above have been cited in order to stress that the way a domain is
represented has a huge bearing on how that domain is "understood". A program which
merely printed out theorems of TNT in a preordained order would have no understanding
of number theory; a program such as Lenat's with its extra layers of knowledge could be
said to have a rudimentary sense of number theory; and one which embeds mathematical
knowledge in a wide context of real-world experience would probably be the most able to
"understand" in the sense that we think we do. It is this' representation of knowledge that
is at the crux of Al.
         In the early days it was assumed that knowledge came in sentence-like packets",
and that the best way to implant knowledge into a program was to develop a simple way
of translating facts into small passive packets of data. Then every fact would simply be a
piece of data, accessible to the programs using it. This is exemplified by chess programs,
where board Positions are coded into matrices or lists of some sort and stored efficiently
in memory where they can be retrieved and acted upon by subroutines.
         The fact that human beings store facts in a more complicated way was




Artificial Intelligence: Retrospects                                                     611

Known to psychologists for quite a while and has only recently been rediscovered by AI
workers, who are now confronting the problems of "chunked" knowledge, and the
difference between procedural and declarative types of knowledge, which is related, as
we saw in Chapter XI, to the difference between knowledge which is accessible to
introspection and knowledge which is inaccessible to introspection.
        The naive assumption that all knowledge should be coded into passive pieces of
data is actually contradicted by the most fundamental fact about computer design: that is,
how to add, subtract, multiply, and so on is not coded into pieces of data and stored in
memory; it is, in fact, represented nowhere in memory, but rather in the wiring patterns of
the hardware. A pocket calculator does not store in its memory knowledge of how to add;
that knowledge is encoded into its "guts". There is no memory location to point to if
somebody demands, "Show me where the knowledge of how to add resides in this
machine!"
        A large amount of work in Al has nevertheless gone into systems in which the
bulk of the knowledge is stored in specific places-that is, declaratively. It goes without
saying that some knowledge has to be embodied in programs; otherwise one would not
have a program at all, but merely an encyclopedia. The question is how to split up
knowledge between program and data. Not that it is always easy to distinguish between
program and data, by any means. I hope that was made clear enough in Chapter XVI. But
in the development of a system, if the programmer intuitively conceives of some
particular item as data (or as program), that may have significant repercussions on the
system's structure, because as one programs one does tend to distinguish between data-
like objects and program-like objects.
        It is important to point out that in principle, any manner of coding information
into data structures or procedures is as good as any other, in the sense that if you are not
too concerned about efficiency, what you can do in one scheme, you can do in the other.
However, reasons can be given which seem to indicate that one method is definitely
superior to the other. For instance, consider the following argument in favor of using
procedural representations only: "As soon as you try to encode features of sufficient
complexity into data, you are forced into developing what amounts to a new language, or
formalism. So in effect your data structures become program-like, with some piece of
your program serving as their interpreter; you might as well represent the same
information directly in procedural form to begin with, and obviate the extra level of
interpretation."

                DNA and Proteins Help Give Some Perspective
This argument sounds quite convincing, and yet, if interpreted a little loosely, it can be
read as an argument for the abolishment of DNA and RNA. Why encode genetic
information in DNA, when by representing it directly in proteins, you could eliminate not
just one, but two levels of interpretation? The answer is: it turns out that it is extremely
useful to have




Artificial Intelligence: Retrospects                                                    612

the same information in several different forms for different purposes. One advantage of
storing genetic information in the modular and data-like form of DNA is that two
individuals' genes can be easily recombined to form a new genotype. This would be very
difficult if the information were only in proteins. A second reason for storing information
in DNA is that it is easy to transcribe and translate it into proteins. When it is not needed,
it does not take up much room; when it is needed, it serves as a template. There is no
mechanism for copying one protein off of another; their folded tertiary structures would
make copying highly unwieldy. Complementarily, it is almost imperative to be able to get
genetic information into three-dimensional structures such as enzymes, because the
recognition and manipulation of molecules is by its nature a three-dimensional operation.
Thus the argument for purely procedural representations is seen to be quite fallacious in
the context of cells. It suggests that there are advantages to being able to switch back and
forth between procedural and declarative representations. This is probably true also in AI.
        This issue was raised by Francis Crick in a conference on communication with
extraterrestrial intelligence:

           We see on Earth that there are two molecules, one of which is good for
   replication [DNA] and one of which is good for action [proteins]. Is it possible to
   devise a system in which one molecule does both jobs, or are there perhaps strong
   arguments, from systems analysis, which might suggest (if they exist) that to divide
   the job into two gives a great advantage, This is a question to which I do not know
   the answer.14

                              Modularity of Knowledge
Another question which comes up in the representation of knowledge is modularity. How
easy is it to insert new knowledge? How easy is it to revise old knowledge? How modular
are books? It all depends. If from a tightly structured book with many cross-references a
single chapter is removed, the rest of the book may become virtually incomprehensible. It
is like trying to pull a single strand out of a spider web-you ruin the whole in doing so.
On the other hand, some books are quite modular, having independent chapters.
         Consider a straightforward theorem-generating program which uses TNT's
axioms and rules of inference. The "knowledge" of such a program has two aspects. It
resides implicitly in the axioms and rules, and explicitly in the body of theorems which
have so far been produced. Depending on which way you look at the knowledge, you will
see it either as modular or as spread all around and completely nonmodular. For instance,
suppose you had written such a program but had forgotten to include TNT's Axiom I in
the list of axioms. After the program had done many thousands of derivations, you
realized your oversight, and inserted the new axiom. The fact that you can do so in a trice
shows that the system's implicit knowledge is modular; but the new axiom's contribution
to the explicit knowledge of the system will only be reflected after a long time-after its
effects have "dif-




Artificial Intelligence: Retrospects                                                      613

fused" outwards, as the odor of perfume slowly diffuses in a room when the bottle is
broken. In that sense the new knowledge takes a long time to be incorporated.
Furthermore, if you wanted to go back and replace Axiom I by its negation, you could not
just do that by itself; you would have to delete all theorems which had involved Axiom 1
in their derivations. Clearly this system's explicit knowledge is not nearly so modular as
its implicit knowledge.
        It would be useful if we learned how to transplant knowledge modularly. Then to
teach everyone French, we would just open up their heads and operate in a fixed way on
their neural structures-then they would know how to speak French. Of course, this is only
a hilarious pipe dream.
        Another aspect of knowledge representation has to do with the way in which one
wishes to use the knowledge. Are inferences supposed to be drawn as pieces of
information arrive? Should analogies and comparisons constantly be being made between
new information and old information? In a chess program, for instance, if you want to
generate look-ahead trees, then a representation which encodes board positions with a
minimum of redundancy will be preferable to one which repeats the information in
several different ways. But if you want your program to "understand" a board position by
looking for patterns and comparing them to known patterns, then representing the same
information several times over in different forms will be more useful.

               Representing Knowledge in a Logical Formalism
There are various schools of thought concerning the best way to represent and manipulate
knowledge. One which has had great influence advocates representations using formal
notations similar to those for TNT-using propositional connectives and quantifiers. The
basic operations in such representations are, not surprisingly, formalizations of deductive
reasoning. Logical deductions can be made using rules of inference analogous to some of
those in TNT. Querying the system about some particular idea sets up a goal in the form
of a string to be derived. For example: "Is MUMON a theorem?" Then the automatic
reasoning mechanisms take over in a goal oriented way, using various methods of
problem reduction.
        For example, suppose that the proposition "All formal arithmetics are incomplete"
were known, and the program were queried, "Is Principia Mathematica incomplete?" In
scanning the list of known facts-often called the data base-the system might notice that if
it could establish that Principia Mathematica is a formal arithmetic, then it could answer
the question. Therefore the proposition "Principia Mathematica is a formal arithmetic"
would be set up as a subgoal, and then problem reduction would take over. If it could find
further things which would help in establishing (or refuting) the goal or the subgoal, it
would work on them-and so on, recursively. This process is given the name of backwards
chaining, since it begins with the goal and works its way backwards, presumably towards
things which may already be known. If one makes a graphic representation of the main
goal,




Artificial Intelligence: Retrospects                                                   614

subsidiary goals, subsubgoals, etc., a tree-like structure will arise, since the main goal
may involve several different subgoals, each of which in turn involves several
subsubgoals, etc.
        Notice that this method is not guaranteed to resolve the question, for there may be
no way of establishing within the system that Principia Mathematica is a formal
arithmetic. This does not imply, however, that either the goal or the subgoal is a false
statement-merely that they cannot be derived with the knowledge currently available to
the system. The system may print out, in such a circumstance, "I do not know" or words
to that effect. The fact that some questions are left open is of course similar to the
incompleteness from which certain well-known formal systems suffer.

                       Deductive vs. Analogical Awareness
This method affords a deductive awareness of the domain that is represented, in that
correct logical conclusions can be drawn from known facts. However, it misses
something of the human ability to spot similarities and to compare situations-it misses
what might be called analogical awareness-a crucial side of human intelligence. This is
not to say that analogical thought processes cannot be forced into such a mold, but they
do not lend themselves naturally to being captured in that kind of formalism. These days,
logic-oriented systems are not so much in vogue as other kinds, which allow complex
forms of comparisons to be carried out rather naturally.
        When you realize that knowledge representation is an altogether different ball
game than mere storage of numbers, then the idea that "a computer has the memory of an
elephant" is an easy myth to explode. What is stored in memory is not necessarily
synonymous with what a program knows; for even if a given piece of knowledge is
encoded somewhere inside a complex system, there may be no procedure, or rule, or
other type of handler of data, which can get at it-it may be inaccessible. In such a case,
you can say that the piece of knowledge has been "forgotten" because access to it has
been temporarily or permanently lost. Thus a computer program may "forget" something
on a high level which it "remembers" on a low level. This is another one of those ever-
recurring level distinctions, from which we can probably learn much about our own
selves. When a human forgets, it most likely means that a high-level pointer has been
lost-not that any information has been deleted or destroyed. This highlights the extreme
importance of keeping track of the ways in which you store incoming experiences, for
you never know in advance under what circumstances, or from what angle, you will want
to pull something out of storage.

                 From Computer Haiku to an RTN-Grammar
The complexity of the knowledge representation in human heads first hit home with me
when I was working on a program to generate English sentences "out of the blue". I had
come to this project in a rather interest-




Artificial Intelligence: Retrospects                                                   615

ing way. I had heard on the radio a few examples of so-called "Computer Haiku".
Something about them struck me deeply. There was a large element of humor and
simultaneously mystery to making a computer generate something which ordinarily
would be considered an artistic creation. I was highly amused by the humorous aspect,
and I was very motivated by the mystery-even contradiction-of programming creative
acts. So I set out to write a program even more mysteriously contradictory and humorous
than the haiku program.
        At first I was concerned with making the grammar flexible and recursive, so that
one would not have the sense that the program was merely filling in the blanks in some
template. At about that time I ran across a Scientific American article by Victor Yngve in
which he described a simple but flexible grammar which could produce a wide variety of
sentences of the type found in some children's books. I modified some of the ideas I'd
gleaned from that article and came up with a set of procedures which formed a Recursive
Transition Network grammar, as described in Chapter V. In this grammar, the selection
of words in a sentence was determined by a process which began by selecting-at random-
the overall structure of the sentence; gradually the decision-making process trickled down
through lower levels of structure until the word level and the letter level were reached. A
lot had to be done below the word level, such as inflecting verbs and making plurals of
nouns; also irregular verb and noun forms were first formed regularly, and then if they
matched entries in a table, substitutions of the proper (irregular) forms were made. As
each word reached its final form, it was printed out. The program was like the proverbial
monkey at a typewriter, but operating on several levels of linguistic structure
simultaneously-not just the letter level.
        In the early stages of developing the program, I used a totally silly vocabulary-
deliberately, since I was aiming at humor. It produced a lot of nonsense sentences, some
of which had very complicated structures, others of which were rather short. Some
excerpts are shown below:

       A male pencil who must laugh clumsily would quack. Must program not
   always crunch girl at memory? The decimal bug which spits clumsily might
   tumble. Cake who does sure take an unexpected man within relationship might
   always dump card.
   Program ought run cheerfully.
   The worthy machine ought not always paste the astronomer.
       Oh, program who ought really run off of the girl writes musician for theater.
   The businesslike relationship quacks.
   The lucky girl which can always quack will never sure quack.
   The game quacks. Professor will write pickle. A bug tumbles. Man takes the box
   who slips.

The effect is strongly surrealistic and at times a little reminiscent of




Artificial Intelligence: Retrospects                                                   616

haiku-for example. the final sample of four consecutive short sentences. At first it seemed
very funny and had a certain charm, but soon it became rather stale. After reading a few
pages of output one could sense the limits of the space in which the program was
operating; and after that, seeing random points inside that space-even though each one
was "new"-was nothing new. This is, it seems to me, a general principle: you get bored
with something not when you have exhausted its repertoire of behavior, but when you
have mapped out the limits of the space that contains its behavior. The behavior space of
a person is just about complex enough that it can continually surprise other people; but
that wasn't true of my program. I realized that my goal of producing truly humorous
output would require that far more subtlety be programmed in. But what, in this case, was
meant by "subtlety It was clear that absurd juxtapositions of words were just too
unsubtle; I needed a way to ensure that words would be used in accordance with the
realities of the world. This was where thoughts about representation of knowledge began
to enter the picture.

                                From RTN's to ATN's
The idea I adopted was to classify each word-noun, verb, preposition, etc.-in several
different "semantic dimensions". Thus, each word was a member of classes of various
sorts; then there were also superclasses-classes of classes (reminiscent of the remark by
Ulam). In principle, such aggregation could continue to any number of levels, but I
stopped at two. At any given moment, the choice of words was now semantically
restricted, because it was required that there should be agreement between the various
parts of the phrase being constructed. The idea was, for instance, that certain kinds of acts
could be performed only by animate objects; that only certain kinds of abstractions could
influence events, and so on. The decisions about what categories were reasonable, and
whether each category was better thought of as a class or a superclass, were quite
complicated. All words were branded in several different dimensions. Common
prepositions-"of", "in", etc.-had several distinct entries, corresponding to their distinct
usages. Now, the output began to be much more comprehensible-and for that reason it
was funny in a new way.

                                 A Little Turing Test
Below, I have reproduced nine selections, carefully culled from many pages of output
from later versions of my program. Along with them are three (seriously intended)
human-written sentences. Which?

   (1) Blurting may be considered as the reciprocal substitution of semiotic material
   (dubbing) for a semiotic dialogical product in a dynamic reflexion.




Artificial Intelligence: Retrospects                                                     617

   (2) Rather think of a pathway of a 'sequence' of gedankenexperiment simpletons
      where heir-lines are a prima facie case of a paradiachronic transitivity.

   (3) Think of that as a chain strength possibility of what, eventually, comes out as a
      product (epistemic conditions:,) and the product is not a Frankfurt-ish packing-
      it-all-in.

   (4) Despite the efforts, the reply, if you will, had been supported by the Orient;
      hence a fallacy will thereafter be suspended by the attitude which will be being
      held by the ambassador.

   (5) Of course, until the upheavals, the ambassador was slightly gradually
      mollycoddling the rabble.

   (6) Supposedly, refined liberty caused the attitudes insofar as peace is distilled by
      the consequences which will not eventually be caused by the command
      irrevocably insofar as peace of it is sometimes causing the intransigency
      infinitesimally surprisingly.

   (7) According to the sophists, the campaigns in the city-states, in other words, have
      been accepted by the Orient cunningly. Of course, the Orient has been separated
      by the states particularly violently.
      The Orient supports the efforts which had been supported by mankind.

   (8) Admittedly, the hierarchical origin of the fallacy, nevertheless, will be
      prophesied by the enemies of it. By the same token, the individualists will have
      testified that intransigency will not have suspended the campaigns.

   (9) Needless to say, during the upheaval which will have warranted the secrecy, the
      replies do not separate the Orient. Of course, the countries, ipso facto, are
      always probing liberty.

   (10) Although a Nobel Prize was being achieved by the humanists, yet in addition,
      it was being achieved by the serf.

   (11) An attitude will often be held by the serfs of a strife-torn nation.

   (12) Moreover, the Nobel Prizes will be achieved. By the same token, despite the
      consequence, the Nobel Prizes which will be achieved will sometimes be
      achieved by a woman.

The human-written sentences are numbers 1 to 3; they were drawn from the
contemporary journal Art-Language15 and are-as far as I can tellcompletely serious
efforts among literate and sane people to communicate something to each other. That
they appear here out of context is not too misleading, since their proper context sounds
just the same as they do.

Artificial Intelligence: Retrospects                                                   618

        My program produced the rest. Numbers 10 to 12 were chosen to show that there
were occasional bursts of total lucidity; numbers 7 to 9 are more typical of the output,
floating, in that curious and provocative netherworld between meaning and no-meaning;
and then numbers 4 to 6 pretty much transcend meaning. In a generous mood, one could
say that they stand on their own as pure "language objects", something like pieces of
abstract sculpture carved out of words instead of stone; alternatively, one could say that
they are pure pseudointellectual drivel.
        My choice of vocabulary was still aimed at producing humorous effects. The
flavor of the output is hard to characterize. Although much of it "makes sense", at least
on a single-sentence level, one definitely gets the feeling that the output is coming from a
source with no understanding of what it is saying and no reason to say it. In particular,
one senses an utter lack of visual imagery behind the words. When I saw such sentences
come pouring out of the line printer, I experienced complex emotions. I was very amused
by the silliness of the output. I was also very proud of my achievement and tried to
describe it to friends as similar to giving rules for building up meaningful stories in
Arabic out of single strokes of the pen-an exaggeration, but it pleased me to think of it
that way. And lastly I was deeply thrilled by the knowledge that this enormously
complicated machine was shunting around long trains of symbols inside it according to
rules, and that these long trains of symbols were something like thoughts in my own head
... something like them.

                            Images of What Thought Is
Of course I didn't fool myself into thinking that there was a conscious being behind those
sentences-far from it. Of all people, I was the most aware of the reasons that this program
was terribly remote from real thought. Tesler's Theorem is quite apt here: as soon as this
level of language handling ability had been mechanized, it was clear that it did not
constitute intelligence. But this strong experience left me with an image: a glimmering
sense that real thought was composed of much longer, much more complicated trains of
symbols in the brain-many trains moving simultaneously down many parallel and
crisscrossing tracks, their cars being pushed and pulled, attached and detached, switched
from track to track by a myriad neural shunting-engines ...
        It was an intangible image which I cannot convey in words, and it was only an
image. But images and intuitions and motivations lie mingled close in the mind, and my
utter fascination with this image was a constant spur to think more deeply about what
thought really could be. I have tried in other parts of this book to communicate some of
the daughter images of this original image-particularly in the Prelude, Ant Fugue.
        What stands out in my mind now, as I look back at this program from the
perspective of a dozen years, is how there is no sense of imagery behind what is being
said. The program had no idea what a serf is, what a person is, or what anything at all is.
The words were empty formal symbols, as empty




Artificial Intelligence: Retrospects                                                    619

  FIGURE 116. A meaningful story in Arabic. [From A. Khatibi and M. S~elmassi, The
           Splendour of Islamic Calligraphy (New York: Rizzoli, 1976).




Artificial Intelligence: Retrospects                                            620

 as-perhaps emptier than-the p and q of the pq-system. My program took advantage of the
 fact that when people read text, they quite naturally tend to imbue each word with its full
   flavor-as if that were necessarily attached to the group of letters which form the word.
      My program could be looked at as a formal system, whose "theorems"-the output
 sentences-had ready-made interpretations (at least to speakers of English). But unlike the
    pq-system, these "theorems" were not all true statements when interpreted that way.
                            Many were false, many were nonsense.
         In its humble way, the pq-system mirrored a tiny corner of the world. But when
my program ran, there was no mirror inside it of how the world works, except for the
small semantic constraints which it had to follow. To create such a mirror of
understanding, I would have had to wrap each concept in layers and layers of knowledge
about the world. To do this would have been another kind of effort from what I had
intended to do. Not that I didn't often think of trying to do it-but I never got around to
trying it out.

                             Higher-Level Grammars ...
In fact, I often pondered whether I could write an ATN-grammar (or some other kind of
sentence-producing program) which would only produce true sentences about the world.
Such a grammar would imbue the words with genuine meanings, in the way it happened
in the pq-system and in TNT. This idea of a language in which false statements are
ungrammatical is an old one, going back to Johann Amos Comenius, in 1633. It is very
appealing because you have a crystal ball embodied in your grammar: just write down the
statement you want to know about, and check to see if it is grammatical... . Actually,
Comenius went even further, for in his language, false statements were not only
ungrammatical-they were inexpressible!
        Carrying this thought in another direction, you might imagine a highlevel
grammar which would produce random koans. Why not? Such a grammar would be
equivalent to a formal system whose theorems are koans. And if you had such a program,
could you not arrange it to produce only genuine koans? My friend Marsha Meredith was
enthusastic about this idea of "Artificial Ism", so she tackled the project of writing a
koan-writing program. One of her early efforts produced this curious quasi-koan:

   A SMALL YOUNG MASTER WANTED A SMALL WHITE GNARLED
   BOWL. "HOW CAN WE LEARN AND UNDERSTAND WITHOUT STUDY"
   THE YOUNG MASTER ASKED A LARGE CONFUSED MASTER. THE
   CONFUSED MASTER WALKED FROM A BROWN HARD MOUNTAIN TO A
   WHITE SOFT MOUNTAIN WITH A SMALL RED STONY BOWL. THE
   CONFUSED MASTER SAW A RED SOFT HUT.THE CONFUSED MASTER
   WANTED THE HUT. "WHY DID BODHIDHARMA COME INTO CHINA?"
   THE CONFUSED MASTER




Artificial Intelligence: Retrospects                                                    621

   ASKED A LARGE ENLIGHTENED STUDENT. "THE PEACHES ARE
   LARGE", THE STUDENT ANSWERED THE CONFUSED MASTER.w "HOW
   CAN WE LEARN AND UNDERSTAND WITHOUT STUDY?" THE
   CONFUSED MASTER ASKED A LARGE OLD MASTER. THE OLD MASTER
   WALKED FROM A WHITE STONY G0025. THE OLD MASTER GOT LOST.

        Your personal decision procedure for koan genuineness probably reached a
verdict without need of the Geometric Code or the Art of Zen Strings. If the lack of
pronouns or the unsophisticated syntax didn't arouse your suspicions, that strange
"00025" towards the end must have. What is it? It is a strange fluke-a manifestation of a
bug which caused the program to print out, in place of the English word for an object, the
program's internal name for the "node" (a LISP atom, in fact) where all information
concerning that particular object was stored. So here we have a "window" onto a lower
level of the underlying Zen mind-a level that should have remained invisible.
Unfortunately, we don't have such clear windows onto the lower levels of human Zen
minds.
        The sequence of actions, though a little arbitrary, comes from a recursive LISP
procedure called "CASCADE", which creates chains of actions linked in a vaguely
causal way to each other. Although the degree of comprehension of the world possessed
by this koan generator is clearly not stupendous, work is in progress to make its output a
little more genuine seeming.

                                Grammars for Music?
Then there is music. This is a domain which you might suppose, on first thought, would
lend itself admirably to being codified in an ATN grammar, or some such program.
Whereas (to continue this naive line of thought) language relies on connections with the
outside world for meaning, music is purely formal. There is no reference to things "out
there" in the sounds of music; there is just pure syntax-note following note, chord
following chord, measure following measure, phrase following phrase...
        But wait. Something is wrong in this analysis. Why is some music so much
deeper and more beautiful than other music? It is because form, in music, is expressive-
expressive to some strange subconscious regions of our minds. The sounds of music do
not refer to serfs or city-states, but they do trigger clouds of emotion in our innermost
selves; in that sense musical meaning is dependent on intangible links from the symbols
to things in the world-those "things", in this case, being secret software structures in our
minds. No, great music will not come out of such an easy formalism as an ATN-
grammar. Pseudomusic, like pseudo-fairy tales, may well come out-and that will be a
valuable exploration for people to make-but the secrets of meaning in music lie far, far
deeper than pure syntax.
        I should clarify one point here: in principle, ATN-grammars have all the power of
any programming formalism, so if musical meaning is captur-




Artificial Intelligence: Retrospects                                                    622

able in any way at all (which I believe it is), it is capturable in an A I N - grammar. True.
But in that case, I maintain, the grammar will be defining not just musical structures, but
the entire structures of the mind of a beholder. The "grammar" will be a full grammar of
thought-not just a grammar of music.

                          Winograd's Program SHRDLU
What kind of program would it take to make human beings admit that it had some
"understanding", even if begrudgingly? What would it take before you wouldn't feel
intuitively that there is "nothing there"?
         In the years 1968-70, Terry Winograd (alias Dr. Tony Earrwig) was a doctoral
student at MIT, working on the joint problems of language and understanding. At that
time at MIT, much Al research involved the so-called blocks world-a relatively simple
domain in which problems concerning both vision and language-handling by computer
could fit easily. The blocks world consists of a table with various kinds of toy-like blocks
on it-square ones, oblong ones, triangular ones, etc., in various colors. (For a "blocks
world" of another kind, see Figure 117: the painting Mental Arithmetic by Magritte. I find
its title singularly appropriate in this context.) The vision problems in the MIT blocks
world are very tricky: how can a computer figure out, from a TV-scan of a scene with
many blocks in it, just what kinds of blocks are present, and what their relationships are?
Some blocks may be perched on top of others, some may be in front of others, there may
be shadows, and so on.

               FIGURE 117. Mental Arithmetic, by Rene Magritte (1931).




Artificial Intelligence: Retrospects                                                     623

       Winograd s work was separate from the issues of vision, however. Beginning with
the assumption that the blocks world was well represented h inside the computer's
memory, he confronted the many-faceted problem of how to get the computer to:

   (1) understand questions in English about the situation;
   (2) give answers in English to questions about the situation;
   (3) understand requests in English to manipulate the blocks;
   (4) break down each request into a sequence of operations it could do;
   (5) understand what it had done, and for what reasons;
   (6) describe its actions and their reasons, in English.

        It might seem reasonable to break up the overall program into modular
subprograms, with one module for each different part of the problem; then, after the
modules have been developed separately, to integrate them smoothly. Winograd found
that this strategy of developing independent modules posed fundamental difficulties. He
developed a radical approach, which challenged the theory that intelligence can be
compartmentalized into independent or semi-independent pieces. His program SHRDLU
named after the old code "ETAOIN SHRDLU", used by linotype operators to mark
typos in a newspaper column-did not separate the problem into clean conceptual parts.
The operations of parsing sentences, producing internal representations, reasoning about
the world represented inside itself, answering questions, and so on, were all deeply and
intricately meshed together in a procedural representation of knowledge. Some critics
have charged that his program is so tangled that it does not represent any "theory" at all
about language, nor does it contribute in any way to our insights about thought processes.
Nothing could be more wrong than such claims, in my opinion. A tour de force such as
SHRDLU may not be isomorphic to what we do-in fact, in no way should you think that
in SHRDLU, the "symbol level" has been attained-but the act of creating it and thinking
about it offers tremendous insight into the way intelligence works.

                             The Structure of SHRDLU
In fact, SHRDLU does consist of separate procedures, each of which contains some
knowledge about the world; but the procedures have such a strong interdependency that
they cannot be cleanly teased apart. The program is like a very tangled knot which resists
untangling; but the fact that you cannot untangle it does not mean that you cannot
understand it. There may be an elegant geometrical description of the entire knot even if
it is physically messy. We could go back to a metaphor from the Mu Offering, and
compare it to looking at an orchard from a "natural" angle.
        Winograd has written lucidly about SHRDLU. I quote here from his article in
Schank and Colby's book:




Artificial Intelligence: Retrospects                                                  624

   One of the basic viewpoints underlying the model is that all language use can be
   thought of as a way of activating procedures within the hearer. We can think of
   any utterance as a program-one that indirectly causes a set of operations to he
   carried out within the hearer's cognitive system. This "program writing" is indirect
   in the sense that we are dealing with an intelligent interpreter, who may take a set
   of actions which are quite different from those the speaker intended. The exact
   form is determined by his knowledge of the world, his expectations about the
   person talking to him, etc. in this program we have a simple version of this
   process of interpretation as it takes place in the robot. Each sentence interpreted by
   the robot is converted to a set of instructions in PLANNER. The program that is
   created is then executed to achieve the desired effect."


                   PLANNER Facilitates Problem Reduction
The language PLANNER, referred to here, is an At language whose principal feature is
that some of the operations necessary for problem reduction are built in-namely, the
recursive process of creating a tree of jsubgoals, subsubgoals, etc. What this means is that
such processes, instead of having to be spelled out time and time again by the
programmer, are automatically implied by so-called GOAL-statements. Someone who
reads a PLANNER program will see no explicit reference to such operations; in argon,
they are user-transparent. If one path-in the tree fails to achieve the desired goal, then the
PLANNER program will "backtrack" and try another route. "Backtracking" is the magic
word as far as PLANNER is concerned.
        Winograd's program made excellent use of these features of PLANNER-more
exactly, of MICROPLANNER, a partial implementation of the plans for PLANNER. In
the past few years, however, people with the goal of developing At have concluded that
automatic backtrack ing, as in PLANNER, has definite disadvantages, and that it will
probably not lead to their goal; therefore they have backed off from it, preferring to try
other routes to AI.
        Let us listen to further comments from Winograd on SHRDLU:

   The definition of every word is a program which is called at an appropriate point in
   the analysis, and which can do arbitrary computations involving the sentence and
   the present physical situation."

Among the examples which Winograd cites is the following:

   The different possibilities for the meaning of "the" are procedures which check
   various facts about the context, then prescribe actions such as "Look for a unique
   object in the data base which fits this description", or "Assert that the object being
   described is unique as far as the speaker is concerned." The program incorporates a
   variety of heuristics for deciding what part of the context is relevant. 18

It is amazing how deep this problem with the word "the" is. It is probably safe to say that
writing a program which can fully handle the top five words

Artificial Intelligence: Retrospects                                                      625

of English-"the", "of", "and", "a", and "to"-would be equivalent to solving the entire
problem of AI, and hence tantamount to knowing what intelligence and consciousness
are. A small digression: the five most common nouns in English are-according to the
Word Frequency Book compiled by John B. Carroll et al-"time", "people", "way",
"water", and "words" (in that order). The amazing thing about this is that most people
have no idea that we think in such abstract terms. Ask your friends, and 10 to 1 they'll
guess such words as "man", "house", "car", "dog", and "money". And while we're on the
subject of frequencies-the top twelve letters in English, in order, according to
Mergenthaler, are: "ETAOIN SHRDLU".
         One amusing feature of SHRDLU which runs totally against the stereotype of
computers as "number crunchers" is this fact, pointed out by Winograd: "Our system does
not accept numbers in numeric form, and has only been taught to count to ten."19 With
all its mathematical underpinning, SHRDLU is a mathematical ignoramus! Just like Aunt
Hillary, SHRDLU doesn't know anything about the lower levels which make it up. Its
knowledge is largely procedural (see particularly the remark by "Dr, Tony Earrwig" in
section 11 of the previous Dialogue).
         It is interesting to contrast the procedural embedding of knowledge in SHRDLU
with the knowledge in my sentence-generation program. All of the syntactical knowledge
in my program was procedurally embedded in Augmented Transition Networks, written
in the language Algol; but the semantic knowledge-the information about semantic class
membership-was static: it was contained in a short list of numbers after each word. There
were a few words, such as the auxiliary verbs "to be", "to have", and others, which were
represented totally in procedures in Algol, but they were the exceptions. By contrast, in
SHRDLU, all words were represented as programs. Here is a case which demonstrates
that, despite the theoretical equivalence of data and programs, in practice the choice of
one over the other has major consequences.

                                Syntax and Semantics
And now, a few more words from Winograd:

   Our program does not operate by first parsing a sentence, then doing semantic
   analysis, and finally by using deduction to produce a response. These three
   activities go on concurrently throughout the understanding of a sentence. As soon
   as a piece of syntactic structure begins to take shape, a semantic program
   is called to see whether it might make sense, and the resultant answer can direct the
   parsing. In deciding whether it makes sense, the semantic routine may call
   deductive processes and ask questions about the real world. As an
   example, in sentence 34 of the Dialogue ("Put the blue pyramid on the block in the
   box"), the parser first comes up with "the blue pyramid on the block" as a candidate
   for a noun group. At this point, semantic analysis is done, and since "the" is
   definite, a check is made in the data base for the object being referred to. When no
   such object is found, the parsing is redirected to find the noun group "the blue
   pyramid". It will then go on to find "on the block




Artificial Intelligence: Retrospects                                                   626

   in the box" as a single phrase indicating a location .-.. Thus there is a continuing
   interplay between the different sorts of analysis, with the results of one affecting
   the others .20

        It is extremely interesting that in natural language, syntax and semantics are so
deeply intertwined. Last Chapter, in discussing the elusive concept of "form", we had
broken the notion into two categories: syntactic form, which is detectable by a
predictably terminating decision procedure, and semantic form, which is not. But here,
Winograd is telling us that-at least when the usual senses of "syntax" and "semantics" are
taken-they merge right into each other, in natural language. The external form of a
sentence-that is, its composition in terms of elementary signs-does not divide up so neatly
into syntactic and semantic aspects. This significant point for linguistics.
        Here are some final comments on SHRDLU by Winograd.

Let us look at what the system would do with a simple description like "a red cube which
supports a pyramid". The description will use concepts like BLOCK, RED, PYRAMID,
and EQUIDIMENSIONAL-all parts of the sys

FIGURE 118. Procedural representation of "a red cube which supports a pyramid."
Adapted from Roger Schank and Kenneth Colby, Computer Models of Thought and
Language (San Francisco: W. H. Freeman, 1973), p. 172.




Artificial Intelligence: Retrospects                                                   627

   tem's underlying categorization of the world. The result can be represented in a
   flow chart like that in Figure 118. Note that this is a program for finding an object
   fitting the description. It would\_ then be incorporated into a command for doing
   something with the object, a question asking something about it, or, if it appeared in
   a statement, it would become part of the program which was generated to represent
   the meaning for later use. Note that this bit of program could also be used as a test
   to see whether an object fit the description, if the first FIND instruction were told
   in advance to look only at that particular object.
            At first glance, it seems that there is too much structure in this program, as
   we don't like to think of the meaning of a simple phrase as explicitly containing
   loops, conditional tests, and other programming details. The solution is to provide
   an internal language that contains the appropriate looping and checking as its
   primitives, and in which the representation of the process is as simple as the
   description. The program described in Figure 11S would be written in PLANNER
   looking something like what is below:
                       (GOAL (IS ?X 1 BLOCK))
                       (GOAL (COLOR-OF ?X1 RED))
                       (GOAL (EQUIDIMENSIONAL ?X1))
                       (GOAL (IS ?X2 PYRAMID))
                       (GOAL (SUPPORT ?X1 ?X2))
   The loops of the flowchart are implicit in PLANNER'S backtrack control structure.
   The description is evaluated by proceeding down the list until some goal fails, at
   which time the system backs up automatically to the last point where a decision
   was made, trying a different possibility. A decision can be made whenever a new
   object name or VARIABLE (indicated by the prefix
   ") such as "?X 1" or "?X2" appears. The variables are used by the pattern matcher.
   If' they have already been assigned to a particular item, it checks to see whether the
   GOAL is true for that item. If not, it checks for all possible items which satisfy the
   GOAL, by choosing one, and then taking successive ones whenever backtracking
   occurs to that point. Thus, even the distinction between testing and choosing is
   implicit.21
One significant strategy decision in devising this program was to not translate all the way
from English into LISP, but only partway-into PLANNER. Thus (since the PLANNER
interpreter is itself written in LISP), a new intermediate level-PLANNER-was inserted
between the top-level language (English) and the bottom-level language (machine
language). Once a PLANNER program had been made from an English sentence
fragment, then it could be sent off to the PLANNER interpreter, and the higher levels of
SHRDLU would be freed up, to work on new tasks.
        This kind of decision constantly crops up: How many levels should a system
have? How much and what kind of "intelligence" should be placed on which level? These
are some of the hardest problems facing AI today. Since we know so little about natural
intelligence, it is hard for us to figure out which level of an artificially intelligent system
should carry out what part of a task.



Artificial Intelligence: Retrospects                                                       628

      This gives you a glimpse behind the scenes of the Dialogue preceding this
Chapter. Next Chapter, we shall meet new and speculative ideas for AI.




Artificial Intelligence: Retrospects                                       629

